\documentclass[11pt,a4paper]{article}

\newcommand{\CourseCode}{CSAI2017P: Applied Machine Learning}
\newcommand{\AssignmentName}{Mid-Sem Question Bank -- Answers}

% Shared settings for Applied Machine Learning lab assignments and marking schemes.
% Layered on the house notes preamble; keep free of \documentclass.
%
% Callers must define \CourseCode, \AssignmentName before \begin{document}.

% Depth-agnostic locate of the house notes preamble: briefs compile from
% public/lab-assignments/LAxx/latex/ and marking schemes from
% private/solutions/lab-assignments/LAxx/latex/, which sit at different depths.
\IfFileExists{../../../../public/_shared/notes-common.tex}{%
  % Shared notes settings for Applied Machine Learning lectures (UPES).
% Keep this file free of \documentclass to allow reuse via \input.

\usepackage[utf8]{inputenc}
\usepackage[T1]{fontenc}
\usepackage{lmodern}          % scalable T1 fonts; without this pdflatex falls back
\input{glyphtounicode}        % to bitmapped EC Type 3 and the PDFs are neither
\pdfgentounicode=1            % searchable nor screen-reader friendly
\usepackage[a4paper,margin=1in]{geometry}
\usepackage{hyperref}
\usepackage{xurl}
\usepackage{booktabs}
\usepackage{amsmath}
\usepackage{amssymb}
\usepackage{listings}
\usepackage{xcolor}
\usepackage{enumitem}
\usepackage{parskip}

% Code listings (Python-oriented for this subject).
\lstset{
  language=Python,
  basicstyle=\ttfamily\small,
  keywordstyle=\color{blue},
  commentstyle=\color{gray},
  stringstyle=\color{teal},
  showstringspaces=false,
  columns=fullflexible,
  frame=single,
  framerule=0.2pt,
  breaklines=true
}

\setlist[itemize]{topsep=0.3em,itemsep=0.2em}
\setlist[enumerate]{topsep=0.3em,itemsep=0.2em}

% Simple per-section source line.
\newcommand{\notesources}[1]{\par\noindent{\small\color{gray}\textit{Sources:} \sloppy #1\par}}

% Unit-wise source strings for this subject (used by slides + notes).
% Path is relative to the lecture `latex/` folder (the compilation CWD).
\IfFileExists{../../../_shared/aml-sources.tex}{%
  % Source strings for Applied Machine Learning (CSAI2017P) — used by slides + notes.
% Prescribed texts and reference books are taken from the course syllabus.
% Support texts are the author-hosted open-access editions:
%   statlearning.com            (ISL with Python)
%   probml.github.io/pml-book   (Murphy, Probabilistic Machine Learning)
%   d2l.ai                      (Dive into Deep Learning)
%   mml-book.github.io          (Mathematics for Machine Learning)

\newcommand{\AMLRepoURL}{https://github.com/tali7c/Applied-Machine-Learning}

% --- prescribed -------------------------------------------------------------
\newcommand{\AMLPrescribed}{%
M\"uller and Guido, \textit{Introduction to Machine Learning with Python}, O'Reilly, 2016.;%
Bishop, \textit{Pattern Recognition and Machine Learning}, Springer, 2006.;%
Murphy, \textit{Machine Learning: A Probabilistic Perspective}, MIT Press, 2012.;%
Han, Kamber and Pei, \textit{Data Mining: Concepts and Techniques} (3e), Morgan Kaufmann, 2012.%
}

% --- unit-wise ---------------------------------------------------------------
\newcommand{\AMLUnitISources}{%
M\"uller and Guido, \textit{Introduction to Machine Learning with Python}, Ch. 1.;%
Bishop, \textit{Pattern Recognition and Machine Learning}, Ch. 1.;%
Murphy, \textit{Probabilistic Machine Learning: An Introduction}, MIT Press, 2022, Ch. 1.;%
Mitchell, \textit{Machine Learning}, McGraw Hill, 1997, Ch. 1.;%
Scikit-learn documentation, accessed 2026-08-04.%
}

\newcommand{\AMLUnitIISources}{%
Bishop, \textit{Pattern Recognition and Machine Learning}, \S1.5, \S4.3, \S7.1.;%
Zhang, Lipton, Li and Smola, \textit{Dive into Deep Learning}, \S3.4.;%
Shalev-Shwartz and Ben-David, \textit{Understanding Machine Learning}, Ch. 15.;%
Murphy, \textit{Probabilistic Machine Learning: An Introduction}, Ch. 5.%
}

\newcommand{\AMLUnitIIISources}{%
Zhang, Lipton, Li and Smola, \textit{Dive into Deep Learning}, Ch. 12 (Optimization Algorithms).;%
Deisenroth, Faisal and Ong, \textit{Mathematics for Machine Learning}, Ch. 7.;%
Kingma and Ba, ``Adam: A Method for Stochastic Optimization,'' ICLR 2015.;%
Ruder, ``An overview of gradient descent optimization algorithms,'' arXiv:1609.04747, 2016.%
}

\newcommand{\AMLUnitIVSources}{%
Han, Kamber and Pei, \textit{Data Mining: Concepts and Techniques} (3e), Ch. 3.;%
M\"uller and Guido, \textit{Introduction to Machine Learning with Python}, Ch. 3--4.;%
James, Witten, Hastie, Tibshirani and Taylor, \textit{An Introduction to Statistical Learning with Python}, Ch. 5--6, 12.;%
Scikit-learn documentation (preprocessing, model selection), accessed 2026-08-04.%
}

\newcommand{\AMLUnitVSources}{%
James et al., \textit{An Introduction to Statistical Learning with Python}, Ch. 3--6.;%
Bishop, \textit{Pattern Recognition and Machine Learning}, Ch. 3.;%
Deisenroth, Faisal and Ong, \textit{Mathematics for Machine Learning}, Ch. 9.;%
M\"uller and Guido, \textit{Introduction to Machine Learning with Python}, \S2.3.%
}

\newcommand{\AMLUnitVISources}{%
Han, Kamber and Pei, \textit{Data Mining: Concepts and Techniques} (3e), Ch. 8--9.;%
James et al., \textit{An Introduction to Statistical Learning with Python}, Ch. 8--9.;%
Bishop, \textit{Pattern Recognition and Machine Learning}, Ch. 4--5, 7, 14.;%
Zhang et al., \textit{Dive into Deep Learning}, Ch. 5.%
}

\newcommand{\AMLUnitVIISources}{%
Han, Kamber and Pei, \textit{Data Mining: Concepts and Techniques} (3e), Ch. 10.;%
James et al., \textit{An Introduction to Statistical Learning with Python}, Ch. 12.;%
M\"uller and Guido, \textit{Introduction to Machine Learning with Python}, \S3.5.;%
Ester, Kriegel, Sander and Xu, ``A density-based algorithm for discovering clusters (DBSCAN),'' KDD 1996.%
}

\newcommand{\AMLLabSources}{%
M\"uller and Guido, \textit{Introduction to Machine Learning with Python}, companion notebooks.;%
McKinney, \textit{Python for Data Analysis} (3e), O'Reilly, 2022.;%
Scikit-learn user guide, accessed 2026-08-04.;%
Matplotlib and Seaborn documentation, accessed 2026-08-04.%
}
%
}{%
  % Faculty copies live under `private/` so their compile CWD is different.
  \IfFileExists{../../../../public/_shared/aml-sources.tex}{%
    % Source strings for Applied Machine Learning (CSAI2017P) — used by slides + notes.
% Prescribed texts and reference books are taken from the course syllabus.
% Support texts are the author-hosted open-access editions:
%   statlearning.com            (ISL with Python)
%   probml.github.io/pml-book   (Murphy, Probabilistic Machine Learning)
%   d2l.ai                      (Dive into Deep Learning)
%   mml-book.github.io          (Mathematics for Machine Learning)

\newcommand{\AMLRepoURL}{https://github.com/tali7c/Applied-Machine-Learning}

% --- prescribed -------------------------------------------------------------
\newcommand{\AMLPrescribed}{%
M\"uller and Guido, \textit{Introduction to Machine Learning with Python}, O'Reilly, 2016.;%
Bishop, \textit{Pattern Recognition and Machine Learning}, Springer, 2006.;%
Murphy, \textit{Machine Learning: A Probabilistic Perspective}, MIT Press, 2012.;%
Han, Kamber and Pei, \textit{Data Mining: Concepts and Techniques} (3e), Morgan Kaufmann, 2012.%
}

% --- unit-wise ---------------------------------------------------------------
\newcommand{\AMLUnitISources}{%
M\"uller and Guido, \textit{Introduction to Machine Learning with Python}, Ch. 1.;%
Bishop, \textit{Pattern Recognition and Machine Learning}, Ch. 1.;%
Murphy, \textit{Probabilistic Machine Learning: An Introduction}, MIT Press, 2022, Ch. 1.;%
Mitchell, \textit{Machine Learning}, McGraw Hill, 1997, Ch. 1.;%
Scikit-learn documentation, accessed 2026-08-04.%
}

\newcommand{\AMLUnitIISources}{%
Bishop, \textit{Pattern Recognition and Machine Learning}, \S1.5, \S4.3, \S7.1.;%
Zhang, Lipton, Li and Smola, \textit{Dive into Deep Learning}, \S3.4.;%
Shalev-Shwartz and Ben-David, \textit{Understanding Machine Learning}, Ch. 15.;%
Murphy, \textit{Probabilistic Machine Learning: An Introduction}, Ch. 5.%
}

\newcommand{\AMLUnitIIISources}{%
Zhang, Lipton, Li and Smola, \textit{Dive into Deep Learning}, Ch. 12 (Optimization Algorithms).;%
Deisenroth, Faisal and Ong, \textit{Mathematics for Machine Learning}, Ch. 7.;%
Kingma and Ba, ``Adam: A Method for Stochastic Optimization,'' ICLR 2015.;%
Ruder, ``An overview of gradient descent optimization algorithms,'' arXiv:1609.04747, 2016.%
}

\newcommand{\AMLUnitIVSources}{%
Han, Kamber and Pei, \textit{Data Mining: Concepts and Techniques} (3e), Ch. 3.;%
M\"uller and Guido, \textit{Introduction to Machine Learning with Python}, Ch. 3--4.;%
James, Witten, Hastie, Tibshirani and Taylor, \textit{An Introduction to Statistical Learning with Python}, Ch. 5--6, 12.;%
Scikit-learn documentation (preprocessing, model selection), accessed 2026-08-04.%
}

\newcommand{\AMLUnitVSources}{%
James et al., \textit{An Introduction to Statistical Learning with Python}, Ch. 3--6.;%
Bishop, \textit{Pattern Recognition and Machine Learning}, Ch. 3.;%
Deisenroth, Faisal and Ong, \textit{Mathematics for Machine Learning}, Ch. 9.;%
M\"uller and Guido, \textit{Introduction to Machine Learning with Python}, \S2.3.%
}

\newcommand{\AMLUnitVISources}{%
Han, Kamber and Pei, \textit{Data Mining: Concepts and Techniques} (3e), Ch. 8--9.;%
James et al., \textit{An Introduction to Statistical Learning with Python}, Ch. 8--9.;%
Bishop, \textit{Pattern Recognition and Machine Learning}, Ch. 4--5, 7, 14.;%
Zhang et al., \textit{Dive into Deep Learning}, Ch. 5.%
}

\newcommand{\AMLUnitVIISources}{%
Han, Kamber and Pei, \textit{Data Mining: Concepts and Techniques} (3e), Ch. 10.;%
James et al., \textit{An Introduction to Statistical Learning with Python}, Ch. 12.;%
M\"uller and Guido, \textit{Introduction to Machine Learning with Python}, \S3.5.;%
Ester, Kriegel, Sander and Xu, ``A density-based algorithm for discovering clusters (DBSCAN),'' KDD 1996.%
}

\newcommand{\AMLLabSources}{%
M\"uller and Guido, \textit{Introduction to Machine Learning with Python}, companion notebooks.;%
McKinney, \textit{Python for Data Analysis} (3e), O'Reilly, 2022.;%
Scikit-learn user guide, accessed 2026-08-04.;%
Matplotlib and Seaborn documentation, accessed 2026-08-04.%
}
%
  }{%
    % Question papers live at `private/<family>/<instrument>/`.
    \IfFileExists{../../../public/_shared/aml-sources.tex}{%
      % Source strings for Applied Machine Learning (CSAI2017P) — used by slides + notes.
% Prescribed texts and reference books are taken from the course syllabus.
% Support texts are the author-hosted open-access editions:
%   statlearning.com            (ISL with Python)
%   probml.github.io/pml-book   (Murphy, Probabilistic Machine Learning)
%   d2l.ai                      (Dive into Deep Learning)
%   mml-book.github.io          (Mathematics for Machine Learning)

\newcommand{\AMLRepoURL}{https://github.com/tali7c/Applied-Machine-Learning}

% --- prescribed -------------------------------------------------------------
\newcommand{\AMLPrescribed}{%
M\"uller and Guido, \textit{Introduction to Machine Learning with Python}, O'Reilly, 2016.;%
Bishop, \textit{Pattern Recognition and Machine Learning}, Springer, 2006.;%
Murphy, \textit{Machine Learning: A Probabilistic Perspective}, MIT Press, 2012.;%
Han, Kamber and Pei, \textit{Data Mining: Concepts and Techniques} (3e), Morgan Kaufmann, 2012.%
}

% --- unit-wise ---------------------------------------------------------------
\newcommand{\AMLUnitISources}{%
M\"uller and Guido, \textit{Introduction to Machine Learning with Python}, Ch. 1.;%
Bishop, \textit{Pattern Recognition and Machine Learning}, Ch. 1.;%
Murphy, \textit{Probabilistic Machine Learning: An Introduction}, MIT Press, 2022, Ch. 1.;%
Mitchell, \textit{Machine Learning}, McGraw Hill, 1997, Ch. 1.;%
Scikit-learn documentation, accessed 2026-08-04.%
}

\newcommand{\AMLUnitIISources}{%
Bishop, \textit{Pattern Recognition and Machine Learning}, \S1.5, \S4.3, \S7.1.;%
Zhang, Lipton, Li and Smola, \textit{Dive into Deep Learning}, \S3.4.;%
Shalev-Shwartz and Ben-David, \textit{Understanding Machine Learning}, Ch. 15.;%
Murphy, \textit{Probabilistic Machine Learning: An Introduction}, Ch. 5.%
}

\newcommand{\AMLUnitIIISources}{%
Zhang, Lipton, Li and Smola, \textit{Dive into Deep Learning}, Ch. 12 (Optimization Algorithms).;%
Deisenroth, Faisal and Ong, \textit{Mathematics for Machine Learning}, Ch. 7.;%
Kingma and Ba, ``Adam: A Method for Stochastic Optimization,'' ICLR 2015.;%
Ruder, ``An overview of gradient descent optimization algorithms,'' arXiv:1609.04747, 2016.%
}

\newcommand{\AMLUnitIVSources}{%
Han, Kamber and Pei, \textit{Data Mining: Concepts and Techniques} (3e), Ch. 3.;%
M\"uller and Guido, \textit{Introduction to Machine Learning with Python}, Ch. 3--4.;%
James, Witten, Hastie, Tibshirani and Taylor, \textit{An Introduction to Statistical Learning with Python}, Ch. 5--6, 12.;%
Scikit-learn documentation (preprocessing, model selection), accessed 2026-08-04.%
}

\newcommand{\AMLUnitVSources}{%
James et al., \textit{An Introduction to Statistical Learning with Python}, Ch. 3--6.;%
Bishop, \textit{Pattern Recognition and Machine Learning}, Ch. 3.;%
Deisenroth, Faisal and Ong, \textit{Mathematics for Machine Learning}, Ch. 9.;%
M\"uller and Guido, \textit{Introduction to Machine Learning with Python}, \S2.3.%
}

\newcommand{\AMLUnitVISources}{%
Han, Kamber and Pei, \textit{Data Mining: Concepts and Techniques} (3e), Ch. 8--9.;%
James et al., \textit{An Introduction to Statistical Learning with Python}, Ch. 8--9.;%
Bishop, \textit{Pattern Recognition and Machine Learning}, Ch. 4--5, 7, 14.;%
Zhang et al., \textit{Dive into Deep Learning}, Ch. 5.%
}

\newcommand{\AMLUnitVIISources}{%
Han, Kamber and Pei, \textit{Data Mining: Concepts and Techniques} (3e), Ch. 10.;%
James et al., \textit{An Introduction to Statistical Learning with Python}, Ch. 12.;%
M\"uller and Guido, \textit{Introduction to Machine Learning with Python}, \S3.5.;%
Ester, Kriegel, Sander and Xu, ``A density-based algorithm for discovering clusters (DBSCAN),'' KDD 1996.%
}

\newcommand{\AMLLabSources}{%
M\"uller and Guido, \textit{Introduction to Machine Learning with Python}, companion notebooks.;%
McKinney, \textit{Python for Data Analysis} (3e), O'Reilly, 2022.;%
Scikit-learn user guide, accessed 2026-08-04.;%
Matplotlib and Seaborn documentation, accessed 2026-08-04.%
}
%
    }{%
      % Marking schemes live at `private/solutions/<family>/<id>/latex/`.
      \IfFileExists{../../../../../public/_shared/aml-sources.tex}{%
        % Source strings for Applied Machine Learning (CSAI2017P) — used by slides + notes.
% Prescribed texts and reference books are taken from the course syllabus.
% Support texts are the author-hosted open-access editions:
%   statlearning.com            (ISL with Python)
%   probml.github.io/pml-book   (Murphy, Probabilistic Machine Learning)
%   d2l.ai                      (Dive into Deep Learning)
%   mml-book.github.io          (Mathematics for Machine Learning)

\newcommand{\AMLRepoURL}{https://github.com/tali7c/Applied-Machine-Learning}

% --- prescribed -------------------------------------------------------------
\newcommand{\AMLPrescribed}{%
M\"uller and Guido, \textit{Introduction to Machine Learning with Python}, O'Reilly, 2016.;%
Bishop, \textit{Pattern Recognition and Machine Learning}, Springer, 2006.;%
Murphy, \textit{Machine Learning: A Probabilistic Perspective}, MIT Press, 2012.;%
Han, Kamber and Pei, \textit{Data Mining: Concepts and Techniques} (3e), Morgan Kaufmann, 2012.%
}

% --- unit-wise ---------------------------------------------------------------
\newcommand{\AMLUnitISources}{%
M\"uller and Guido, \textit{Introduction to Machine Learning with Python}, Ch. 1.;%
Bishop, \textit{Pattern Recognition and Machine Learning}, Ch. 1.;%
Murphy, \textit{Probabilistic Machine Learning: An Introduction}, MIT Press, 2022, Ch. 1.;%
Mitchell, \textit{Machine Learning}, McGraw Hill, 1997, Ch. 1.;%
Scikit-learn documentation, accessed 2026-08-04.%
}

\newcommand{\AMLUnitIISources}{%
Bishop, \textit{Pattern Recognition and Machine Learning}, \S1.5, \S4.3, \S7.1.;%
Zhang, Lipton, Li and Smola, \textit{Dive into Deep Learning}, \S3.4.;%
Shalev-Shwartz and Ben-David, \textit{Understanding Machine Learning}, Ch. 15.;%
Murphy, \textit{Probabilistic Machine Learning: An Introduction}, Ch. 5.%
}

\newcommand{\AMLUnitIIISources}{%
Zhang, Lipton, Li and Smola, \textit{Dive into Deep Learning}, Ch. 12 (Optimization Algorithms).;%
Deisenroth, Faisal and Ong, \textit{Mathematics for Machine Learning}, Ch. 7.;%
Kingma and Ba, ``Adam: A Method for Stochastic Optimization,'' ICLR 2015.;%
Ruder, ``An overview of gradient descent optimization algorithms,'' arXiv:1609.04747, 2016.%
}

\newcommand{\AMLUnitIVSources}{%
Han, Kamber and Pei, \textit{Data Mining: Concepts and Techniques} (3e), Ch. 3.;%
M\"uller and Guido, \textit{Introduction to Machine Learning with Python}, Ch. 3--4.;%
James, Witten, Hastie, Tibshirani and Taylor, \textit{An Introduction to Statistical Learning with Python}, Ch. 5--6, 12.;%
Scikit-learn documentation (preprocessing, model selection), accessed 2026-08-04.%
}

\newcommand{\AMLUnitVSources}{%
James et al., \textit{An Introduction to Statistical Learning with Python}, Ch. 3--6.;%
Bishop, \textit{Pattern Recognition and Machine Learning}, Ch. 3.;%
Deisenroth, Faisal and Ong, \textit{Mathematics for Machine Learning}, Ch. 9.;%
M\"uller and Guido, \textit{Introduction to Machine Learning with Python}, \S2.3.%
}

\newcommand{\AMLUnitVISources}{%
Han, Kamber and Pei, \textit{Data Mining: Concepts and Techniques} (3e), Ch. 8--9.;%
James et al., \textit{An Introduction to Statistical Learning with Python}, Ch. 8--9.;%
Bishop, \textit{Pattern Recognition and Machine Learning}, Ch. 4--5, 7, 14.;%
Zhang et al., \textit{Dive into Deep Learning}, Ch. 5.%
}

\newcommand{\AMLUnitVIISources}{%
Han, Kamber and Pei, \textit{Data Mining: Concepts and Techniques} (3e), Ch. 10.;%
James et al., \textit{An Introduction to Statistical Learning with Python}, Ch. 12.;%
M\"uller and Guido, \textit{Introduction to Machine Learning with Python}, \S3.5.;%
Ester, Kriegel, Sander and Xu, ``A density-based algorithm for discovering clusters (DBSCAN),'' KDD 1996.%
}

\newcommand{\AMLLabSources}{%
M\"uller and Guido, \textit{Introduction to Machine Learning with Python}, companion notebooks.;%
McKinney, \textit{Python for Data Analysis} (3e), O'Reilly, 2022.;%
Scikit-learn user guide, accessed 2026-08-04.;%
Matplotlib and Seaborn documentation, accessed 2026-08-04.%
}
%
      }{%
        % Source strings for Applied Machine Learning (CSAI2017P) — used by slides + notes.
% Prescribed texts and reference books are taken from the course syllabus.
% Support texts are the author-hosted open-access editions:
%   statlearning.com            (ISL with Python)
%   probml.github.io/pml-book   (Murphy, Probabilistic Machine Learning)
%   d2l.ai                      (Dive into Deep Learning)
%   mml-book.github.io          (Mathematics for Machine Learning)

\newcommand{\AMLRepoURL}{https://github.com/tali7c/Applied-Machine-Learning}

% --- prescribed -------------------------------------------------------------
\newcommand{\AMLPrescribed}{%
M\"uller and Guido, \textit{Introduction to Machine Learning with Python}, O'Reilly, 2016.;%
Bishop, \textit{Pattern Recognition and Machine Learning}, Springer, 2006.;%
Murphy, \textit{Machine Learning: A Probabilistic Perspective}, MIT Press, 2012.;%
Han, Kamber and Pei, \textit{Data Mining: Concepts and Techniques} (3e), Morgan Kaufmann, 2012.%
}

% --- unit-wise ---------------------------------------------------------------
\newcommand{\AMLUnitISources}{%
M\"uller and Guido, \textit{Introduction to Machine Learning with Python}, Ch. 1.;%
Bishop, \textit{Pattern Recognition and Machine Learning}, Ch. 1.;%
Murphy, \textit{Probabilistic Machine Learning: An Introduction}, MIT Press, 2022, Ch. 1.;%
Mitchell, \textit{Machine Learning}, McGraw Hill, 1997, Ch. 1.;%
Scikit-learn documentation, accessed 2026-08-04.%
}

\newcommand{\AMLUnitIISources}{%
Bishop, \textit{Pattern Recognition and Machine Learning}, \S1.5, \S4.3, \S7.1.;%
Zhang, Lipton, Li and Smola, \textit{Dive into Deep Learning}, \S3.4.;%
Shalev-Shwartz and Ben-David, \textit{Understanding Machine Learning}, Ch. 15.;%
Murphy, \textit{Probabilistic Machine Learning: An Introduction}, Ch. 5.%
}

\newcommand{\AMLUnitIIISources}{%
Zhang, Lipton, Li and Smola, \textit{Dive into Deep Learning}, Ch. 12 (Optimization Algorithms).;%
Deisenroth, Faisal and Ong, \textit{Mathematics for Machine Learning}, Ch. 7.;%
Kingma and Ba, ``Adam: A Method for Stochastic Optimization,'' ICLR 2015.;%
Ruder, ``An overview of gradient descent optimization algorithms,'' arXiv:1609.04747, 2016.%
}

\newcommand{\AMLUnitIVSources}{%
Han, Kamber and Pei, \textit{Data Mining: Concepts and Techniques} (3e), Ch. 3.;%
M\"uller and Guido, \textit{Introduction to Machine Learning with Python}, Ch. 3--4.;%
James, Witten, Hastie, Tibshirani and Taylor, \textit{An Introduction to Statistical Learning with Python}, Ch. 5--6, 12.;%
Scikit-learn documentation (preprocessing, model selection), accessed 2026-08-04.%
}

\newcommand{\AMLUnitVSources}{%
James et al., \textit{An Introduction to Statistical Learning with Python}, Ch. 3--6.;%
Bishop, \textit{Pattern Recognition and Machine Learning}, Ch. 3.;%
Deisenroth, Faisal and Ong, \textit{Mathematics for Machine Learning}, Ch. 9.;%
M\"uller and Guido, \textit{Introduction to Machine Learning with Python}, \S2.3.%
}

\newcommand{\AMLUnitVISources}{%
Han, Kamber and Pei, \textit{Data Mining: Concepts and Techniques} (3e), Ch. 8--9.;%
James et al., \textit{An Introduction to Statistical Learning with Python}, Ch. 8--9.;%
Bishop, \textit{Pattern Recognition and Machine Learning}, Ch. 4--5, 7, 14.;%
Zhang et al., \textit{Dive into Deep Learning}, Ch. 5.%
}

\newcommand{\AMLUnitVIISources}{%
Han, Kamber and Pei, \textit{Data Mining: Concepts and Techniques} (3e), Ch. 10.;%
James et al., \textit{An Introduction to Statistical Learning with Python}, Ch. 12.;%
M\"uller and Guido, \textit{Introduction to Machine Learning with Python}, \S3.5.;%
Ester, Kriegel, Sander and Xu, ``A density-based algorithm for discovering clusters (DBSCAN),'' KDD 1996.%
}

\newcommand{\AMLLabSources}{%
M\"uller and Guido, \textit{Introduction to Machine Learning with Python}, companion notebooks.;%
McKinney, \textit{Python for Data Analysis} (3e), O'Reilly, 2022.;%
Scikit-learn user guide, accessed 2026-08-04.;%
Matplotlib and Seaborn documentation, accessed 2026-08-04.%
}
%
      }%
    }%
  }%
}
%
}{%
  \IfFileExists{../../../../../public/_shared/notes-common.tex}{%
    % Shared notes settings for Applied Machine Learning lectures (UPES).
% Keep this file free of \documentclass to allow reuse via \input.

\usepackage[utf8]{inputenc}
\usepackage[T1]{fontenc}
\usepackage{lmodern}          % scalable T1 fonts; without this pdflatex falls back
\input{glyphtounicode}        % to bitmapped EC Type 3 and the PDFs are neither
\pdfgentounicode=1            % searchable nor screen-reader friendly
\usepackage[a4paper,margin=1in]{geometry}
\usepackage{hyperref}
\usepackage{xurl}
\usepackage{booktabs}
\usepackage{amsmath}
\usepackage{amssymb}
\usepackage{listings}
\usepackage{xcolor}
\usepackage{enumitem}
\usepackage{parskip}

% Code listings (Python-oriented for this subject).
\lstset{
  language=Python,
  basicstyle=\ttfamily\small,
  keywordstyle=\color{blue},
  commentstyle=\color{gray},
  stringstyle=\color{teal},
  showstringspaces=false,
  columns=fullflexible,
  frame=single,
  framerule=0.2pt,
  breaklines=true
}

\setlist[itemize]{topsep=0.3em,itemsep=0.2em}
\setlist[enumerate]{topsep=0.3em,itemsep=0.2em}

% Simple per-section source line.
\newcommand{\notesources}[1]{\par\noindent{\small\color{gray}\textit{Sources:} \sloppy #1\par}}

% Unit-wise source strings for this subject (used by slides + notes).
% Path is relative to the lecture `latex/` folder (the compilation CWD).
\IfFileExists{../../../_shared/aml-sources.tex}{%
  % Source strings for Applied Machine Learning (CSAI2017P) — used by slides + notes.
% Prescribed texts and reference books are taken from the course syllabus.
% Support texts are the author-hosted open-access editions:
%   statlearning.com            (ISL with Python)
%   probml.github.io/pml-book   (Murphy, Probabilistic Machine Learning)
%   d2l.ai                      (Dive into Deep Learning)
%   mml-book.github.io          (Mathematics for Machine Learning)

\newcommand{\AMLRepoURL}{https://github.com/tali7c/Applied-Machine-Learning}

% --- prescribed -------------------------------------------------------------
\newcommand{\AMLPrescribed}{%
M\"uller and Guido, \textit{Introduction to Machine Learning with Python}, O'Reilly, 2016.;%
Bishop, \textit{Pattern Recognition and Machine Learning}, Springer, 2006.;%
Murphy, \textit{Machine Learning: A Probabilistic Perspective}, MIT Press, 2012.;%
Han, Kamber and Pei, \textit{Data Mining: Concepts and Techniques} (3e), Morgan Kaufmann, 2012.%
}

% --- unit-wise ---------------------------------------------------------------
\newcommand{\AMLUnitISources}{%
M\"uller and Guido, \textit{Introduction to Machine Learning with Python}, Ch. 1.;%
Bishop, \textit{Pattern Recognition and Machine Learning}, Ch. 1.;%
Murphy, \textit{Probabilistic Machine Learning: An Introduction}, MIT Press, 2022, Ch. 1.;%
Mitchell, \textit{Machine Learning}, McGraw Hill, 1997, Ch. 1.;%
Scikit-learn documentation, accessed 2026-08-04.%
}

\newcommand{\AMLUnitIISources}{%
Bishop, \textit{Pattern Recognition and Machine Learning}, \S1.5, \S4.3, \S7.1.;%
Zhang, Lipton, Li and Smola, \textit{Dive into Deep Learning}, \S3.4.;%
Shalev-Shwartz and Ben-David, \textit{Understanding Machine Learning}, Ch. 15.;%
Murphy, \textit{Probabilistic Machine Learning: An Introduction}, Ch. 5.%
}

\newcommand{\AMLUnitIIISources}{%
Zhang, Lipton, Li and Smola, \textit{Dive into Deep Learning}, Ch. 12 (Optimization Algorithms).;%
Deisenroth, Faisal and Ong, \textit{Mathematics for Machine Learning}, Ch. 7.;%
Kingma and Ba, ``Adam: A Method for Stochastic Optimization,'' ICLR 2015.;%
Ruder, ``An overview of gradient descent optimization algorithms,'' arXiv:1609.04747, 2016.%
}

\newcommand{\AMLUnitIVSources}{%
Han, Kamber and Pei, \textit{Data Mining: Concepts and Techniques} (3e), Ch. 3.;%
M\"uller and Guido, \textit{Introduction to Machine Learning with Python}, Ch. 3--4.;%
James, Witten, Hastie, Tibshirani and Taylor, \textit{An Introduction to Statistical Learning with Python}, Ch. 5--6, 12.;%
Scikit-learn documentation (preprocessing, model selection), accessed 2026-08-04.%
}

\newcommand{\AMLUnitVSources}{%
James et al., \textit{An Introduction to Statistical Learning with Python}, Ch. 3--6.;%
Bishop, \textit{Pattern Recognition and Machine Learning}, Ch. 3.;%
Deisenroth, Faisal and Ong, \textit{Mathematics for Machine Learning}, Ch. 9.;%
M\"uller and Guido, \textit{Introduction to Machine Learning with Python}, \S2.3.%
}

\newcommand{\AMLUnitVISources}{%
Han, Kamber and Pei, \textit{Data Mining: Concepts and Techniques} (3e), Ch. 8--9.;%
James et al., \textit{An Introduction to Statistical Learning with Python}, Ch. 8--9.;%
Bishop, \textit{Pattern Recognition and Machine Learning}, Ch. 4--5, 7, 14.;%
Zhang et al., \textit{Dive into Deep Learning}, Ch. 5.%
}

\newcommand{\AMLUnitVIISources}{%
Han, Kamber and Pei, \textit{Data Mining: Concepts and Techniques} (3e), Ch. 10.;%
James et al., \textit{An Introduction to Statistical Learning with Python}, Ch. 12.;%
M\"uller and Guido, \textit{Introduction to Machine Learning with Python}, \S3.5.;%
Ester, Kriegel, Sander and Xu, ``A density-based algorithm for discovering clusters (DBSCAN),'' KDD 1996.%
}

\newcommand{\AMLLabSources}{%
M\"uller and Guido, \textit{Introduction to Machine Learning with Python}, companion notebooks.;%
McKinney, \textit{Python for Data Analysis} (3e), O'Reilly, 2022.;%
Scikit-learn user guide, accessed 2026-08-04.;%
Matplotlib and Seaborn documentation, accessed 2026-08-04.%
}
%
}{%
  % Faculty copies live under `private/` so their compile CWD is different.
  \IfFileExists{../../../../public/_shared/aml-sources.tex}{%
    % Source strings for Applied Machine Learning (CSAI2017P) — used by slides + notes.
% Prescribed texts and reference books are taken from the course syllabus.
% Support texts are the author-hosted open-access editions:
%   statlearning.com            (ISL with Python)
%   probml.github.io/pml-book   (Murphy, Probabilistic Machine Learning)
%   d2l.ai                      (Dive into Deep Learning)
%   mml-book.github.io          (Mathematics for Machine Learning)

\newcommand{\AMLRepoURL}{https://github.com/tali7c/Applied-Machine-Learning}

% --- prescribed -------------------------------------------------------------
\newcommand{\AMLPrescribed}{%
M\"uller and Guido, \textit{Introduction to Machine Learning with Python}, O'Reilly, 2016.;%
Bishop, \textit{Pattern Recognition and Machine Learning}, Springer, 2006.;%
Murphy, \textit{Machine Learning: A Probabilistic Perspective}, MIT Press, 2012.;%
Han, Kamber and Pei, \textit{Data Mining: Concepts and Techniques} (3e), Morgan Kaufmann, 2012.%
}

% --- unit-wise ---------------------------------------------------------------
\newcommand{\AMLUnitISources}{%
M\"uller and Guido, \textit{Introduction to Machine Learning with Python}, Ch. 1.;%
Bishop, \textit{Pattern Recognition and Machine Learning}, Ch. 1.;%
Murphy, \textit{Probabilistic Machine Learning: An Introduction}, MIT Press, 2022, Ch. 1.;%
Mitchell, \textit{Machine Learning}, McGraw Hill, 1997, Ch. 1.;%
Scikit-learn documentation, accessed 2026-08-04.%
}

\newcommand{\AMLUnitIISources}{%
Bishop, \textit{Pattern Recognition and Machine Learning}, \S1.5, \S4.3, \S7.1.;%
Zhang, Lipton, Li and Smola, \textit{Dive into Deep Learning}, \S3.4.;%
Shalev-Shwartz and Ben-David, \textit{Understanding Machine Learning}, Ch. 15.;%
Murphy, \textit{Probabilistic Machine Learning: An Introduction}, Ch. 5.%
}

\newcommand{\AMLUnitIIISources}{%
Zhang, Lipton, Li and Smola, \textit{Dive into Deep Learning}, Ch. 12 (Optimization Algorithms).;%
Deisenroth, Faisal and Ong, \textit{Mathematics for Machine Learning}, Ch. 7.;%
Kingma and Ba, ``Adam: A Method for Stochastic Optimization,'' ICLR 2015.;%
Ruder, ``An overview of gradient descent optimization algorithms,'' arXiv:1609.04747, 2016.%
}

\newcommand{\AMLUnitIVSources}{%
Han, Kamber and Pei, \textit{Data Mining: Concepts and Techniques} (3e), Ch. 3.;%
M\"uller and Guido, \textit{Introduction to Machine Learning with Python}, Ch. 3--4.;%
James, Witten, Hastie, Tibshirani and Taylor, \textit{An Introduction to Statistical Learning with Python}, Ch. 5--6, 12.;%
Scikit-learn documentation (preprocessing, model selection), accessed 2026-08-04.%
}

\newcommand{\AMLUnitVSources}{%
James et al., \textit{An Introduction to Statistical Learning with Python}, Ch. 3--6.;%
Bishop, \textit{Pattern Recognition and Machine Learning}, Ch. 3.;%
Deisenroth, Faisal and Ong, \textit{Mathematics for Machine Learning}, Ch. 9.;%
M\"uller and Guido, \textit{Introduction to Machine Learning with Python}, \S2.3.%
}

\newcommand{\AMLUnitVISources}{%
Han, Kamber and Pei, \textit{Data Mining: Concepts and Techniques} (3e), Ch. 8--9.;%
James et al., \textit{An Introduction to Statistical Learning with Python}, Ch. 8--9.;%
Bishop, \textit{Pattern Recognition and Machine Learning}, Ch. 4--5, 7, 14.;%
Zhang et al., \textit{Dive into Deep Learning}, Ch. 5.%
}

\newcommand{\AMLUnitVIISources}{%
Han, Kamber and Pei, \textit{Data Mining: Concepts and Techniques} (3e), Ch. 10.;%
James et al., \textit{An Introduction to Statistical Learning with Python}, Ch. 12.;%
M\"uller and Guido, \textit{Introduction to Machine Learning with Python}, \S3.5.;%
Ester, Kriegel, Sander and Xu, ``A density-based algorithm for discovering clusters (DBSCAN),'' KDD 1996.%
}

\newcommand{\AMLLabSources}{%
M\"uller and Guido, \textit{Introduction to Machine Learning with Python}, companion notebooks.;%
McKinney, \textit{Python for Data Analysis} (3e), O'Reilly, 2022.;%
Scikit-learn user guide, accessed 2026-08-04.;%
Matplotlib and Seaborn documentation, accessed 2026-08-04.%
}
%
  }{%
    % Question papers live at `private/<family>/<instrument>/`.
    \IfFileExists{../../../public/_shared/aml-sources.tex}{%
      % Source strings for Applied Machine Learning (CSAI2017P) — used by slides + notes.
% Prescribed texts and reference books are taken from the course syllabus.
% Support texts are the author-hosted open-access editions:
%   statlearning.com            (ISL with Python)
%   probml.github.io/pml-book   (Murphy, Probabilistic Machine Learning)
%   d2l.ai                      (Dive into Deep Learning)
%   mml-book.github.io          (Mathematics for Machine Learning)

\newcommand{\AMLRepoURL}{https://github.com/tali7c/Applied-Machine-Learning}

% --- prescribed -------------------------------------------------------------
\newcommand{\AMLPrescribed}{%
M\"uller and Guido, \textit{Introduction to Machine Learning with Python}, O'Reilly, 2016.;%
Bishop, \textit{Pattern Recognition and Machine Learning}, Springer, 2006.;%
Murphy, \textit{Machine Learning: A Probabilistic Perspective}, MIT Press, 2012.;%
Han, Kamber and Pei, \textit{Data Mining: Concepts and Techniques} (3e), Morgan Kaufmann, 2012.%
}

% --- unit-wise ---------------------------------------------------------------
\newcommand{\AMLUnitISources}{%
M\"uller and Guido, \textit{Introduction to Machine Learning with Python}, Ch. 1.;%
Bishop, \textit{Pattern Recognition and Machine Learning}, Ch. 1.;%
Murphy, \textit{Probabilistic Machine Learning: An Introduction}, MIT Press, 2022, Ch. 1.;%
Mitchell, \textit{Machine Learning}, McGraw Hill, 1997, Ch. 1.;%
Scikit-learn documentation, accessed 2026-08-04.%
}

\newcommand{\AMLUnitIISources}{%
Bishop, \textit{Pattern Recognition and Machine Learning}, \S1.5, \S4.3, \S7.1.;%
Zhang, Lipton, Li and Smola, \textit{Dive into Deep Learning}, \S3.4.;%
Shalev-Shwartz and Ben-David, \textit{Understanding Machine Learning}, Ch. 15.;%
Murphy, \textit{Probabilistic Machine Learning: An Introduction}, Ch. 5.%
}

\newcommand{\AMLUnitIIISources}{%
Zhang, Lipton, Li and Smola, \textit{Dive into Deep Learning}, Ch. 12 (Optimization Algorithms).;%
Deisenroth, Faisal and Ong, \textit{Mathematics for Machine Learning}, Ch. 7.;%
Kingma and Ba, ``Adam: A Method for Stochastic Optimization,'' ICLR 2015.;%
Ruder, ``An overview of gradient descent optimization algorithms,'' arXiv:1609.04747, 2016.%
}

\newcommand{\AMLUnitIVSources}{%
Han, Kamber and Pei, \textit{Data Mining: Concepts and Techniques} (3e), Ch. 3.;%
M\"uller and Guido, \textit{Introduction to Machine Learning with Python}, Ch. 3--4.;%
James, Witten, Hastie, Tibshirani and Taylor, \textit{An Introduction to Statistical Learning with Python}, Ch. 5--6, 12.;%
Scikit-learn documentation (preprocessing, model selection), accessed 2026-08-04.%
}

\newcommand{\AMLUnitVSources}{%
James et al., \textit{An Introduction to Statistical Learning with Python}, Ch. 3--6.;%
Bishop, \textit{Pattern Recognition and Machine Learning}, Ch. 3.;%
Deisenroth, Faisal and Ong, \textit{Mathematics for Machine Learning}, Ch. 9.;%
M\"uller and Guido, \textit{Introduction to Machine Learning with Python}, \S2.3.%
}

\newcommand{\AMLUnitVISources}{%
Han, Kamber and Pei, \textit{Data Mining: Concepts and Techniques} (3e), Ch. 8--9.;%
James et al., \textit{An Introduction to Statistical Learning with Python}, Ch. 8--9.;%
Bishop, \textit{Pattern Recognition and Machine Learning}, Ch. 4--5, 7, 14.;%
Zhang et al., \textit{Dive into Deep Learning}, Ch. 5.%
}

\newcommand{\AMLUnitVIISources}{%
Han, Kamber and Pei, \textit{Data Mining: Concepts and Techniques} (3e), Ch. 10.;%
James et al., \textit{An Introduction to Statistical Learning with Python}, Ch. 12.;%
M\"uller and Guido, \textit{Introduction to Machine Learning with Python}, \S3.5.;%
Ester, Kriegel, Sander and Xu, ``A density-based algorithm for discovering clusters (DBSCAN),'' KDD 1996.%
}

\newcommand{\AMLLabSources}{%
M\"uller and Guido, \textit{Introduction to Machine Learning with Python}, companion notebooks.;%
McKinney, \textit{Python for Data Analysis} (3e), O'Reilly, 2022.;%
Scikit-learn user guide, accessed 2026-08-04.;%
Matplotlib and Seaborn documentation, accessed 2026-08-04.%
}
%
    }{%
      % Marking schemes live at `private/solutions/<family>/<id>/latex/`.
      \IfFileExists{../../../../../public/_shared/aml-sources.tex}{%
        % Source strings for Applied Machine Learning (CSAI2017P) — used by slides + notes.
% Prescribed texts and reference books are taken from the course syllabus.
% Support texts are the author-hosted open-access editions:
%   statlearning.com            (ISL with Python)
%   probml.github.io/pml-book   (Murphy, Probabilistic Machine Learning)
%   d2l.ai                      (Dive into Deep Learning)
%   mml-book.github.io          (Mathematics for Machine Learning)

\newcommand{\AMLRepoURL}{https://github.com/tali7c/Applied-Machine-Learning}

% --- prescribed -------------------------------------------------------------
\newcommand{\AMLPrescribed}{%
M\"uller and Guido, \textit{Introduction to Machine Learning with Python}, O'Reilly, 2016.;%
Bishop, \textit{Pattern Recognition and Machine Learning}, Springer, 2006.;%
Murphy, \textit{Machine Learning: A Probabilistic Perspective}, MIT Press, 2012.;%
Han, Kamber and Pei, \textit{Data Mining: Concepts and Techniques} (3e), Morgan Kaufmann, 2012.%
}

% --- unit-wise ---------------------------------------------------------------
\newcommand{\AMLUnitISources}{%
M\"uller and Guido, \textit{Introduction to Machine Learning with Python}, Ch. 1.;%
Bishop, \textit{Pattern Recognition and Machine Learning}, Ch. 1.;%
Murphy, \textit{Probabilistic Machine Learning: An Introduction}, MIT Press, 2022, Ch. 1.;%
Mitchell, \textit{Machine Learning}, McGraw Hill, 1997, Ch. 1.;%
Scikit-learn documentation, accessed 2026-08-04.%
}

\newcommand{\AMLUnitIISources}{%
Bishop, \textit{Pattern Recognition and Machine Learning}, \S1.5, \S4.3, \S7.1.;%
Zhang, Lipton, Li and Smola, \textit{Dive into Deep Learning}, \S3.4.;%
Shalev-Shwartz and Ben-David, \textit{Understanding Machine Learning}, Ch. 15.;%
Murphy, \textit{Probabilistic Machine Learning: An Introduction}, Ch. 5.%
}

\newcommand{\AMLUnitIIISources}{%
Zhang, Lipton, Li and Smola, \textit{Dive into Deep Learning}, Ch. 12 (Optimization Algorithms).;%
Deisenroth, Faisal and Ong, \textit{Mathematics for Machine Learning}, Ch. 7.;%
Kingma and Ba, ``Adam: A Method for Stochastic Optimization,'' ICLR 2015.;%
Ruder, ``An overview of gradient descent optimization algorithms,'' arXiv:1609.04747, 2016.%
}

\newcommand{\AMLUnitIVSources}{%
Han, Kamber and Pei, \textit{Data Mining: Concepts and Techniques} (3e), Ch. 3.;%
M\"uller and Guido, \textit{Introduction to Machine Learning with Python}, Ch. 3--4.;%
James, Witten, Hastie, Tibshirani and Taylor, \textit{An Introduction to Statistical Learning with Python}, Ch. 5--6, 12.;%
Scikit-learn documentation (preprocessing, model selection), accessed 2026-08-04.%
}

\newcommand{\AMLUnitVSources}{%
James et al., \textit{An Introduction to Statistical Learning with Python}, Ch. 3--6.;%
Bishop, \textit{Pattern Recognition and Machine Learning}, Ch. 3.;%
Deisenroth, Faisal and Ong, \textit{Mathematics for Machine Learning}, Ch. 9.;%
M\"uller and Guido, \textit{Introduction to Machine Learning with Python}, \S2.3.%
}

\newcommand{\AMLUnitVISources}{%
Han, Kamber and Pei, \textit{Data Mining: Concepts and Techniques} (3e), Ch. 8--9.;%
James et al., \textit{An Introduction to Statistical Learning with Python}, Ch. 8--9.;%
Bishop, \textit{Pattern Recognition and Machine Learning}, Ch. 4--5, 7, 14.;%
Zhang et al., \textit{Dive into Deep Learning}, Ch. 5.%
}

\newcommand{\AMLUnitVIISources}{%
Han, Kamber and Pei, \textit{Data Mining: Concepts and Techniques} (3e), Ch. 10.;%
James et al., \textit{An Introduction to Statistical Learning with Python}, Ch. 12.;%
M\"uller and Guido, \textit{Introduction to Machine Learning with Python}, \S3.5.;%
Ester, Kriegel, Sander and Xu, ``A density-based algorithm for discovering clusters (DBSCAN),'' KDD 1996.%
}

\newcommand{\AMLLabSources}{%
M\"uller and Guido, \textit{Introduction to Machine Learning with Python}, companion notebooks.;%
McKinney, \textit{Python for Data Analysis} (3e), O'Reilly, 2022.;%
Scikit-learn user guide, accessed 2026-08-04.;%
Matplotlib and Seaborn documentation, accessed 2026-08-04.%
}
%
      }{%
        % Source strings for Applied Machine Learning (CSAI2017P) — used by slides + notes.
% Prescribed texts and reference books are taken from the course syllabus.
% Support texts are the author-hosted open-access editions:
%   statlearning.com            (ISL with Python)
%   probml.github.io/pml-book   (Murphy, Probabilistic Machine Learning)
%   d2l.ai                      (Dive into Deep Learning)
%   mml-book.github.io          (Mathematics for Machine Learning)

\newcommand{\AMLRepoURL}{https://github.com/tali7c/Applied-Machine-Learning}

% --- prescribed -------------------------------------------------------------
\newcommand{\AMLPrescribed}{%
M\"uller and Guido, \textit{Introduction to Machine Learning with Python}, O'Reilly, 2016.;%
Bishop, \textit{Pattern Recognition and Machine Learning}, Springer, 2006.;%
Murphy, \textit{Machine Learning: A Probabilistic Perspective}, MIT Press, 2012.;%
Han, Kamber and Pei, \textit{Data Mining: Concepts and Techniques} (3e), Morgan Kaufmann, 2012.%
}

% --- unit-wise ---------------------------------------------------------------
\newcommand{\AMLUnitISources}{%
M\"uller and Guido, \textit{Introduction to Machine Learning with Python}, Ch. 1.;%
Bishop, \textit{Pattern Recognition and Machine Learning}, Ch. 1.;%
Murphy, \textit{Probabilistic Machine Learning: An Introduction}, MIT Press, 2022, Ch. 1.;%
Mitchell, \textit{Machine Learning}, McGraw Hill, 1997, Ch. 1.;%
Scikit-learn documentation, accessed 2026-08-04.%
}

\newcommand{\AMLUnitIISources}{%
Bishop, \textit{Pattern Recognition and Machine Learning}, \S1.5, \S4.3, \S7.1.;%
Zhang, Lipton, Li and Smola, \textit{Dive into Deep Learning}, \S3.4.;%
Shalev-Shwartz and Ben-David, \textit{Understanding Machine Learning}, Ch. 15.;%
Murphy, \textit{Probabilistic Machine Learning: An Introduction}, Ch. 5.%
}

\newcommand{\AMLUnitIIISources}{%
Zhang, Lipton, Li and Smola, \textit{Dive into Deep Learning}, Ch. 12 (Optimization Algorithms).;%
Deisenroth, Faisal and Ong, \textit{Mathematics for Machine Learning}, Ch. 7.;%
Kingma and Ba, ``Adam: A Method for Stochastic Optimization,'' ICLR 2015.;%
Ruder, ``An overview of gradient descent optimization algorithms,'' arXiv:1609.04747, 2016.%
}

\newcommand{\AMLUnitIVSources}{%
Han, Kamber and Pei, \textit{Data Mining: Concepts and Techniques} (3e), Ch. 3.;%
M\"uller and Guido, \textit{Introduction to Machine Learning with Python}, Ch. 3--4.;%
James, Witten, Hastie, Tibshirani and Taylor, \textit{An Introduction to Statistical Learning with Python}, Ch. 5--6, 12.;%
Scikit-learn documentation (preprocessing, model selection), accessed 2026-08-04.%
}

\newcommand{\AMLUnitVSources}{%
James et al., \textit{An Introduction to Statistical Learning with Python}, Ch. 3--6.;%
Bishop, \textit{Pattern Recognition and Machine Learning}, Ch. 3.;%
Deisenroth, Faisal and Ong, \textit{Mathematics for Machine Learning}, Ch. 9.;%
M\"uller and Guido, \textit{Introduction to Machine Learning with Python}, \S2.3.%
}

\newcommand{\AMLUnitVISources}{%
Han, Kamber and Pei, \textit{Data Mining: Concepts and Techniques} (3e), Ch. 8--9.;%
James et al., \textit{An Introduction to Statistical Learning with Python}, Ch. 8--9.;%
Bishop, \textit{Pattern Recognition and Machine Learning}, Ch. 4--5, 7, 14.;%
Zhang et al., \textit{Dive into Deep Learning}, Ch. 5.%
}

\newcommand{\AMLUnitVIISources}{%
Han, Kamber and Pei, \textit{Data Mining: Concepts and Techniques} (3e), Ch. 10.;%
James et al., \textit{An Introduction to Statistical Learning with Python}, Ch. 12.;%
M\"uller and Guido, \textit{Introduction to Machine Learning with Python}, \S3.5.;%
Ester, Kriegel, Sander and Xu, ``A density-based algorithm for discovering clusters (DBSCAN),'' KDD 1996.%
}

\newcommand{\AMLLabSources}{%
M\"uller and Guido, \textit{Introduction to Machine Learning with Python}, companion notebooks.;%
McKinney, \textit{Python for Data Analysis} (3e), O'Reilly, 2022.;%
Scikit-learn user guide, accessed 2026-08-04.;%
Matplotlib and Seaborn documentation, accessed 2026-08-04.%
}
%
      }%
    }%
  }%
}
%
  }{%
    % Shared notes settings for Applied Machine Learning lectures (UPES).
% Keep this file free of \documentclass to allow reuse via \input.

\usepackage[utf8]{inputenc}
\usepackage[T1]{fontenc}
\usepackage{lmodern}          % scalable T1 fonts; without this pdflatex falls back
\input{glyphtounicode}        % to bitmapped EC Type 3 and the PDFs are neither
\pdfgentounicode=1            % searchable nor screen-reader friendly
\usepackage[a4paper,margin=1in]{geometry}
\usepackage{hyperref}
\usepackage{xurl}
\usepackage{booktabs}
\usepackage{amsmath}
\usepackage{amssymb}
\usepackage{listings}
\usepackage{xcolor}
\usepackage{enumitem}
\usepackage{parskip}

% Code listings (Python-oriented for this subject).
\lstset{
  language=Python,
  basicstyle=\ttfamily\small,
  keywordstyle=\color{blue},
  commentstyle=\color{gray},
  stringstyle=\color{teal},
  showstringspaces=false,
  columns=fullflexible,
  frame=single,
  framerule=0.2pt,
  breaklines=true
}

\setlist[itemize]{topsep=0.3em,itemsep=0.2em}
\setlist[enumerate]{topsep=0.3em,itemsep=0.2em}

% Simple per-section source line.
\newcommand{\notesources}[1]{\par\noindent{\small\color{gray}\textit{Sources:} \sloppy #1\par}}

% Unit-wise source strings for this subject (used by slides + notes).
% Path is relative to the lecture `latex/` folder (the compilation CWD).
\IfFileExists{../../../_shared/aml-sources.tex}{%
  % Source strings for Applied Machine Learning (CSAI2017P) — used by slides + notes.
% Prescribed texts and reference books are taken from the course syllabus.
% Support texts are the author-hosted open-access editions:
%   statlearning.com            (ISL with Python)
%   probml.github.io/pml-book   (Murphy, Probabilistic Machine Learning)
%   d2l.ai                      (Dive into Deep Learning)
%   mml-book.github.io          (Mathematics for Machine Learning)

\newcommand{\AMLRepoURL}{https://github.com/tali7c/Applied-Machine-Learning}

% --- prescribed -------------------------------------------------------------
\newcommand{\AMLPrescribed}{%
M\"uller and Guido, \textit{Introduction to Machine Learning with Python}, O'Reilly, 2016.;%
Bishop, \textit{Pattern Recognition and Machine Learning}, Springer, 2006.;%
Murphy, \textit{Machine Learning: A Probabilistic Perspective}, MIT Press, 2012.;%
Han, Kamber and Pei, \textit{Data Mining: Concepts and Techniques} (3e), Morgan Kaufmann, 2012.%
}

% --- unit-wise ---------------------------------------------------------------
\newcommand{\AMLUnitISources}{%
M\"uller and Guido, \textit{Introduction to Machine Learning with Python}, Ch. 1.;%
Bishop, \textit{Pattern Recognition and Machine Learning}, Ch. 1.;%
Murphy, \textit{Probabilistic Machine Learning: An Introduction}, MIT Press, 2022, Ch. 1.;%
Mitchell, \textit{Machine Learning}, McGraw Hill, 1997, Ch. 1.;%
Scikit-learn documentation, accessed 2026-08-04.%
}

\newcommand{\AMLUnitIISources}{%
Bishop, \textit{Pattern Recognition and Machine Learning}, \S1.5, \S4.3, \S7.1.;%
Zhang, Lipton, Li and Smola, \textit{Dive into Deep Learning}, \S3.4.;%
Shalev-Shwartz and Ben-David, \textit{Understanding Machine Learning}, Ch. 15.;%
Murphy, \textit{Probabilistic Machine Learning: An Introduction}, Ch. 5.%
}

\newcommand{\AMLUnitIIISources}{%
Zhang, Lipton, Li and Smola, \textit{Dive into Deep Learning}, Ch. 12 (Optimization Algorithms).;%
Deisenroth, Faisal and Ong, \textit{Mathematics for Machine Learning}, Ch. 7.;%
Kingma and Ba, ``Adam: A Method for Stochastic Optimization,'' ICLR 2015.;%
Ruder, ``An overview of gradient descent optimization algorithms,'' arXiv:1609.04747, 2016.%
}

\newcommand{\AMLUnitIVSources}{%
Han, Kamber and Pei, \textit{Data Mining: Concepts and Techniques} (3e), Ch. 3.;%
M\"uller and Guido, \textit{Introduction to Machine Learning with Python}, Ch. 3--4.;%
James, Witten, Hastie, Tibshirani and Taylor, \textit{An Introduction to Statistical Learning with Python}, Ch. 5--6, 12.;%
Scikit-learn documentation (preprocessing, model selection), accessed 2026-08-04.%
}

\newcommand{\AMLUnitVSources}{%
James et al., \textit{An Introduction to Statistical Learning with Python}, Ch. 3--6.;%
Bishop, \textit{Pattern Recognition and Machine Learning}, Ch. 3.;%
Deisenroth, Faisal and Ong, \textit{Mathematics for Machine Learning}, Ch. 9.;%
M\"uller and Guido, \textit{Introduction to Machine Learning with Python}, \S2.3.%
}

\newcommand{\AMLUnitVISources}{%
Han, Kamber and Pei, \textit{Data Mining: Concepts and Techniques} (3e), Ch. 8--9.;%
James et al., \textit{An Introduction to Statistical Learning with Python}, Ch. 8--9.;%
Bishop, \textit{Pattern Recognition and Machine Learning}, Ch. 4--5, 7, 14.;%
Zhang et al., \textit{Dive into Deep Learning}, Ch. 5.%
}

\newcommand{\AMLUnitVIISources}{%
Han, Kamber and Pei, \textit{Data Mining: Concepts and Techniques} (3e), Ch. 10.;%
James et al., \textit{An Introduction to Statistical Learning with Python}, Ch. 12.;%
M\"uller and Guido, \textit{Introduction to Machine Learning with Python}, \S3.5.;%
Ester, Kriegel, Sander and Xu, ``A density-based algorithm for discovering clusters (DBSCAN),'' KDD 1996.%
}

\newcommand{\AMLLabSources}{%
M\"uller and Guido, \textit{Introduction to Machine Learning with Python}, companion notebooks.;%
McKinney, \textit{Python for Data Analysis} (3e), O'Reilly, 2022.;%
Scikit-learn user guide, accessed 2026-08-04.;%
Matplotlib and Seaborn documentation, accessed 2026-08-04.%
}
%
}{%
  % Faculty copies live under `private/` so their compile CWD is different.
  \IfFileExists{../../../../public/_shared/aml-sources.tex}{%
    % Source strings for Applied Machine Learning (CSAI2017P) — used by slides + notes.
% Prescribed texts and reference books are taken from the course syllabus.
% Support texts are the author-hosted open-access editions:
%   statlearning.com            (ISL with Python)
%   probml.github.io/pml-book   (Murphy, Probabilistic Machine Learning)
%   d2l.ai                      (Dive into Deep Learning)
%   mml-book.github.io          (Mathematics for Machine Learning)

\newcommand{\AMLRepoURL}{https://github.com/tali7c/Applied-Machine-Learning}

% --- prescribed -------------------------------------------------------------
\newcommand{\AMLPrescribed}{%
M\"uller and Guido, \textit{Introduction to Machine Learning with Python}, O'Reilly, 2016.;%
Bishop, \textit{Pattern Recognition and Machine Learning}, Springer, 2006.;%
Murphy, \textit{Machine Learning: A Probabilistic Perspective}, MIT Press, 2012.;%
Han, Kamber and Pei, \textit{Data Mining: Concepts and Techniques} (3e), Morgan Kaufmann, 2012.%
}

% --- unit-wise ---------------------------------------------------------------
\newcommand{\AMLUnitISources}{%
M\"uller and Guido, \textit{Introduction to Machine Learning with Python}, Ch. 1.;%
Bishop, \textit{Pattern Recognition and Machine Learning}, Ch. 1.;%
Murphy, \textit{Probabilistic Machine Learning: An Introduction}, MIT Press, 2022, Ch. 1.;%
Mitchell, \textit{Machine Learning}, McGraw Hill, 1997, Ch. 1.;%
Scikit-learn documentation, accessed 2026-08-04.%
}

\newcommand{\AMLUnitIISources}{%
Bishop, \textit{Pattern Recognition and Machine Learning}, \S1.5, \S4.3, \S7.1.;%
Zhang, Lipton, Li and Smola, \textit{Dive into Deep Learning}, \S3.4.;%
Shalev-Shwartz and Ben-David, \textit{Understanding Machine Learning}, Ch. 15.;%
Murphy, \textit{Probabilistic Machine Learning: An Introduction}, Ch. 5.%
}

\newcommand{\AMLUnitIIISources}{%
Zhang, Lipton, Li and Smola, \textit{Dive into Deep Learning}, Ch. 12 (Optimization Algorithms).;%
Deisenroth, Faisal and Ong, \textit{Mathematics for Machine Learning}, Ch. 7.;%
Kingma and Ba, ``Adam: A Method for Stochastic Optimization,'' ICLR 2015.;%
Ruder, ``An overview of gradient descent optimization algorithms,'' arXiv:1609.04747, 2016.%
}

\newcommand{\AMLUnitIVSources}{%
Han, Kamber and Pei, \textit{Data Mining: Concepts and Techniques} (3e), Ch. 3.;%
M\"uller and Guido, \textit{Introduction to Machine Learning with Python}, Ch. 3--4.;%
James, Witten, Hastie, Tibshirani and Taylor, \textit{An Introduction to Statistical Learning with Python}, Ch. 5--6, 12.;%
Scikit-learn documentation (preprocessing, model selection), accessed 2026-08-04.%
}

\newcommand{\AMLUnitVSources}{%
James et al., \textit{An Introduction to Statistical Learning with Python}, Ch. 3--6.;%
Bishop, \textit{Pattern Recognition and Machine Learning}, Ch. 3.;%
Deisenroth, Faisal and Ong, \textit{Mathematics for Machine Learning}, Ch. 9.;%
M\"uller and Guido, \textit{Introduction to Machine Learning with Python}, \S2.3.%
}

\newcommand{\AMLUnitVISources}{%
Han, Kamber and Pei, \textit{Data Mining: Concepts and Techniques} (3e), Ch. 8--9.;%
James et al., \textit{An Introduction to Statistical Learning with Python}, Ch. 8--9.;%
Bishop, \textit{Pattern Recognition and Machine Learning}, Ch. 4--5, 7, 14.;%
Zhang et al., \textit{Dive into Deep Learning}, Ch. 5.%
}

\newcommand{\AMLUnitVIISources}{%
Han, Kamber and Pei, \textit{Data Mining: Concepts and Techniques} (3e), Ch. 10.;%
James et al., \textit{An Introduction to Statistical Learning with Python}, Ch. 12.;%
M\"uller and Guido, \textit{Introduction to Machine Learning with Python}, \S3.5.;%
Ester, Kriegel, Sander and Xu, ``A density-based algorithm for discovering clusters (DBSCAN),'' KDD 1996.%
}

\newcommand{\AMLLabSources}{%
M\"uller and Guido, \textit{Introduction to Machine Learning with Python}, companion notebooks.;%
McKinney, \textit{Python for Data Analysis} (3e), O'Reilly, 2022.;%
Scikit-learn user guide, accessed 2026-08-04.;%
Matplotlib and Seaborn documentation, accessed 2026-08-04.%
}
%
  }{%
    % Question papers live at `private/<family>/<instrument>/`.
    \IfFileExists{../../../public/_shared/aml-sources.tex}{%
      % Source strings for Applied Machine Learning (CSAI2017P) — used by slides + notes.
% Prescribed texts and reference books are taken from the course syllabus.
% Support texts are the author-hosted open-access editions:
%   statlearning.com            (ISL with Python)
%   probml.github.io/pml-book   (Murphy, Probabilistic Machine Learning)
%   d2l.ai                      (Dive into Deep Learning)
%   mml-book.github.io          (Mathematics for Machine Learning)

\newcommand{\AMLRepoURL}{https://github.com/tali7c/Applied-Machine-Learning}

% --- prescribed -------------------------------------------------------------
\newcommand{\AMLPrescribed}{%
M\"uller and Guido, \textit{Introduction to Machine Learning with Python}, O'Reilly, 2016.;%
Bishop, \textit{Pattern Recognition and Machine Learning}, Springer, 2006.;%
Murphy, \textit{Machine Learning: A Probabilistic Perspective}, MIT Press, 2012.;%
Han, Kamber and Pei, \textit{Data Mining: Concepts and Techniques} (3e), Morgan Kaufmann, 2012.%
}

% --- unit-wise ---------------------------------------------------------------
\newcommand{\AMLUnitISources}{%
M\"uller and Guido, \textit{Introduction to Machine Learning with Python}, Ch. 1.;%
Bishop, \textit{Pattern Recognition and Machine Learning}, Ch. 1.;%
Murphy, \textit{Probabilistic Machine Learning: An Introduction}, MIT Press, 2022, Ch. 1.;%
Mitchell, \textit{Machine Learning}, McGraw Hill, 1997, Ch. 1.;%
Scikit-learn documentation, accessed 2026-08-04.%
}

\newcommand{\AMLUnitIISources}{%
Bishop, \textit{Pattern Recognition and Machine Learning}, \S1.5, \S4.3, \S7.1.;%
Zhang, Lipton, Li and Smola, \textit{Dive into Deep Learning}, \S3.4.;%
Shalev-Shwartz and Ben-David, \textit{Understanding Machine Learning}, Ch. 15.;%
Murphy, \textit{Probabilistic Machine Learning: An Introduction}, Ch. 5.%
}

\newcommand{\AMLUnitIIISources}{%
Zhang, Lipton, Li and Smola, \textit{Dive into Deep Learning}, Ch. 12 (Optimization Algorithms).;%
Deisenroth, Faisal and Ong, \textit{Mathematics for Machine Learning}, Ch. 7.;%
Kingma and Ba, ``Adam: A Method for Stochastic Optimization,'' ICLR 2015.;%
Ruder, ``An overview of gradient descent optimization algorithms,'' arXiv:1609.04747, 2016.%
}

\newcommand{\AMLUnitIVSources}{%
Han, Kamber and Pei, \textit{Data Mining: Concepts and Techniques} (3e), Ch. 3.;%
M\"uller and Guido, \textit{Introduction to Machine Learning with Python}, Ch. 3--4.;%
James, Witten, Hastie, Tibshirani and Taylor, \textit{An Introduction to Statistical Learning with Python}, Ch. 5--6, 12.;%
Scikit-learn documentation (preprocessing, model selection), accessed 2026-08-04.%
}

\newcommand{\AMLUnitVSources}{%
James et al., \textit{An Introduction to Statistical Learning with Python}, Ch. 3--6.;%
Bishop, \textit{Pattern Recognition and Machine Learning}, Ch. 3.;%
Deisenroth, Faisal and Ong, \textit{Mathematics for Machine Learning}, Ch. 9.;%
M\"uller and Guido, \textit{Introduction to Machine Learning with Python}, \S2.3.%
}

\newcommand{\AMLUnitVISources}{%
Han, Kamber and Pei, \textit{Data Mining: Concepts and Techniques} (3e), Ch. 8--9.;%
James et al., \textit{An Introduction to Statistical Learning with Python}, Ch. 8--9.;%
Bishop, \textit{Pattern Recognition and Machine Learning}, Ch. 4--5, 7, 14.;%
Zhang et al., \textit{Dive into Deep Learning}, Ch. 5.%
}

\newcommand{\AMLUnitVIISources}{%
Han, Kamber and Pei, \textit{Data Mining: Concepts and Techniques} (3e), Ch. 10.;%
James et al., \textit{An Introduction to Statistical Learning with Python}, Ch. 12.;%
M\"uller and Guido, \textit{Introduction to Machine Learning with Python}, \S3.5.;%
Ester, Kriegel, Sander and Xu, ``A density-based algorithm for discovering clusters (DBSCAN),'' KDD 1996.%
}

\newcommand{\AMLLabSources}{%
M\"uller and Guido, \textit{Introduction to Machine Learning with Python}, companion notebooks.;%
McKinney, \textit{Python for Data Analysis} (3e), O'Reilly, 2022.;%
Scikit-learn user guide, accessed 2026-08-04.;%
Matplotlib and Seaborn documentation, accessed 2026-08-04.%
}
%
    }{%
      % Marking schemes live at `private/solutions/<family>/<id>/latex/`.
      \IfFileExists{../../../../../public/_shared/aml-sources.tex}{%
        % Source strings for Applied Machine Learning (CSAI2017P) — used by slides + notes.
% Prescribed texts and reference books are taken from the course syllabus.
% Support texts are the author-hosted open-access editions:
%   statlearning.com            (ISL with Python)
%   probml.github.io/pml-book   (Murphy, Probabilistic Machine Learning)
%   d2l.ai                      (Dive into Deep Learning)
%   mml-book.github.io          (Mathematics for Machine Learning)

\newcommand{\AMLRepoURL}{https://github.com/tali7c/Applied-Machine-Learning}

% --- prescribed -------------------------------------------------------------
\newcommand{\AMLPrescribed}{%
M\"uller and Guido, \textit{Introduction to Machine Learning with Python}, O'Reilly, 2016.;%
Bishop, \textit{Pattern Recognition and Machine Learning}, Springer, 2006.;%
Murphy, \textit{Machine Learning: A Probabilistic Perspective}, MIT Press, 2012.;%
Han, Kamber and Pei, \textit{Data Mining: Concepts and Techniques} (3e), Morgan Kaufmann, 2012.%
}

% --- unit-wise ---------------------------------------------------------------
\newcommand{\AMLUnitISources}{%
M\"uller and Guido, \textit{Introduction to Machine Learning with Python}, Ch. 1.;%
Bishop, \textit{Pattern Recognition and Machine Learning}, Ch. 1.;%
Murphy, \textit{Probabilistic Machine Learning: An Introduction}, MIT Press, 2022, Ch. 1.;%
Mitchell, \textit{Machine Learning}, McGraw Hill, 1997, Ch. 1.;%
Scikit-learn documentation, accessed 2026-08-04.%
}

\newcommand{\AMLUnitIISources}{%
Bishop, \textit{Pattern Recognition and Machine Learning}, \S1.5, \S4.3, \S7.1.;%
Zhang, Lipton, Li and Smola, \textit{Dive into Deep Learning}, \S3.4.;%
Shalev-Shwartz and Ben-David, \textit{Understanding Machine Learning}, Ch. 15.;%
Murphy, \textit{Probabilistic Machine Learning: An Introduction}, Ch. 5.%
}

\newcommand{\AMLUnitIIISources}{%
Zhang, Lipton, Li and Smola, \textit{Dive into Deep Learning}, Ch. 12 (Optimization Algorithms).;%
Deisenroth, Faisal and Ong, \textit{Mathematics for Machine Learning}, Ch. 7.;%
Kingma and Ba, ``Adam: A Method for Stochastic Optimization,'' ICLR 2015.;%
Ruder, ``An overview of gradient descent optimization algorithms,'' arXiv:1609.04747, 2016.%
}

\newcommand{\AMLUnitIVSources}{%
Han, Kamber and Pei, \textit{Data Mining: Concepts and Techniques} (3e), Ch. 3.;%
M\"uller and Guido, \textit{Introduction to Machine Learning with Python}, Ch. 3--4.;%
James, Witten, Hastie, Tibshirani and Taylor, \textit{An Introduction to Statistical Learning with Python}, Ch. 5--6, 12.;%
Scikit-learn documentation (preprocessing, model selection), accessed 2026-08-04.%
}

\newcommand{\AMLUnitVSources}{%
James et al., \textit{An Introduction to Statistical Learning with Python}, Ch. 3--6.;%
Bishop, \textit{Pattern Recognition and Machine Learning}, Ch. 3.;%
Deisenroth, Faisal and Ong, \textit{Mathematics for Machine Learning}, Ch. 9.;%
M\"uller and Guido, \textit{Introduction to Machine Learning with Python}, \S2.3.%
}

\newcommand{\AMLUnitVISources}{%
Han, Kamber and Pei, \textit{Data Mining: Concepts and Techniques} (3e), Ch. 8--9.;%
James et al., \textit{An Introduction to Statistical Learning with Python}, Ch. 8--9.;%
Bishop, \textit{Pattern Recognition and Machine Learning}, Ch. 4--5, 7, 14.;%
Zhang et al., \textit{Dive into Deep Learning}, Ch. 5.%
}

\newcommand{\AMLUnitVIISources}{%
Han, Kamber and Pei, \textit{Data Mining: Concepts and Techniques} (3e), Ch. 10.;%
James et al., \textit{An Introduction to Statistical Learning with Python}, Ch. 12.;%
M\"uller and Guido, \textit{Introduction to Machine Learning with Python}, \S3.5.;%
Ester, Kriegel, Sander and Xu, ``A density-based algorithm for discovering clusters (DBSCAN),'' KDD 1996.%
}

\newcommand{\AMLLabSources}{%
M\"uller and Guido, \textit{Introduction to Machine Learning with Python}, companion notebooks.;%
McKinney, \textit{Python for Data Analysis} (3e), O'Reilly, 2022.;%
Scikit-learn user guide, accessed 2026-08-04.;%
Matplotlib and Seaborn documentation, accessed 2026-08-04.%
}
%
      }{%
        % Source strings for Applied Machine Learning (CSAI2017P) — used by slides + notes.
% Prescribed texts and reference books are taken from the course syllabus.
% Support texts are the author-hosted open-access editions:
%   statlearning.com            (ISL with Python)
%   probml.github.io/pml-book   (Murphy, Probabilistic Machine Learning)
%   d2l.ai                      (Dive into Deep Learning)
%   mml-book.github.io          (Mathematics for Machine Learning)

\newcommand{\AMLRepoURL}{https://github.com/tali7c/Applied-Machine-Learning}

% --- prescribed -------------------------------------------------------------
\newcommand{\AMLPrescribed}{%
M\"uller and Guido, \textit{Introduction to Machine Learning with Python}, O'Reilly, 2016.;%
Bishop, \textit{Pattern Recognition and Machine Learning}, Springer, 2006.;%
Murphy, \textit{Machine Learning: A Probabilistic Perspective}, MIT Press, 2012.;%
Han, Kamber and Pei, \textit{Data Mining: Concepts and Techniques} (3e), Morgan Kaufmann, 2012.%
}

% --- unit-wise ---------------------------------------------------------------
\newcommand{\AMLUnitISources}{%
M\"uller and Guido, \textit{Introduction to Machine Learning with Python}, Ch. 1.;%
Bishop, \textit{Pattern Recognition and Machine Learning}, Ch. 1.;%
Murphy, \textit{Probabilistic Machine Learning: An Introduction}, MIT Press, 2022, Ch. 1.;%
Mitchell, \textit{Machine Learning}, McGraw Hill, 1997, Ch. 1.;%
Scikit-learn documentation, accessed 2026-08-04.%
}

\newcommand{\AMLUnitIISources}{%
Bishop, \textit{Pattern Recognition and Machine Learning}, \S1.5, \S4.3, \S7.1.;%
Zhang, Lipton, Li and Smola, \textit{Dive into Deep Learning}, \S3.4.;%
Shalev-Shwartz and Ben-David, \textit{Understanding Machine Learning}, Ch. 15.;%
Murphy, \textit{Probabilistic Machine Learning: An Introduction}, Ch. 5.%
}

\newcommand{\AMLUnitIIISources}{%
Zhang, Lipton, Li and Smola, \textit{Dive into Deep Learning}, Ch. 12 (Optimization Algorithms).;%
Deisenroth, Faisal and Ong, \textit{Mathematics for Machine Learning}, Ch. 7.;%
Kingma and Ba, ``Adam: A Method for Stochastic Optimization,'' ICLR 2015.;%
Ruder, ``An overview of gradient descent optimization algorithms,'' arXiv:1609.04747, 2016.%
}

\newcommand{\AMLUnitIVSources}{%
Han, Kamber and Pei, \textit{Data Mining: Concepts and Techniques} (3e), Ch. 3.;%
M\"uller and Guido, \textit{Introduction to Machine Learning with Python}, Ch. 3--4.;%
James, Witten, Hastie, Tibshirani and Taylor, \textit{An Introduction to Statistical Learning with Python}, Ch. 5--6, 12.;%
Scikit-learn documentation (preprocessing, model selection), accessed 2026-08-04.%
}

\newcommand{\AMLUnitVSources}{%
James et al., \textit{An Introduction to Statistical Learning with Python}, Ch. 3--6.;%
Bishop, \textit{Pattern Recognition and Machine Learning}, Ch. 3.;%
Deisenroth, Faisal and Ong, \textit{Mathematics for Machine Learning}, Ch. 9.;%
M\"uller and Guido, \textit{Introduction to Machine Learning with Python}, \S2.3.%
}

\newcommand{\AMLUnitVISources}{%
Han, Kamber and Pei, \textit{Data Mining: Concepts and Techniques} (3e), Ch. 8--9.;%
James et al., \textit{An Introduction to Statistical Learning with Python}, Ch. 8--9.;%
Bishop, \textit{Pattern Recognition and Machine Learning}, Ch. 4--5, 7, 14.;%
Zhang et al., \textit{Dive into Deep Learning}, Ch. 5.%
}

\newcommand{\AMLUnitVIISources}{%
Han, Kamber and Pei, \textit{Data Mining: Concepts and Techniques} (3e), Ch. 10.;%
James et al., \textit{An Introduction to Statistical Learning with Python}, Ch. 12.;%
M\"uller and Guido, \textit{Introduction to Machine Learning with Python}, \S3.5.;%
Ester, Kriegel, Sander and Xu, ``A density-based algorithm for discovering clusters (DBSCAN),'' KDD 1996.%
}

\newcommand{\AMLLabSources}{%
M\"uller and Guido, \textit{Introduction to Machine Learning with Python}, companion notebooks.;%
McKinney, \textit{Python for Data Analysis} (3e), O'Reilly, 2022.;%
Scikit-learn user guide, accessed 2026-08-04.;%
Matplotlib and Seaborn documentation, accessed 2026-08-04.%
}
%
      }%
    }%
  }%
}
%
  }%
}

\usepackage{fancyhdr}
\usepackage{tcolorbox}
\tcbuselibrary{skins,breakable}

\usepackage{array}
\usepackage[htt]{hyphenat}   % allow \texttt{...} to hyphenate, else long
                             % identifiers overflow the measure

\hypersetup{colorlinks=true, linkcolor=blue, urlcolor=blue, breaklinks=true}
\setlist{itemsep=0.35em, topsep=0.35em}
\sloppy
\emergencystretch=2em

\providecommand{\CourseCode}{CSAI2017P: Applied Machine Learning (Lab)}
\providecommand{\AssignmentName}{Lab Assignment}

\pagestyle{fancy}
\fancyhf{}
\lhead{\small\CourseCode}
\rhead{\small\AssignmentName}
\cfoot{\thepage}
\setlength{\headheight}{14pt}

% Per-question mark tag, right aligned.
\newcommand{\QMarks}[1]{\hfill \textbf{[#1 marks]}}

% Title block for a brief.
\newcommand{\LAtitle}[4]{%
  \begin{center}
    {\LARGE \textbf{#1}}\\[0.25em]
    {\large \CourseCode}\\[0.25em]
    \normalsize #2
  \end{center}
  \noindent
  \textbf{Instructor:} Dr.\ Tofik Ali \hfill \textbf{Due:} #3\\
  \textbf{Student Name:} \rule{5cm}{0.4pt} \hfill \textbf{SAP ID:} \rule{3.5cm}{0.4pt}\\
  \textbf{Batch:} \rule{2.5cm}{0.4pt} \hfill \textbf{Lab quiz:} #4
  \vspace{0.6em}\hrule\vspace{0.8em}}

% Title block for a marking scheme (no student fields).
\newcommand{\MStitle}[3]{%
  \begin{center}
    {\LARGE \textbf{#1}}\\[0.25em]
    {\large \CourseCode}\\[0.25em]
    \normalsize #2\\[0.4em]
    {\color{red}\textbf{MARKING SCHEME --- NOT FOR CIRCULATION}}
  \end{center}
  \noindent \textbf{Instructor:} Dr.\ Tofik Ali \hfill \textbf{Assignment due:} #3
  \vspace{0.6em}\hrule\vspace{0.8em}}

\newtcolorbox{infobox}[1][Note]{
  breakable, enhanced, colback=gray!6, colframe=gray!55,
  fonttitle=\bfseries\small, title={#1}, boxrule=0.7pt, arc=2pt,
  left=6pt,right=6pt,top=4pt,bottom=4pt}

\newtcolorbox{answerbox}[1][Expected answer]{
  breakable, enhanced, colback=green!4, colframe=green!35!black,
  fonttitle=\bfseries\small, title={#1}, boxrule=0.7pt, arc=2pt,
  left=6pt,right=6pt,top=4pt,bottom=4pt}


\newcommand{\A}[1]{\item[\textbf{#1}]}
\newcommand{\NUMT}{{\footnotesize\textbf{[NUM]}}}
\setlist[enumerate]{itemsep=0.4em, topsep=0.35em}

\begin{document}

\MStitle{Mid-Semester Question Bank --- Answers}{Units I--IV and Unit V (part) \textbullet\ L1--L10 \textbullet\ 314 questions}{}

\begin{infobox}[How to use this]
These are the answers to all 314 questions in the bank, with the working for
every numerical one.

\textbf{Work the question first.} Reading a solution produces the feeling of
understanding and almost none of the substance --- that is why the bank told
you to do the numericals with a pen, and it is still true now that the answers
are in your hands.

\textbf{Numerical values are exact to the digits shown.} If you rounded
intermediate steps you may differ in the last decimal or two; that is rounding,
not an error. If you differ by more than that, the working is here so you can
find where the two of you parted company.

\textbf{Theory answers say what has to be present, not words to match.} Your
wording will differ from the wording here, and that is expected. What matters
is whether the named ideas are in your answer.

\textbf{Statements to judge} carry a verdict and the reason behind it. The
reason is the part that does the work --- a verdict on its own does not answer
the question that was asked.

Several questions have more than one good route, and a route different from the
printed one is not automatically worse. If you reached a different answer and
can show your reasoning, bring it to office hours.
\end{infobox}

\noindent If you change a question's numbers, change them in the script first
and paste the new output here. A key that has drifted from its source is worse
than no key.

\vspace{0.5em}\hrule

% =====================================================================
\section*{Unit I --- Introduction and the Python ML stack}

\subsection*{§A --- Two-mark questions}

\begin{enumerate}[leftmargin=1.9cm]

\A{U1.1} ``A program learns from experience $E$ with respect to a class of
tasks $T$ and performance measure $P$, if its performance at $T$, as measured
by $P$, improves with $E$.'' For pass prediction: $T$ = classify a student as
pass/fail; $E$ = records of past students with their outcomes; $P$ = accuracy
or recall on students not in that record.

\A{U1.2} (a) Supervised regression --- signal is the recorded rainfall for each
past month. (b) Unsupervised --- no target; the signal is structure in the
feature space. (c) Reinforcement --- a delayed scalar reward (the score), not a
per-step label. \emph{0.5 each classification, 0.5 for naming signals.}

\A{U1.3} $\text{AI} \supset \text{ML} \supset \text{DL}$. Not right for deep
learning: a 500-row tabular credit dataset --- a gradient-boosted tree usually
matches or beats a neural network and trains in seconds.

\A{U1.4} Any three of: the rule is already known and writable; there are no
examples of the answer; the cost of an error is unbounded and the model cannot
explain itself; the data distribution at deployment differs from training;
a simple baseline already meets the requirement.

\A{U1.5} NumPy --- $n$-dimensional numeric arrays and vectorised arithmetic.
Pandas --- labelled heterogeneous tables, I/O, cleaning. scikit-learn ---
estimators, transformers, model selection.

\A{U1.6} \texttt{fit} learns parameters from data and stores them;
\texttt{transform} applies stored parameters; \texttt{predict} applies a
learned model. \texttt{fit\_transform} on test data would recompute the
parameters \emph{from} the test data --- the definition of leakage.
\emph{The last sentence is the mark.}

\A{U1.7} \texttt{head} (what the rows look like), \texttt{info} (dtypes and
non-null counts), \texttt{describe} (location, spread, range per column),
\texttt{isna().sum()} (where the holes are).

\A{U1.8} ``The axis I name is the one that disappears.'' $(3,4)$:
\texttt{sum(axis=0)} $\to (4,)$; \texttt{sum(axis=1)} $\to (3,)$.

\A{U1.9} Training error low and test error high. The diagnostic is the
\emph{gap} between them, not either number alone.

\A{U1.10} Use, in training or model selection, of information not available at
prediction time --- in practice, information from the evaluation data. Two
examples from: scaler/imputer fitted on everything; feature selection on
everything; resampling before the split; target encoding over all rows.

\A{U1.11} Once you compare models on the test set, you have selected using it,
and the winner's score is optimistically biased by the selection. The
validation set absorbs the selection; the test set is looked at once, at the
end.

\A{U1.12} Two decimal places on a score from 38 test rows implies a precision
the sample cannot support --- one row is worth 2.6 accuracy points. Report
``95\%'', or better, a cross-validated mean with its spread.

\A{U1.13} Need scaling: $k$-NN, RBF-kernel SVM, K-means, and any
gradient-trained linear model (distance or step size depends on units). Do not
need it: decision trees and random forests --- a split asks
``is $x_j \le t$?'', and a monotone rescaling $x \mapsto (x-a)/b$ with $b>0$
moves $t$ to $(t-a)/b$ and sends exactly the same rows left.

\end{enumerate}

\subsection*{§B --- Five-mark questions}

\begin{enumerate}[leftmargin=1.9cm]

\A{U1.14} \textbf{Verified.}
\begin{answerbox}
(a) Distances from $(4.6, 1.4)$:
\begin{center}\begin{tabular}{@{}lrl@{}}\toprule
point & distance & species\\ \midrule
$(1.5,0.3)$ & \textbf{3.2894} & setosa\\
$(1.4,0.2)$ & \textbf{3.4176} & setosa\\
$(4.8,1.5)$ & \textbf{0.2236} & versicolor\\
$(4.2,1.2)$ & \textbf{0.4472} & versicolor\\
$(5.9,2.4)$ & \textbf{1.6401} & virginica\\
\bottomrule\end{tabular}\end{center}
Working for the third: $\sqrt{0.2^2 + 0.1^2} = \sqrt{0.05} = 0.2236$.

(b) 1-NN $\to$ \textbf{versicolor} (0.2236 is smallest).

(c) Three nearest are versicolor, versicolor, virginica $\to$ vote 2--1 $\to$
\textbf{versicolor}.

(d) In millimetres, petal width is multiplied by 10 while length stays in cm.
Width differences then dominate the squared distance --- the metric is no
longer measuring what it was. Nothing about the flowers changed; only the
units did. This is the whole case for standardising before any distance-based
method.
\end{answerbox}

\A{U1.15} \textbf{Verified.}
\begin{answerbox}
(a) $e = (0.5, -0.5, 1.0, -1.0)$. $\text{MSE} = (0.25+0.25+1+1)/4 =
\mathbf{0.625}$; $\text{MAE} = 3/4 = \mathbf{0.75}$;
$\text{RMSE} = \sqrt{0.625} = \mathbf{0.7906}$.

(b) $\hat y_3 = 1.0 \Rightarrow e = (0.5, -0.5, 7.0, -1.0)$.
$\text{MSE} = (0.25+0.25+49+1)/4 = \mathbf{12.625}$;
$\text{MAE} = 9/4 = \mathbf{2.25}$;
$\text{RMSE} = \sqrt{12.625} = \mathbf{3.5532}$.

(c) MSE grew $\times 20.2$; MAE $\times 3.0$; RMSE $\times 4.49$. The error
went from 1 to 7, a factor of 7; MAE is linear so it grew by roughly that
factor spread over four points, MSE is quadratic so it grew by roughly $7^2$
spread over four, and RMSE --- being the square root of MSE --- sits between
them and is back in the units of $y$.
\end{answerbox}

\A{U1.16} \textbf{Verified.} (a) $0.2 \times 569 = \mathbf{114}$ test,
$\mathbf{455}$ train. (b) 10-fold on 569: validation $\approx \mathbf{57}$
rows, training $\approx \mathbf{512}$. (c) \textbf{10} fits. (d)
$n(k-1) = 569 \times 9 = \mathbf{5121}$ row-fits.
\emph{1.25 each.}

\A{U1.17} (a) \textbf{Degree 3} --- lowest test RMSE (1.10), which is the only
column that estimates future performance. (b) Degree 12 has driven training
error to 0.57 by fitting the noise; test error explodes to 40.52, so the fitted
curve is wild between the training points. (c) Training error can always be
driven to zero by adding capacity --- it measures memorisation, not
generalisation, so it cannot arbitrate between models. (d) Training curve falls
monotonically; test curve is U-shaped with its minimum at degree 3; everything
to the right of that minimum is the overfitting region.

\A{U1.18} (a) \textbf{0.50} --- the labels are coin flips and the features are
noise, so no procedure can beat chance in expectation. (b) The selector
computed its $F$-statistics using every row's target, including rows that later
became validation folds. Among 5000 pure-noise columns, roughly 20 will
correlate with $y$ across the whole sample by chance alone; the selector finds
exactly those, and the correlation it exploited is still present in the
validation rows because those rows helped choose the columns. The classifier
then ``confirms'' a pattern that the selection step planted. (c) Move
\texttt{SelectKBest} inside a \texttt{Pipeline} with the classifier and
cross-validate the pipeline.

\A{U1.19} Expect: \texttt{load\_*}; \texttt{train\_test\_split(test\_size=0.25,
random\_state=0, stratify=y)}; \texttt{make\_pipeline(StandardScaler(),
KNeighborsClassifier(5))}; \texttt{.fit(X\_train, y\_train)};
\texttt{.score(X\_test, y\_test)}. The two leakage-preventing lines are
\texttt{stratify=y}/\texttt{random\_state} (an honest, reproducible split) and
the pipeline construction (the scaler is fitted on training folds only).

\A{U1.20} A Python loop executes interpreted bytecode and boxes every element
as a Python object; NumPy dispatches once into a compiled C loop over a
contiguous typed buffer, with no per-element interpreter overhead and with SIMD
available. The loop wins when the operation is not expressible elementwise,
when it must exit early, or when the array is small enough that NumPy's
call overhead dominates.

\end{enumerate}

\subsection*{§C --- Ten-mark questions}

\noindent\small A complete §C answer covers the listed points, arrives somewhere rather than trailing off, and carries at least one concrete numerical illustration. An answer with no numbers anywhere
caps at 6 of 10 regardless of prose quality.\normalsize

\begin{enumerate}[leftmargin=1.9cm]

\A{U1.21} Must cover: training score measures memorisation and can be driven to
zero (the degree-12 row); a held-out set estimates performance on unseen rows;
CV averages over $k$ disjoint held-out sets and so has lower variance than one
split; three-way split --- train fits, validation selects, test estimates once;
two bias mechanisms, e.g.\ selecting a model on the test set, and any
preprocessing step fitted before the split. Numerical illustration: the
overfitting table, or the 0.81-on-noise leak.

\A{U1.22} Expected chain: inspect (\texttt{head/info/describe/isna}) $\to$
split with \texttt{stratify} and a seed $\to$ \texttt{ColumnTransformer} with
imputers, scaler, encoder $\to$ model inside a \texttt{Pipeline} $\to$
\texttt{StratifiedKFold} cross-validation on the training portion $\to$ tune on
validation or by nested CV $\to$ one look at the test set $\to$ report a mean
with its spread. At each stage the failure: unseen categories
(\texttt{handle\_unknown}), \texttt{NaN} reaching the scaler, statistics
computed over test rows, selecting on the test set. The pipeline prevents the
middle three by construction.

\A{U1.23} Rule-based: written by hand, adapts only when a human edits it,
debugged by reading, fully explainable, degrades predictably. Learned: induced
from labelled examples, adapts by retraining, debugged by examining data and
errors, explainable only with effort, can fail silently under distribution
shift. Choose rules when the rule is known, stable, and must be auditable;
choose learning when you have examples of the answer but no statement of the
rule.

\A{U1.24} The uniform API means any estimator can be dropped into any
pipeline, any cross-validator, any grid search, without those tools knowing
what the estimator is. Concretely, it makes the leakage bug impossible: because
a \texttt{Pipeline} is itself an estimator with \texttt{fit}/\texttt{predict},
\texttt{cross\_val\_score} refits the \emph{whole} chain inside each fold, so a
scaler cannot accidentally be fitted once outside the loop.

\end{enumerate}

\vspace{0.4em}\hrule

% =====================================================================
\section*{Unit II --- Loss functions}

\subsection*{§A --- Two-mark questions}

\begin{enumerate}[leftmargin=1.9cm]
\A{U2.1} $\text{MSE} = \frac1n\sum r_i^2$, $\text{MAE} = \frac1n\sum|r_i|$,
$\text{RMSE} = \sqrt{\text{MSE}}$. MSE is most sensitive: doubling a residual
quadruples its contribution; equivalently $\partial\text{MSE}/\partial\hat y$
grows linearly with the residual while MAE's is constant $\pm 1$.

\A{U2.2} \emph{Loss} --- the penalty on one example. \emph{Cost/objective} ---
the aggregate the optimiser minimises. \emph{Metric} --- the number reported to
a human, which need not be differentiable. Example: train with cross-entropy,
report $F_1$; train with Huber, report RMSE.

\A{U2.3} It must be minimised by a correct prediction, and it must be
differentiable (or subdifferentiable) so a gradient method can move on it.

\A{U2.4} Squared $\to$ the \textbf{mean}. Absolute $\to$ the \textbf{median}.
0--1 $\to$ the \textbf{mode}.

\A{U2.5} $L_\delta(r) = \tfrac12 r^2$ for $|r| \le \delta$, else
$\delta(|r| - \tfrac12\delta)$. The second branch is the tangent to the first
at $r = \delta$: written that way the value \emph{and} the slope match at the
join, so the loss is differentiable everywhere. Plain $\delta|r|$ would match
neither.

\A{U2.6} $\delta$ is in the units of $y$. Wrong way: copying a value from a
tutorial. Right: cross-validation, or a robust estimate of the residual scale.

\A{U2.7} It is piecewise constant --- zero gradient almost everywhere, so
nothing to descend --- and it takes very few distinct values on a small batch,
so it cannot distinguish a barely-right model from a confidently-right one.

\A{U2.8} $\text{BCE} = -\frac1n\sum_i [y_i\log p_i + (1-y_i)\log(1-p_i)]$.
For a single example it is $-\log(\text{probability given to the correct
class})$.

\A{U2.9} Categorical expects one-hot labels, shape $(n,K)$; sparse expects
integer indices, shape $(n,)$. They give \textbf{identical} numbers ---
2-D means categorical, 1-D means sparse.

\A{U2.10} Memory and speed. $K = 10{,}000$ with a batch of 256 means a one-hot
matrix of 2.56 million numbers of which 256 are non-zero; the sparse form
indexes directly and never builds it.

\A{U2.11} $\log 0$ is undefined and $-\log$ of a value near zero overflows to
$\infty$, which propagates \texttt{nan} through the gradients. Clipping bounds
the loss without materially changing it.

\A{U2.12} $m_i = y_i f(x_i)$ --- positive when correct, and its magnitude is
the confidence. Hinge $= \max(0, 1 - m_i)$.

\A{U2.13} Hinge becomes \emph{exactly} zero once $m \ge 1$ and stays there;
logistic loss approaches zero but never arrives. Consequence: only points with
$m < 1$ contribute gradient, so an SVM's solution is determined by a handful of
points --- the support vectors.

\A{U2.14} $L = \max(0,\ d(a,p) - d(a,n) + \alpha)$ on an anchor, a positive and
a negative. $\alpha$ is the margin: how much closer the positive must be than
the negative before the triplet costs nothing.

\A{U2.15} Most random triplets already satisfy the margin, giving zero loss and
zero gradient, so nearly all the compute is wasted. Remedy: \textbf{triplet
mining} --- selecting hard or semi-hard triplets.

\A{U2.16} A differentiable, tractable loss that upper-bounds the loss you
actually care about. The two standard surrogates for 0--1 are cross-entropy
(logistic) and hinge.

\A{U2.17} (a) Huber, or MAE. (b) Sparse categorical cross-entropy.
(c) Triplet loss. (d) Binary cross-entropy.
\emph{0.5 each.}
\end{enumerate}

\subsection*{§B --- Five-mark questions}

\begin{enumerate}[leftmargin=1.9cm]

\A{U2.18} \textbf{Verified.}
\begin{answerbox}
(a) $r = y - \hat y = (-0.5,\ 1.0,\ -2.5,\ 1.0,\ 0)$;
$r^2 = (0.25,\ 1,\ 6.25,\ 1,\ 0)$, $\sum = 8.5$;
$|r| = (0.5,\ 1,\ 2.5,\ 1,\ 0)$, $\sum = 5$.

(b) $\text{MSE} = 8.5/5 = \mathbf{1.7}$; $\text{MAE} = 5/5 = \mathbf{1.0}$;
$\text{RMSE} = \sqrt{1.7} = \mathbf{1.3038}$.

(c) Worst point $r = -2.5$: share of SSE $= 6.25/8.5 = \mathbf{0.7353}$;
share of SAE $= 2.5/5 = \mathbf{0.5000}$.

(d) One point in five carries 74\% of the squared error but only 50\% of the
absolute error. Under MSE that single row effectively decides the fit; under
MAE it is one voice among five. This is the whole difference between the two
losses, in two numbers.
\end{answerbox}

\A{U2.19} \textbf{Verified.} Residuals as above.
\begin{answerbox}
\begin{center}\begin{tabular}{@{}lllr@{}}\toprule
$\delta$ & per-point loss & quadratic branch & mean\\
\midrule
1.0 & $0.125,\ 0.5,\ 2.0,\ 0.5,\ 0$ & all but $r=-2.5$ & \textbf{0.625}\\
1.5 & $0.125,\ 0.5,\ 2.625,\ 0.5,\ 0$ & all but $r=-2.5$ & \textbf{0.75}\\
3.0 & $0.125,\ 0.5,\ 3.125,\ 0.5,\ 0$ & \emph{all five} & \textbf{0.85}\\
\bottomrule\end{tabular}\end{center}
Check at $\delta = 1.5$: $1.5 \times (2.5 - 0.75) = 2.625$.
At $\delta = 3$ every $|r| \le 3$, so every point takes $\tfrac12 r^2$ and the
mean is $8.5/(2\times5) = 0.85$ --- exactly half the MSE.

As $\delta \to \infty$ the Huber loss becomes $\tfrac12 \times$ MSE: no point
ever reaches the linear branch, and all robustness is lost. As $\delta \to 0$
it becomes $\delta \times$ MAE.
\end{answerbox}

\A{U2.20} \textbf{Verified.}
\begin{answerbox}
(a) Clean: mean $= 146/6 = \mathbf{24.3333}$, median $= (24+25)/2 =
\mathbf{24.5}$.

(b) With 148: mean $= 266/6 = \mathbf{44.3333}$, median $= \mathbf{24.5}$
--- unmoved.

(c) Squared error is minimised by the mean, $\mathbf{44.3333}$; absolute error
by the median. A grid search over $[0,200]$ returns $\mathbf{24.0}$ for the
absolute error rather than 24.5, because with an even $n$ the absolute loss is
\emph{flat} on the whole interval $[24, 25]$ and \texttt{argmin} returns the
left end; every point in that interval is a minimiser.

(d) The constant a loss chooses \emph{is} what the model becomes when it has no
features to work with. A model trained on MSE is a model of the average case
and inherits the mean's zero breakdown point; a model trained on MAE is a model
of the typical case. The one corrupted mark moved the MSE answer by 20 marks
and the MAE answer by nothing.
\end{answerbox}

\A{U2.21} \textbf{Verified.}
\begin{answerbox}
(a) Probability on the correct class: $(0.8,\ 0.7,\ 0.6,\ 0.9)$.

(b) $-\log$: $(0.2231,\ 0.3567,\ 0.5108,\ 0.1054)$; sum $1.1960$;
$\text{BCE} = \mathbf{0.2990}$.

(c) With $p_3 = 0.02$: $-\log 0.02 = 3.9120$, so the sum becomes
$0.2231+0.3567+3.9120+0.1054 = 4.5972$ and
$\text{BCE} = \mathbf{1.1493}$ --- an increase of $\mathbf{284\%}$.

(d) The loss is a negative logarithm, which is unbounded as $p \to 0$. A
prediction of 0.01 on the correct class costs 4.61 against 0.01 for a
prediction of 0.99 --- a factor of \textbf{458}. One confidently wrong row can
therefore outweigh several hundred quietly correct ones.
\end{answerbox}

\A{U2.22} \textbf{Verified.}
\begin{answerbox}
(a) Correct-class probabilities are the diagonal: $0.6,\ 0.7,\ 0.6$.
$-\log$: $0.5108,\ 0.3567,\ 0.5108$; mean $= \mathbf{0.459442}$.

(b) Identical: $\mathbf{0.459442}$.

(c) They must agree --- one-hot labels zero out every term except the correct
class, which is exactly what integer indexing selects. Same loss, two label
representations.

(d) One-hot for $K = 10{,}000$ and a batch of 256 is
$2{,}560{,}000$ numbers of which \textbf{256} are non-zero; the sparse form
stores 256 integers. Roughly a $10{,}000\times$ saving on the label tensor.
\end{answerbox}

\A{U2.23} \textbf{Verified.}
\begin{answerbox}
\begin{center}\begin{tabular}{@{}rrr@{}}\toprule
$p$ & $-\log p$ & ratio to $p=0.99$\\
\midrule
0.99 & 0.0101 & 1.0\\ 0.90 & 0.1054 & 10.5\\ 0.50 & \textbf{0.6931} & 69.0\\
0.10 & 2.3026 & 229.1\\ 0.01 & 4.6052 & \textbf{458.2}\\
\bottomrule\end{tabular}\end{center}
The $p = 0.50$ row is exactly $\log 2 = 0.6931$ --- the loss of a model that
has learnt nothing and says so honestly. It is the natural baseline for binary
cross-entropy, and a training run whose loss sits at 0.693 has not started.
\end{answerbox}

\A{U2.24} \textbf{Verified.}
\begin{answerbox}
(a) $m = y f = (1.5,\ 0.4,\ 0.2,\ 0.9)$ --- all positive.

(b) Hinge $= \max(0, 1-m) = (0,\ 0.6,\ 0.8,\ 0.1)$; mean $= 1.5/4 =
\mathbf{0.375}$.

(c) Squared hinge $= (0,\ 0.36,\ 0.64,\ 0.01)$; mean $= 1.01/4 =
\mathbf{0.2525}$.

(d) Every margin is positive, so every point is on the correct side of the
boundary --- but three of them sit inside the margin band $m < 1$. Hinge loss
does not ask for correctness; it asks for correctness \emph{by a stated
margin}. Only the first point, at $m = 1.5$, has bought enough room to stop
contributing. This is why an SVM's solution depends on so few points: everything
with $m \ge 1$ has zero loss and zero gradient and drops out of the problem.
\end{answerbox}

\A{U2.25} \textbf{Verified.} (a) $\max(0, 0.9-1.6+1.0) = \mathbf{0.3}$ ---
non-zero, so there is a gradient. (b) $\max(0, 0.4-2.1+1.0) = \max(0,-0.7) =
\mathbf{0}$ --- the margin is already satisfied, no gradient, the triplet
teaches nothing. (c) $\max(0, 1.2-1.4+0.5) = \mathbf{0.3}$ --- non-zero even
though the ordering is correct, because the gap 0.2 is below $\alpha = 0.5$.
Case (b) is the common case under random sampling, which is why almost all
compute is wasted without mining.
\emph{1 each + 2 for the sampling conclusion.}

\A{U2.26} The two numbers are not comparable: MAE and RMSE are different
functions of the same residuals, and MAE $\le$ RMSE always. On the L2 worked
example, identical predictions gave MAE 1.0 and MSE 1.375 --- neither model was
better; only the ruler changed. To compare the two models, evaluate both under
\emph{one} metric on \emph{one} held-out set, chosen for what the application
cares about.

\A{U2.27} $\partial\text{MSE}/\partial r = 2r$: continuous, and it shrinks as
the fit improves, so steps get gentler near the optimum.
$\partial\text{MAE}/\partial r = \text{sign}(r)$: constant magnitude, so the
step size never eases off, and it is undefined at $r = 0$ --- the optimiser
oscillates about the minimum. On an even-sized sample the absolute loss is flat
across the whole interval between the two middle values, so the gradient is
exactly zero on a set of minimisers and the answer returned depends on the
solver.

\A{U2.28} (a) BCE on a softmax treats the $K$ outputs as $K$ independent binary
problems; the softmax constraint $\sum_k p_k = 1$ makes them dependent, so the
loss is mis-specified and the gradients fight each other. (b) Categorical CE
expects shape $(n,4)$; integer labels of shape $(n,)$ either raise a shape
error or --- worse --- broadcast silently and compute nonsense. (c) Four
sigmoids do not sum to 1, so ``probabilities'' can total 3.2; categorical CE
divides by nothing and the loss is no longer a proper scoring rule.

\A{U2.29} All three surrogates upper-bound the 0--1 step, which is what makes
minimising them a sensible proxy. Hinge reaches \emph{exactly} zero at $m = 1$
and is flat thereafter; logistic approaches zero asymptotically and never
arrives, so it keeps pushing confident points further. For very negative
margins hinge grows linearly and squared hinge quadratically --- so squared
hinge is far more sensitive to a badly misclassified point, in the same way MSE
is more sensitive than MAE.

\end{enumerate}

\subsection*{§C --- Ten-mark questions}

\begin{enumerate}[leftmargin=1.9cm]

\A{U2.30} Required coverage: the three constant minimisers (mean / median /
mode) with the argument; a worked contaminated sample showing the mean move and
the median hold (U2.20 numbers are fine); Huber as the piecewise compromise,
$\delta$ in units of $y$, chosen by CV; a loss/metric divergence example. The
thesis to land: the loss determines what the model \emph{becomes}, not merely
how it is scored.

\A{U2.31} Coverage: BCE, CCE and sparse CCE with formulae, label formats and
output activations; the identity between the last two; $-\log$ of the correct
class as the unifying form; the confidently-wrong table with the $458\times$
figure and $\log 2 = 0.6931$ as the honest-ignorance baseline; why accuracy
cannot serve --- piecewise constant, few distinct values, no gradient.

\A{U2.32} A comparison table plus a worked five-point example. Expect: MSE
quadratic / outlier-sensitive / fits the mean / smooth / gradient $2r$ / no
hyperparameter / clean symmetric data. MAE linear / robust / fits the median /
kink at 0 / gradient $\pm 1$ / none / untrusted outliers. Huber piecewise /
bounded influence / between the two / smooth everywhere / one hyperparameter
$\delta$ in units of $y$ / unavoidable outliers.

\A{U2.33} Coverage: $m = yf$; hinge geometry and the margin band; support
vectors as the points with $m<1$; squared hinge and its quadratic tail; triplet
loss for metric learning; mining because random triplets give zero gradient.

\end{enumerate}

\subsection*{§D --- Statements}

\begin{enumerate}[leftmargin=1.9cm]
\A{U2.34} \textbf{False.} They are different functions of the same residuals
and MAE $\le$ RMSE by construction; the smaller number reflects the ruler, not
the model.
\A{U2.35} \textbf{True.} The linear branch is the tangent to the quadratic at
$r = \delta$, so value \emph{and} first derivative agree at the join.
\A{U2.36} \textbf{False.} Same loss, different label encoding; they return
identical values (U2.22: both 0.459442). Only memory differs.
\A{U2.37} \textbf{False.} Hinge is zero only for $m \ge 1$. U2.24 has four
correctly classified points and three of them cost something.
\A{U2.38} \textbf{False on the second clause.} RMSE is a monotone function of
MSE, so on \emph{one} dataset the ranking is identical --- but the claim as
written is loose: RMSE is in the units of $y$ while MSE is in units of $y^2$,
which matters whenever the two are compared across differently scaled targets
or combined with another term.
\end{enumerate}

\vspace{0.4em}\hrule

% =====================================================================
\section*{Unit III --- Optimiser functions}

\subsection*{§A --- Two-mark questions}

\begin{enumerate}[leftmargin=1.9cm]
\A{U3.1} $w \leftarrow w - \eta\,\nabla J(w)$, where $w$ is the parameter
vector, $\eta > 0$ the learning rate and
$\nabla J = [\partial J/\partial w_1, \dots, \partial J/\partial w_p]^\top$.

\A{U3.2} Any two: the closed form costs $O(np^2 + p^3)$ and is infeasible for
large $p$; $X^\top X$ may be singular or badly conditioned; most losses of
interest have no closed form at all; gradient methods stream over data that
does not fit in memory.

\A{U3.3} An \emph{epoch} is one complete pass over the training set; an
\emph{update} is one application of the update rule. Batch: 1 update/epoch.
SGD: $n$. Mini-batch: $n/B$.

\A{U3.4} $0 < \eta < 2/L$. At $\eta = 1/L$ the multiplier $|1 - \eta f''|$ is
exactly zero, so a quadratic is solved in one step --- that step is
\textbf{Newton's method}.

\A{U3.5} Pick $i$ at random; $w \leftarrow w - \eta\nabla\ell_i(w)$. It works
because the per-example gradient is \emph{unbiased}:
$\mathbb E_i[\nabla\ell_i(w)] = \nabla J(w)$.

\A{U3.6} $\sum_t \eta_t = \infty$ (able to travel any distance) and
$\sum_t \eta_t^2 < \infty$ (noise eventually suppressed).
$\eta_t = \eta_0/(1+ct)$ satisfies both.

\A{U3.7} It falls as $1/\sqrt B$. So quality improves as $\sqrt B$ while cost
grows as $B$ --- doubling $B$ buys 29\% less noise for 100\% more work.

\A{U3.8} $v \leftarrow \beta v + \nabla J(w)$; $w \leftarrow w - \eta v$.
$\beta$ sets how much accumulated velocity is retained --- the memory of past
gradients. As $\beta \to 1$ the effective step $\eta/(1-\beta) \to \infty$:
oscillation across a ravine is damped but the iterate overshoots badly.

\A{U3.9} $\eta_{\text{eff}} \approx \eta/(1-\beta)$. Going from $0.9$ to
$0.99$ multiplies the effective step by \textbf{10}, so $\eta$ must be divided
by 10 to keep the same behaviour.

\A{U3.10} $\kappa = \lambda_{\max}/\lambda_{\min}$ of the Hessian. GD converges
at rate $(\kappa-1)/(\kappa+1) \approx 1 - 2/\kappa$; momentum improves this to
$\approx 1 - 2/\sqrt\kappa$.

\A{U3.11} On a well-conditioned problem ($\kappa \approx 1$) there is no
oscillation to damp, so the retained velocity only causes overshoot. On the
1-D quadratic, $\beta = 0.9$ was measured 2.5$\times$ \emph{slower} than
$\beta = 0$.

\A{U3.12} $s_t = s_{t-1} + g_t \odot g_t$;
$w \leftarrow w - \frac{\eta}{\sqrt{s_t}+\varepsilon} \odot g_t$.
$s_t$ accumulates the \emph{sum of squared gradients}, per coordinate.

\A{U3.13} The sum only ever grows, so the effective learning rate
$\eta/\sqrt{s_t} \approx \eta/(\bar g\sqrt t) \propto 1/\sqrt t$ decays without
bound and training stalls before convergence on long runs.

\A{U3.14} $s_t = \gamma s_{t-1} + (1-\gamma)g_t\odot g_t$; same second line.
The sum becomes an \textbf{exponentially weighted moving average}, so old
gradients decay out and the denominator can shrink again.

\A{U3.15} Adagrad's convergence guarantee. Its $1/\sqrt t$ decay \emph{is} a
Robbins--Monro schedule on a convex problem; RMSprop's average never settles,
so it tracks a moving target but has no such proof.

\A{U3.16} That the update should have the same units as the parameter it
updates. It removes $\eta$ entirely, replacing it with the RMS of recent
updates.

\A{U3.17} With $d_0 = 0$ the first numerator is $\sqrt\varepsilon$, so the
first step is proportional to $\sqrt\varepsilon$ and the whole trajectory is
set by it. A hyperparameter has been disguised as a stability constant.

\A{U3.18} $m_t = \beta_1 m_{t-1} + (1-\beta_1)g_t$;
$v_t = \beta_2 v_{t-1} + (1-\beta_2)g_t^2$;
$\hat m_t = m_t/(1-\beta_1^t)$, $\hat v_t = v_t/(1-\beta_2^t)$;
$w \leftarrow w - \eta\hat m_t/(\sqrt{\hat v_t}+\varepsilon)$.

\A{U3.19} $m_0 = v_0 = 0$, so for a constant gradient
$m_t = (1-\beta_1^t)g$ --- both moments are shrunk towards zero early on, and
by \emph{different} factors because $\beta_1 \ne \beta_2$. It is the mismatch,
not the shrinkage, that breaks the ratio.

\A{U3.20} At $t=1$: $\hat m_1 = (1-\beta_1)g_1/(1-\beta_1) = g_1$ and
$\hat v_1 = (1-\beta_2)g_1^2/(1-\beta_2) = g_1^2$. So
$\hat m_1/\sqrt{\hat v_1} = g_1/|g_1| = \pm 1$ and the step is exactly
$\pm\eta$.

\A{U3.21} Not a better final answer --- on a well-conditioned convex problem
with a tuned $\eta$ the optimiser is worth about one per cent. What they buy is
a \textbf{wider window of learning rates that work}: Adagrad worked over 2.24
decades of $\eta$ against SGD's 1.31.

\A{U3.22} \textbf{Curvature} (a ravine: coordinates want different step sizes
because of differing $\lambda_i$) and \textbf{frequency} (sparsity: a rare
feature is updated rarely). Momentum addresses the first and does nothing for
the second.

\A{U3.23} Still falling steeply $\to$ $\eta$ too small, multiply by 3. Falls
fast then flattens $\to$ healthy. Falls then bounces $\to$ $\eta$ slightly too
large, or the SGD noise floor. Rises or \texttt{nan} $\to$ past the stability
threshold, divide $\eta$ by 10.

\A{U3.24} That it is genuinely $B = 1$ --- true stochastic gradient descent,
not the mini-batch method that the name ``SGD'' usually denotes elsewhere.

\A{U3.25} $v \leftarrow \beta v + g$ (PyTorch, D2L) versus
$v \leftarrow \beta v + (1-\beta)g$ (Ng, the Adam paper). They differ by a
factor of $1-\beta$, so an $\eta$ tuned under one is wrong by $10\times$ under
the other at $\beta = 0.9$.
\end{enumerate}

\subsection*{§B --- Five-mark questions}

\begin{enumerate}[leftmargin=1.9cm]

\A{U3.26} \textbf{Verified.}
\begin{answerbox}
(a) $f'(w) = 2(w-4)$, so $w \leftarrow w - 0.15\cdot 2(w-4) = 0.7w + 1.2$.

(b) \begin{center}\begin{tabular}{@{}rrrr@{}}\toprule
$k$ & $w_k$ & $f(w_k)$ & $f'(w_k)$\\
\midrule
0 & 0.0000 & 16.0000 & $-8.0000$\\
1 & 1.2000 & 7.8400 & $-5.6000$\\
2 & 2.0400 & 3.8416 & $-3.9200$\\
3 & 2.6280 & 1.8824 & $-2.7440$\\
\bottomrule\end{tabular}\end{center}
$w_4 = \mathbf{3.0396}$.

(c) Gradients fall in the fixed ratio $5.6/8 = 3.92/5.6 = \mathbf{0.7}$, and
$|1 - 2\eta| = |1 - 0.30| = 0.7$. Convergence is \emph{geometric}: every step
buys the same percentage of the remaining distance.

(d) $w_k = 4 - 4(0.7)^k$; $w_4 = 4 - 4(0.2401) = \mathbf{3.0396}$. Agrees.
\end{answerbox}

\A{U3.27} \textbf{Verified.} (a) $f'' = 6$, so $0 < \eta < 2/6 =
\mathbf{0.3333}$. (b) $\eta = 1/6 = \mathbf{0.1667}$. (c) Multiplier
$|1 - 0.4\times6| = 1.4 > 1$: $w_1 = \mathbf{-1.4}$, $w_2 = \mathbf{1.96}$,
$w_3 = \mathbf{-2.744}$, $w_4 = \mathbf{3.8416}$ --- oscillating with growing
magnitude, i.e.\ diverging. (d) The loss curve rises, then produces
\texttt{inf} or \texttt{nan}. First action: divide $\eta$ by 10. A diverging
loss is almost always a learning-rate problem before it is anything else.

\A{U3.28} \textbf{Verified.}
\begin{answerbox}
(a) $\nabla J(w) = -\frac{2}{3}\sum x_i(y_i - wx_i)$. At $w_0 = 0$:
$\sum x_iy_i = 3 + 10 + 36 = 49$, so $g = -\frac23(49) = \mathbf{-32.6667}$ and
$w = 0 + 0.05(32.6667) = \mathbf{1.6333}$.

(b) SGD, per-example gradient $2x_i(wx_i - y_i)$:\\
$x=1$: $g = 2(0-3) = -6.0$, $w = \mathbf{0.3}$.\\
$x=2$: $g = 4(0.6-5) = -17.6$, $w = \mathbf{1.18}$.\\
$x=4$: $g = 8(4.72-9) = -34.24$, $w = \mathbf{2.892}$.

(c) $w^\star = 49/(1+4+16) = 49/21 = \mathbf{2.3333}$.

(d) One batch step reached 1.63; one epoch of SGD --- the same single pass over
the data --- reached 2.89, overshooting 2.33 slightly. The batch step is the
better \emph{single} step but SGD took three of them for the same data cost.
This is why $n$ mediocre steps beat one good one.
\end{answerbox}

\A{U3.29} \textbf{Verified.}
\begin{answerbox}
\begin{center}\begin{tabular}{@{}rrrr@{}}\toprule
$t$ & $g_t$ & $v_t$ & $w_t$\\
\midrule
1 & $-8.0000$ & $-8.0000$ & 1.2000\\
2 & $-5.6000$ & $-12.8000$ & 3.1200\\
3 & $-1.7600$ & $-13.2800$ & \textbf{5.1120}\\
\bottomrule\end{tabular}\end{center}
Plain GD reached $w_3 = 2.6280$; the optimum is $4$. Momentum has
\textbf{overshot to 5.11} --- past the minimum by more than plain GD was short
of it. The velocity still carries most of the huge first gradient, so step 3 is
almost as large as step 1 even though the true gradient has collapsed to
$-1.76$. It will now oscillate inward.
\end{answerbox}
\emph{The alternative convention $v \leftarrow \beta v - \eta g$, $w \leftarrow w + v$ reaches the same $w_t$; either is fine, so long as you stay consistent within a question.}

\A{U3.30} \textbf{Verified.} (a) $\eta_{\text{eff}} = 0.01/0.1 = 0.1$ at
$\beta = 0.9$; at $\beta = 0.99$ the same $\eta$ would give $0.01/0.01 = 1.0$,
ten times larger. (b) $\eta_{\text{new}} = \eta_{\text{old}}(1-\beta_{\text{new}})
/(1-\beta_{\text{old}}) = 0.01 \times 0.01/0.1 = \mathbf{0.001}$.
(c) Without rescaling, the effective step jumps tenfold and will very likely
cross the stability ceiling $2(1+\beta)/L$ --- the loss rises or goes
\texttt{nan}, and the cause looks like a momentum problem when it is a
learning-rate problem.

\A{U3.31} \textbf{Verified.}
\begin{answerbox}
\begin{center}\begin{tabular}{@{}rrrrrr@{}}\toprule
$t$ & $g_t$ & $s_t$ & $\sqrt{s_t}$ & step & $w_t$\\
\midrule
1 & $-8.0000$ & 64.0000 & 8.0000 & $-0.5000$ & 0.5000\\
2 & $-7.0000$ & 113.0000 & 10.6301 & $-0.3293$ & 0.8293\\
3 & $-6.3415$ & 153.2146 & 12.3780 & $-0.2562$ & \textbf{1.0854}\\
\bottomrule\end{tabular}\end{center}
Step shrank from 0.5000 to 0.2562: a fall of \textbf{48.8\%}. The gradient
shrank from 8.0000 to 6.3415: a fall of only \textbf{20.7\%}. The extra
shrinkage is the denominator: $s_t$ accumulates and never decays, so
$\sqrt{s_t}$ rose from 8.0 to 12.4 and the effective learning rate fell with
it. Note the first step is exactly $\eta$, because $s_1 = g_1^2$ makes the
ratio $g_1/|g_1| = \pm 1$.
\end{answerbox}

\A{U3.32} \textbf{Verified.}
\begin{answerbox}
(a) \begin{center}\begin{tabular}{@{}rrrrrr@{}}\toprule
$t$ & $g_t$ & $s_t$ & $\sqrt{s_t}$ & step & $w_t$\\
\midrule
1 & $-8.0000$ & 6.4000 & 2.5298 & $-1.5811$ & 1.5811\\
2 & $-4.8377$ & 8.1004 & 2.8461 & $-0.8499$ & 2.4310\\
3 & $-3.1380$ & \textbf{8.2750} & 2.8766 & $-0.5454$ & \textbf{2.9764}\\
\bottomrule\end{tabular}\end{center}

(b) $\eta/\sqrt{1-\gamma} = 0.5/\sqrt{0.1} = 0.5/0.3162 = \mathbf{1.5811}$ ---
matches the first step exactly, since $s_1 = (1-\gamma)g_1^2$ gives
$\sqrt{s_1} = |g_1|\sqrt{1-\gamma}$.

(c) $s_t$ has not yet fallen by $t=3$ (6.40 $\to$ 8.10 $\to$ 8.28, still
rising but decelerating); it falls as soon as $(1-\gamma)g_t^2 < (1-\gamma)s_{t-1}$,
i.e.\ once $g_t^2 < s_{t-1}$ --- here at $t = 4$. Adagrad's $s_t$ can
\emph{never} fall because it is a sum of non-negative terms with no decay
factor. That single difference is the whole of RMSprop.
\end{answerbox}

\A{U3.33} \textbf{Verified.}
\begin{answerbox}
\begin{center}\footnotesize\begin{tabular}{@{}rrrrrrrrr@{}}\toprule
$t$ & $g_t$ & $m_t$ & $v_t$ & $\hat m_t$ & $\hat v_t$ & $\sqrt{\hat v_t}$ & step & $w_t$\\
\midrule
1 & $-8.0000$ & $-0.8000$ & 0.064000 & $-8.0000$ & 64.0000 & 8.0000 & $-0.2000$ & 0.2000\\
2 & $-7.6000$ & $-1.4800$ & 0.121696 & $-7.7895$ & 60.8784 & 7.8025 & $-0.1997$ & 0.3997\\
3 & $-7.2007$ & $-2.0521$ & 0.173424 & $-7.5722$ & 57.8658 & 7.6070 & $-0.1991$ & \textbf{0.5988}\\
\bottomrule\end{tabular}\end{center}
Steps: $0.2000$, $0.1997$, $0.1991$ --- essentially constant at $\eta = 0.2$
regardless of the gradient's size. Adam sets the \emph{relative} step between
coordinates; the overall scale remains $\eta$'s job. Note $\hat m_1 = g_1$ and
$\hat v_1 = g_1^2$ exactly.
\end{answerbox}

\A{U3.34} \textbf{Verified.}
\begin{answerbox}
(a) Uncorrected steps: $\mathbf{-0.6325}$, $\mathbf{-0.8430}$,
$\mathbf{-0.9586}$, giving $w = 0.6325,\ 1.4755,\ 2.4341$.

(b) $(1-\beta_1)/\sqrt{1-\beta_2} = 0.1/\sqrt{0.001} = 0.1/0.031623 =
\mathbf{3.1623}$. The corrected first step was $0.2000$; the uncorrected was
$0.6325 = 0.2 \times 3.1623$. Exactly the predicted factor.

(c) With $m_0 = v_0 = 0$ and a constant $g$, $m_t = (1-\beta_1^t)g$ and
$v_t = (1-\beta_2^t)g^2$. The step uses $m/\sqrt v$, so the two biases do not
cancel: they enter as $(1-\beta_1^t)$ over $\sqrt{1-\beta_2^t}$. If both
moments shrank by the \emph{same} factor the ratio would be untouched and no
correction would be needed. It is the mismatch --- $\beta_1 = 0.9$ against
$\beta_2 = 0.999$ --- that makes the early uncorrected steps 3.16$\times$ too
large.
\end{answerbox}

\A{U3.35} \textbf{Verified.} (a) Require $1/(1-\beta^t) \le 1.01$, i.e.\
$\beta^t \le 1/101$, i.e.\ $t \ge \log(1/101)/\log\beta$. (b)
$\beta = 0.9 \to t \ge 43.8$, so $\mathbf{t = 44}$;
$\beta = 0.99 \to \mathbf{t = 460}$;
$\beta = 0.999 \to t \ge 4612.8$, so $\mathbf{t = 4613}$.
(c) A 2000-step run \emph{never} escapes the second moment's bias: correction
is active for the entire training run, not just a warm-up. Switching it off
would corrupt every step.

\A{U3.36} \textbf{Verified.} $s_1 = 0.1 \times 64 = 6.4$; $\sqrt{s_1 + \varepsilon}
= 2.5298$. Then $\Delta w_1 = \sqrt\varepsilon \times 8 / 2.5298$:
\begin{center}\begin{tabular}{@{}rrr@{}}\toprule
$\varepsilon$ & $\sqrt\varepsilon$ & $\Delta w_1$\\
\midrule
$10^{-6}$ & $10^{-3}$ & $\mathbf{3.1623\times10^{-3}}$\\
$10^{-8}$ & $10^{-4}$ & $\mathbf{3.1623\times10^{-4}}$\\
$10^{-12}$ & $10^{-6}$ & $\mathbf{3.16\times10^{-6}}$\\
\bottomrule\end{tabular}\end{center}
$\Delta w_1 \propto \sqrt\varepsilon$ exactly. So Adadelta does have a
hyperparameter controlling its step scale --- it is simply spelled
$\varepsilon$ and presented as a numerical-stability constant. The claim ``no
learning rate'' is true of the \emph{form} of the update and misleading about
the practice.

\A{U3.37} \textbf{Verified.}
\begin{center}\begin{tabular}{@{}rrr@{}}\toprule
$B$ & $1/\sqrt B$ & work per update\\
\midrule
1 & 1.0000 & $1\times$\\ 4 & 0.5000 & $4\times$\\ 16 & 0.2500 & $16\times$\\
64 & 0.1250 & $64\times$\\ 256 & 0.0625 & $256\times$\\
\bottomrule\end{tabular}\end{center}
Doubling $B$ reduces noise by a factor $\sqrt2 = 1.414$ --- that is, by
\textbf{29\%} --- for \textbf{100\%} more work. Conclusion: there is a broad
flat optimum. Start at $B = 32$, use powers of two, and let memory rather than
statistics decide the upper limit; tune $\eta$ before $B$.
\emph{2.5 table, 1 for the $\sqrt2$/29\% figure, 1.5 for the conclusion.}

\A{U3.38} Batch: $B = n$, 1 update/epoch, no gradient noise, $O(np)$ per
update, vectorises well. Mini-batch: $B = 32$--$512$, $n/B$ updates, moderate
noise, $O(Bp)$, vectorises well. Stochastic: $B = 1$, $n$ updates, large noise,
$O(p)$, does \emph{not} vectorise. Comments should note that mini-batch is
strictly better than both extremes, which is why it is the universal default.

\A{U3.39} SGD: $w \leftarrow w - \eta g$. Momentum: adds a velocity ---
$v \leftarrow \beta v + g$, $w \leftarrow w - \eta v$ --- adapting the
\emph{direction}. Adagrad: adds per-coordinate scaling by
$\sqrt{\sum g^2}$ --- adapting the \emph{size}. RMSprop: replaces Adagrad's sum
with an EWMA so the denominator can shrink. Adam: puts momentum's first moment
on top of RMSprop's second, plus bias correction.

\A{U3.40} Adagrad: $s_i = \sum_t g_{i,t}^2$ (a sum). RMSprop:
$s_i = \gamma s_i + (1-\gamma)g_i^2$ (an EWMA). Adam: $s_i = \hat v_i$, an
EWMA with bias correction, and additionally the numerator $g_i$ is replaced by
$\hat m_i$. So Adam is the only one that also changes the numerator --- worth
saying, and worth a mark.

\A{U3.41} Plain SGD's step $-\eta g$ has units $[\eta][J]/[w]$ --- it depends
on how the loss is scaled. Adagrad's and RMSprop's $-\eta g/\sqrt s$ has units
$[\eta]$, because $\sqrt s$ cancels $g$ exactly. Newton's $-H^{-1}g$ has units
$[w]$ --- dimensionally correct, and therefore the only one whose step is
genuinely a distance in parameter space. It is not used because forming and
inverting $H$ costs $O(p^3)$.

\A{U3.42} A dense feature is updated every batch and its accumulated squared
gradient grows fast, so it wants a small step; a feature firing in 18 rows of
4000 is updated rarely and needs a large step when it is. One scalar $\eta$
must serve both and will be wrong for one of them. Reach for \textbf{Adagrad}:
because $s_i$ accumulates \emph{per coordinate}, the rare feature's denominator
stays small and its effective step stays large --- measured at 31.8$\times$ the
median dense step.

\A{U3.43} Standardising changes the Hessian's eigenvalues and so its condition
number --- measured at $470 \to 1.17\times10^8$ when one feature was multiplied
by 1000, with the stability limit $2/L$ collapsing from 0.2485 to
$4.42\times10^{-4}$. Since $\kappa$ sets the convergence rate and $2/L$ sets
the largest usable $\eta$, scaling is a decision about the optimisation
problem, not about tidiness.

\A{U3.44} (1) Standardise. (2) $B = 32$. (3) Find $\eta$ by bisection on a log
scale --- try $10^{-1}, 10^{-2}, 10^{-3}$, take the largest that does not
diverge, back off by $3\times$. (4) Add momentum $\beta = 0.9$, dividing $\eta$
by 10 if it was tuned without. (5) Watch the loss curve and read its signature.
(6) Decay $\eta$ at the end.
\end{enumerate}

\subsection*{§C --- Ten-mark questions}

\begin{enumerate}[leftmargin=1.9cm]
\A{U3.45} Coverage: the update rule; the four regimes with the multiplier
$|1-\eta f''|$ evaluated in each; the derivation of $2/L$ (the iterate
contracts iff the multiplier has magnitude below 1); the four loss-curve
signatures; bisection on a log scale. Worked example: U3.26 or U3.27 numbers.
\emph{The derivation of $2/L$ is the one part students omit; require it.}

\A{U3.46} Expected chain, each with problem/fix/rule/cost:
batch (exact but $O(np)$ per update) $\to$ SGD (cheap steps, but a noise floor)
$\to$ mini-batch (noise as $1/\sqrt B$, vectorises; introduces $B$) $\to$
momentum (ravines; introduces $\beta$ and overshoot) $\to$ Adagrad
(per-coordinate scale; stalls) $\to$ RMSprop (EWMA fixes the stall; loses the
convergence proof) $\to$ Adam (adds the first moment and bias correction; three
more hyperparameters).

\A{U3.47} Coverage: the ravine and $\kappa$; $v \leftarrow \beta v + g$;
the unrolled EWA $v_t = \sum_j \beta^j g_{t-j}$ with total weight $1/(1-\beta)$;
$\eta_{\text{eff}} = \eta/(1-\beta)$; the ceiling $\eta < 2(1+\beta)/L$;
overshoot with numbers (U3.29 gives $w_3 = 5.11$ against an optimum of 4);
rate $1-2/\kappa \to 1-2/\sqrt\kappa$; and the caution that momentum does not
help when $\kappa \approx 1$.

\A{U3.48} Coverage: the two moments and their roles (direction / scale); the
constant-gradient derivation $m_t = (1-\beta_1^t)g$; the first-step result
$\pm\eta$; the $3.1623\times$ error if correction is omitted, growing to
$6.57\times$ at $t = 12$; the defaults $\beta_1 = 0.9$, $\beta_2 = 0.999$,
$\varepsilon = 10^{-8}$, $\eta = 10^{-3}$; and the misconception that
``adaptive'' means schedule-free.

\A{U3.49} A four-way table plus at least two hand-computed steps of two of
them. Expect: Adagrad accumulates a sum, never decays, first step $\eta$, one
hyperparameter, fixes per-coordinate scale, right for sparse data. RMSprop
accumulates an EWMA, decays, first step $\eta/\sqrt{1-\gamma}$, adds $\gamma$,
fixes the stall, right for non-stationary objectives. Adadelta accumulates two
EWMAs, first step $\propto\sqrt\varepsilon$, fixes the units, right when you
cannot tune $\eta$. Adam accumulates both moments with correction, first step
$\eta$, adds $\beta_1,\beta_2$, is the general default.

\A{U3.50} Coverage: the claim that final performance differs by about one per
cent on a tuned convex problem; the sensitivity experiment --- sweep $\eta$
over four decades, count the grid points reaching within 5\% of the optimum,
report the width in decades (Adagrad 2.24 against SGD 1.31); and why a wider
window matters more in practice than a marginally better optimum, because tuning
budget is the scarce resource.
\end{enumerate}

\subsection*{§D --- Statements}

\begin{enumerate}[leftmargin=1.9cm]
\A{U3.51} \textbf{False.} On the diabetes benchmark SGD with momentum finished
\emph{best} (0.4846) and Adam \emph{last} (0.4929) over 30 epochs. With a tuned
$\eta$ the choice of optimiser is worth about one per cent; Adam's real
advantage is robustness to $\eta$, not final loss.
\A{U3.52} \textbf{False.} Adam adapts the \emph{relative} step between
coordinates; the overall scale is still $\eta$. Measured: with decay off, the
loss got \emph{worse} between epoch 5 (0.4888) and epoch 30 (0.4929); with
decay on it improved to 0.4850.
\end{enumerate}

\vspace{0.4em}\hrule

% =====================================================================
\section*{Unit IV --- Data preprocessing}

\subsection*{§A --- Two-mark questions}

\begin{enumerate}[leftmargin=1.9cm]
\A{U4.1} Data cleaning, data integration, data reduction, data transformation.
\A{U4.2} $P = (1-p)^d$. At $p=0.05$, $d=30$: $0.95^{30} = \mathbf{0.2146}$ ---
79\% of rows lost.
\A{U4.3} Numeric: mean (symmetric, few holes), median (skewed or
outlier-prone), a constant with an indicator (missingness may be informative),
KNN or MICE (holes are predictable from other columns). Categorical:
most-frequent, or an explicit \texttt{"Missing"} level.
\A{U4.4} $\sigma_{\text{after}}/\sigma_{\text{before}} = \sqrt{m/n} =
\sqrt{1-p}$, where $m$ of $n$ are observed.
\A{U4.5} MCAR: $P(R)$ depends on nothing (a dropped test tube). MAR: depends on
observed columns only (BMI recorded only for high-BP patients). MNAR: depends
on the missing value itself (the heaviest patients decline to be weighed).
\A{U4.6} MAR --- a multivariate imputer recovered 88\% of the bias
($+1.3473 \to -0.1588$). MNAR --- nothing in the data can help, because the
information needed is precisely what is absent.
\A{U4.7} A binary column recording \emph{that} a value was missing. It is
leakage when the reason for the hole will not exist at deployment --- ask when
in the real timeline the missingness is decided.
\A{U4.8} $|z_i| = |x_i - \bar x|/s > 3$; $x_i < Q_1 - 1.5\,\text{IQR}$ or
$x_i > Q_3 + 1.5\,\text{IQR}$; $|M_i| = |0.6745(x_i - \tilde x)/\text{MAD}| >
3.5$.
\A{U4.9} $z$-score \textbf{0\%} (one point moves both $\bar x$ and $s$);
Tukey/IQR \textbf{25\%}; modified $z$ \textbf{50\%} (both statistics are
medians).
\A{U4.10} An extreme value inflates the standard deviation it is measured
against, so its own $z$-score shrinks and it hides behind its own effect. With
$k$ contaminants present, each masks the others.
\A{U4.11} $|z_1| \le (n-1)/\sqrt n$. It exceeds 3 first at
$\mathbf{n = 11}$ ($3.0151$); below that the $|z|>3$ test is not merely weak,
it is \emph{inert} --- \texttt{False} for every possible dataset.
\A{U4.12} $0.6745$ is the 0.75 quantile of the standard Normal, so
$\text{MAD} \approx 0.6745\sigma$ for Normal data and the modified $z$ is on
the same scale as an ordinary $z$. The multiplier $1.5$ makes
$Q_3 + 1.5\,\text{IQR} \approx 2.70\sigma$, flagging about 0.70\% of Normal
data.
\A{U4.13} Skew is a property of the true distribution; contamination is a
property of the recording. An outlier rule applied to a skewed variable is
mostly a skewness detector --- on \texttt{'area error'} the IQR fence flagged
11.4\% of rows, falling to 1 row after Yeo--Johnson. \textbf{Transform first,
then detect.}
\A{U4.14} Leave alone (it is real and it matters); trim/drop (it is provably an
error); winsorise or cap at a fence (keep the row, bound the influence); use a
robust method instead (\texttt{RobustScaler}, MAE, Huber).
\A{U4.15} Equal-width cuts the \emph{range} into intervals of equal size;
equal-depth puts an equal \emph{count} in each bin. Smoothing: by bin means, by
bin medians, by bin boundaries.
\A{U4.16} $x^{(\lambda)} = (x^\lambda - 1)/\lambda$ for $\lambda \ne 0$,
$\log x$ for $\lambda = 0$. It requires $x > 0$; \textbf{Yeo--Johnson} removes
that restriction.
\A{U4.17} To make the family continuous at $\lambda = 0$: as $\lambda \to 0$,
$(x^\lambda - 1)/\lambda \to \log x$ by l'H\^opital. Without the $-1$ and the
division, the limit would be the constant 1.
\A{U4.18} By \emph{maximum likelihood} --- \texttt{PowerTransformer} and
\texttt{scipy.stats.boxcox} choose the $\lambda$ maximising the Normal
likelihood. It is \textbf{not} chosen to zero the skewness: on the same column
those two criteria gave $-0.4357$ and $-0.4688$.
\A{U4.19} \texttt{StandardScaler} $(x-\bar x)/\sigma$, mean 0 sd 1, unbounded.
\texttt{MinMaxScaler} $(x-x_{\min})/(x_{\max}-x_{\min})$, exactly $[0,1]$.
\texttt{RobustScaler} $(x-\tilde x)/\text{IQR}$, median 0 IQR 1, unbounded.
\texttt{MaxAbsScaler} $x/\max|x|$, $[-1,1]$, preserves zeros.
\A{U4.20} \texttt{MinMaxScaler} --- it is \emph{defined} by the two extremes,
so one bad value sets the entire range. \texttt{RobustScaler} is safest: median
and IQR are order statistics with a 25\% breakdown point.
\A{U4.21} A split asks ``is $x_j \le t$?''. Under $x \mapsto (x-a)/b$ with
$b > 0$ the threshold moves to $(t-a)/b$ and exactly the same rows go left.
Predictions, structure and accuracy are identical; only the printed thresholds
become unreadable.
\A{U4.22} Scaling changes location and width only --- skewness and kurtosis are
unchanged. A transformation (log, Box--Cox, Yeo--Johnson) changes the
\emph{shape}, and so changes skewness.
\A{U4.23} Nominal categories have no order (colour, city); ordinal do
(small/medium/large). Nominal $\to$ one-hot. Ordinal $\to$ ordinal encoding
\emph{with the order supplied explicitly}.
\A{U4.24} With an intercept and all $k$ dummies, the dummies sum to 1 in every
row and duplicate the intercept, so the design matrix is rank deficient. It
breaks coefficients, standard errors, $t$-statistics and $p$-values. It does
\emph{not} break predictions --- both coefficient vectors gave identical
predictions and $R^2 = 0.986231$.
\A{U4.25} Alphabetical. So \texttt{large} sorts before \texttt{medium} before
\texttt{small}, producing the codes 0, 1, 2 in exactly the wrong order --- and
it does so silently, with no error and no warning.
\A{U4.26} It emits an all-zero row for a category not seen during
\texttt{fit}. Without it, \texttt{OneHotEncoder} raises
\texttt{ValueError: Found unknown categories} at \texttt{transform} time --- in
production, or inside a CV fold.
\A{U4.27} Dimensionality reduction (fewer columns), numerosity reduction (fewer
rows), and compression.
\A{U4.28} As $d$ grows, the ratio of the farthest to the nearest distance tends
to 1 --- measured $9135$ at $d=1$ against $1.11$ at $d=1000$. Every
distance-based method degrades, because ``nearest'' stops meaning anything.
\A{U4.29} Centre the data; form the covariance
$S = \frac{1}{n-1}X_c^\top X_c$; take its eigenvalues and eigenvectors, sorted
by eigenvalue; project, $T = X_c U_k$.
\A{U4.30} $\lambda_j$ \emph{is} the variance along component $j$. The trace
equals $\sum_j \lambda_j$ --- the total variance, which a rotation moves
around but never creates or destroys.
\A{U4.31} PCA maximises variance, and variance has units. Unstandardised on
breast\_cancer, PC1 explained 98.20\% and was essentially the single column
\texttt{worst area} (sd 568.9); standardised, PC1 explained 44.27\% and drew on
many columns. Without standardising you are ranking columns by their units.
\A{U4.32} PCA keeps \emph{combinations} of all the features, not a subset of
them. On breast\_cancer, 26 of 30 loadings on PC1 exceeded 0.10. Saying you
``only need ten components'' does not mean you only need ten measurements ---
you still need all thirty to compute them.
\A{U4.33} Variance is not relevance. In the designed example a $\mathcal
N(0,10^2)$ noise column carried 99.05\% of the variance while the label
depended entirely on a $\mathcal N(0,1)$ column: logistic regression scored
0.9867 on both features and \textbf{0.5217} on PCA's one component. Supervised
alternative: \textbf{LDA} (or PLS).
\A{U4.34} \emph{Filter} (variance threshold, ANOVA $F$, mutual information) ---
cheapest, but blind to interactions and redundancy. \emph{Wrapper} (RFE, forward
selection) --- most expensive, and can overfit the selection.
\emph{Embedded} (L1/lasso) --- nearly free, but the answer is specific to that
model.
\A{U4.35} $F$ measures \emph{linear} association; mutual information measures
\emph{any} statistical dependence. A feature entering only through $|b|$ ranked
1st by MI and 8th of 10 by $F$ ($F = 0.122$, below several pure-noise columns).
\A{U4.36} L1 has a corner at zero, so the optimum frequently lands exactly on
it and coefficients become \emph{exactly} zero --- selection. L2 is smooth, so
coefficients shrink towards zero without reaching it --- ridge at
$\alpha = 1$ kept all 10 non-zero while lasso at $\alpha = 5$ kept 5.
\A{U4.37} Any step that \emph{learns} anything from the data --- a mean, a
median, a category vocabulary, a feature subset, a hyperparameter, a resampled
row --- must be fitted on the training portion only, and therefore belongs
inside the cross-validation loop.
\A{U4.38} \textbf{Feature selection}, by a wide margin, because it looks at
$y$. Measured: scaling on everything leaked $+0.0034$ at $n_{\text{train}} = 20$
and $+0.0002$ at 427; selection on everything leaked $+0.3133$. Scaling and PCA
never see the labels; selection does.
\A{U4.39} So that the class proportions of the whole are preserved in each
part. Without it a fold can contain no minority rows at all, making recall
undefined. Argument: \texttt{stratify=y}, or use \texttt{StratifiedKFold}.
\A{U4.40} \texttt{shuffle=False}. On a file sorted by label, the folds
partition by class: measured, 200 rows with 10 positives gave positives per
test fold $[10,0,0,0,0]$, mean recall $0.0000$ and mean accuracy $1.0000$ with
no error message.
\A{U4.41} The unit that could not have produced another row --- the level at
which observations are genuinely independent. Not the row when: many windows
per patient; many transactions per cardholder; many sentences per document.
\A{U4.42} \texttt{GroupKFold}; \texttt{TimeSeriesSplit}; \texttt{StratifiedKFold};
\texttt{StratifiedGroupKFold}.
\A{U4.43} $k$ fits; each on $n(1 - 1/k)$ rows; total row-fits $n(k-1)$.
\A{U4.44} Cost: $n$ fits --- on breast\_cancer, 2.07\,s against 0.02\,s for
5-fold, a factor of 98 for the same accuracy (0.9789 both). Variance: the
across-fold sd is $\sqrt{p(1-p)} = 0.1437$, which is a property of scoring one
row at a time, not of the estimator, and tells you nothing.
\A{U4.45} It is the spread of the $k$ fold scores, not the standard error of
their mean. LOOCV's 0.1437 comes from each fold being a single 0/1 outcome.
Report the mean with the spread across repeated CV runs if you need an
uncertainty.
\A{U4.46} The validation set \emph{selects} --- look at it as often as you
like. The test set \emph{estimates} --- look at it once, after everything is
fixed.
\A{U4.47} It is the maximum over the grid of noisy scores, so it is a selection
and is biased upward --- the expected maximum of $m$ noisy scores grows like
$\mu + \sigma\sqrt{2\ln m}$. Measured on pure noise: 1 combination gave
optimism exactly 0.0000, 60 combinations gave $+0.0783$. The estimate is
\textbf{nested cross-validation}, or a held-out test set touched once.
\A{U4.48} $\text{acc} = \frac{TP+TN}{TP+TN+FP+FN}$;
$\text{prec} = \frac{TP}{TP+FP}$; $\text{rec} = \frac{TP}{TP+FN}$;
$F_1 = \frac{2TP}{2TP+FP+FN}$;
$\text{bal} = \frac12\left(\frac{TN}{TN+FP} + \frac{TP}{TP+FN}\right)$.
\A{U4.49} ROC--AUC: \textbf{0.5} always. PR--AUC: the \textbf{positive rate}
--- 0.01 on a 1\%-positive problem. This is why a PR--AUC of 0.19 there can be
a real result.
\A{U4.50} $x_{\text{new}} = A + u(B-A)$, $u \sim \mathcal U(0,1)$, where $A$ is
a minority point and $B$ is one of $A$'s $k$ nearest \emph{minority}
neighbours.
\A{U4.51} Any two: categorical features (the interpolation is meaningless);
multi-modal minority classes (the line between two modes passes through empty
space); minority points sitting inside a majority region (synthetic points land
squarely in majority territory).
\A{U4.52} Inside each fold, on the \emph{training portion only}, after the
transformers are fitted and immediately before the model fit. Earlier, and
copies of the same row land in both training and validation --- measured
$F_1$ 0.8942 against 0.1133, an overstatement of $+0.7809$, with 100\% of the
minority test rows being exact copies of training rows.
\A{U4.53} \textbf{Moving the decision threshold} using the precision--recall
curve. It is cheaper (no refitting, no extra rows), deterministic, and leaves
the model's probabilities calibrated.
\end{enumerate}

\subsection*{§B --- Five-mark questions}

\begin{enumerate}[leftmargin=1.9cm]

\A{U4.54} \textbf{Verified.}
\begin{answerbox}
(a) Observed $48,55,60,67,72,78$: mean $= 380/6 = \mathbf{63.3333}$;
median $= (60+67)/2 = \mathbf{63.5}$; sd $\text{ddof}{=}1 = \mathbf{11.1295}$,
$\text{ddof}{=}0 = \mathbf{10.1598}$.

(b) Mean-imputed (two extra values at 63.3333): mean $= \mathbf{63.3333}$
(unchanged, by construction), population sd $= \mathbf{8.7987}$.

(c) Median-imputed (two extra at 63.5): mean $= \mathbf{63.375}$,
population sd $= \mathbf{8.7990}$.

(d) $m = 6$, $n = 8$, so $\sqrt{m/n} = \sqrt{0.75} = \mathbf{0.8660}$.
Ratio of imputed to observed population sd $= 8.7987/10.1598 = \mathbf{0.8660}$.
\emph{Why:} the imputed values sit exactly at the mean, so they contribute
\emph{zero} to $\sum(x_i - \bar x)^2$ while still counting in the denominator
$n$. The sum of squares is unchanged and the divisor grew from $m$ to $n$,
giving a variance factor of exactly $m/n$ and an sd factor of $\sqrt{m/n}$.
\end{answerbox}

\A{U4.55} \textbf{Verified.}
\begin{answerbox}
(a) \begin{center}\begin{tabular}{@{}rrrrr@{}}\toprule
$p$ & $d=5$ & $d=20$ & $d=30$ & $d=50$\\
\midrule
0.01 & 0.9510 & 0.8179 & 0.7397 & 0.6050\\
0.05 & 0.7738 & 0.3585 & \textbf{0.2146} & 0.0769\\
0.10 & 0.5905 & 0.1216 & 0.0424 & \textbf{0.0052}\\
\bottomrule\end{tabular}\end{center}

(b) $(0.95)^d = 0.5 \Rightarrow d = \log 0.5/\log 0.95 = \mathbf{13.5}$
columns.

(c) $0.9^{50} = 0.005154$, so $1000/0.005154 = \mathbf{194{,}033}$ raw rows.

(d) Listwise deletion is right when the rows are few relative to $n$, the
missingness is MCAR (so the survivors are unbiased), and you want an analysis
whose validity does not depend on an imputation model. It is never free: with
$d = 30$ and $p = 5\%$ you are discarding 79\% of your evidence.
\end{answerbox}

\A{U4.56} \textbf{Verified.}
\begin{answerbox}
(a) $n = 12$. Mean $= 452/12 = \mathbf{37.6667}$;
sd ($\text{ddof}{=}1$) $= \mathbf{39.6584}$; median $= \mathbf{21.0}$;
absolute deviations from the median are $1,0,2,1,1,0,2,2,1,1,99,104$, so
$\text{MAD} = \mathbf{1.0}$.

(b) $Q_1 = \mathbf{20.0}$, $Q_3 = \mathbf{22.25}$, $\text{IQR} = \mathbf{2.25}$;
fences $20 - 3.375 = \mathbf{16.625}$ and $22.25 + 3.375 = \mathbf{25.625}$.

(c) For $x = 120$: $z = (120 - 37.6667)/39.6584 = \mathbf{2.0761}$ --- under
$|z|>3$, \textbf{not flagged}. Tukey: $120 > 25.625$ --- \textbf{flagged}.
Modified $z = 0.6745(120-21)/1.0 = \mathbf{66.78}$ --- \textbf{flagged}, by a
factor of 19 over the threshold.
Flag counts across the column: $|z|>3$ \textbf{0}; IQR fence \textbf{2};
$|M|>3.5$ \textbf{2}.

(d) Standard deviation of the ten genuine readings alone: $\mathbf{1.3375}$.
The value 120 would then score $z = (120-20.7)/1.3375 = \mathbf{74.24}$.
The discrepancy is \textbf{masking}: the two contaminants inflated $s$ from
1.34 to 39.66 --- a factor of 30 --- so each hides behind the variance the
\emph{pair} of them created. Note also the ceiling: with $n = 12$ the largest
attainable $|z|$ is $11/\sqrt{12} = 3.1754$, so the rule had almost no room to
fire even in principle.
\end{answerbox}

\A{U4.57} \textbf{Verified.}
\begin{answerbox}
\begin{center}\begin{tabular}{@{}rrl@{}}\toprule
$n$ & $(n-1)/\sqrt n$ & can $|z|>3$ fire?\\
\midrule
5 & 1.7889 & \textbf{no}\\ 8 & 2.4749 & \textbf{no}\\ 10 & 2.8460 & \textbf{no}\\
11 & 3.0151 & yes\\ 20 & 4.2485 & yes\\
\bottomrule\end{tabular}\end{center}
\emph{Derivation.} Let $d_i = x_i - \bar x$. Then $\sum d_i = 0$ and
$\sum d_i^2 = (n-1)s^2$. From $\sum_{i\ge2} d_i = -d_1$ and Cauchy--Schwarz,
$\sum_{i\ge2} d_i^2 \ge d_1^2/(n-1)$. Hence
$(n-1)s^2 \ge d_1^2 + d_1^2/(n-1) = d_1^2\frac{n}{n-1}$, giving
$|z_1| = |d_1|/s \le (n-1)/\sqrt n$.

\emph{Consequence.} On batches of eight, \texttt{if abs(z) > 3} evaluates to
\texttt{False} for \emph{every possible input}. The line is not a weak test; it
is dead code that looks like a test, and it will never raise an alert or an
error.
\end{answerbox}
\emph{1.5 table, 2 derivation, 1.5 consequence.}

\A{U4.58} \textbf{Verified.}
\begin{answerbox}
(a) \begin{center}\begin{tabular}{@{}llrrl@{}}\toprule
bin & values & mean & median & boundaries\\
\midrule
1 & $5, 8, 14$ & 9 & 8 & 5, 14\\
2 & $17, 20, 26$ & 21 & 20 & 17, 26\\
3 & $27, 30, 36$ & 31 & 30 & 27, 36\\
\bottomrule\end{tabular}\end{center}

(b) By means: $9,9,9,\ 21,21,21,\ 31,31,31$.
By medians: $8,8,8,\ 20,20,20,\ 30,30,30$.
By boundaries (each value to its nearer boundary):
$5,5,14,\ 17,17,26,\ 27,27,36$.

(c) Equal-width on $[5,36]$: three intervals of width
$31/3 = \mathbf{10.33}$ --- $[5, 15.33)$, $[15.33, 25.67)$, $[25.67, 36]$ ---
with bin sizes $\mathbf{3, 2, 4}$ rather than $3,3,3$.

(d) Gains: local noise is smoothed away, and a linear model can express a
non-monotone relationship it otherwise could not (measured $+0.0914$ $R^2$).
Costs: within-bin resolution is destroyed, the cut points are arbitrary, and a
tree --- which can already find its own thresholds --- is actively harmed
($-0.0534$).
\end{answerbox}

\A{U4.59} \textbf{Verified.}
\begin{answerbox}
\begin{center}\begin{tabular}{@{}lrrl@{}}\toprule
treatment & mean & sd & comment\\
\midrule
leave alone & 37.6667 & 39.6584 & useless\\
winsorise 10/90 & 35.6500 & 34.8891 & \textbf{did almost nothing}\\
cap at IQR fence & 21.5208 & 2.2669 & close\\
drop by modified $z$ & \textbf{20.7000} & \textbf{1.3375} & exact\\
\bottomrule\end{tabular}\end{center}
Truth: mean 20.7, sd 1.34.

\emph{Ranking:} drop-by-modified-$z$ (recovers the truth exactly, because the
rule identified precisely the two contaminants); cap at the fence (close, but
the two capped values sit at 25.625 and still pull the mean up by 0.8); winsorise
10/90 (a 10\% tail on $n=12$ is 1.2 values, so it clips one contaminant down
towards the other and leaves the damage almost intact); leave alone.

\emph{The lesson:} a fixed percentile is a \emph{guess} about how much dirt
there is. When the guess is wrong the treatment does nothing while appearing to
have been applied.
\end{answerbox}

\A{U4.60} \textbf{Verified.} (a) Gaps $4,\ 12,\ 36$ --- growing by a factor of
3. (b) $\log$: $0.6931,\ 1.7918,\ 2.8904,\ 3.9890$; gaps
$\mathbf{1.0986,\ 1.0986,\ 1.0986}$. (c) Exactly $\ln 3 = 1.0986$, because
$\log(3x) = \log x + \log 3$: a constant \emph{ratio} becomes a constant
\emph{difference}. (d) Most quantities with a long right tail are multiplicative
in origin --- incomes, populations, areas, prices. The log converts that
multiplicative structure into the additive structure a linear model assumes,
and it does so with a single rung of the ladder of powers and no fitted
parameter.

\A{U4.61} \textbf{Verified.}
\begin{answerbox}
\begin{center}\begin{tabular}{@{}lrrrr@{}}\toprule
$\lambda$ & $x=1$ & $4$ & $16$ & $64$\\
\midrule
$0.5$ & 0.0000 & 2.0000 & 6.0000 & 14.0000\\
$-0.5$ & 0.0000 & 1.0000 & 1.5000 & 1.7500\\
$0.001$ & 0.0000 & 1.3873 & 2.7764 & 4.1675\\
$\log x$ & 0.0000 & 1.3863 & 2.7726 & 4.1589\\
\bottomrule\end{tabular}\end{center}
The $\lambda = 0.001$ row agrees with $\log x$ to three decimals and would
agree to more as $\lambda \to 0$. This is why the $\lambda = 0$ case of
Box--Cox is \emph{defined} to be $\log x$: the $-1$ and the division by
$\lambda$ exist precisely so that the limit is the logarithm rather than the
constant 1.

Note also $\lambda = -0.5$ compresses the range hardest (1 to 1.75) --- the
further down the ladder, the more a long right tail is pulled in.
\end{answerbox}

\A{U4.62} \textbf{Verified.}
\begin{answerbox}
(a) Mean $= 48/8 = \mathbf{6.0}$; population sd $= \mathbf{2.0}$;
min $\mathbf{3}$, max $\mathbf{10}$; median $= \mathbf{5.5}$;
$Q_1 = \mathbf{5.0}$, $Q_3 = \mathbf{6.5}$, $\text{IQR} = \mathbf{1.5}$.

(b) \begin{center}\footnotesize\begin{tabular}{@{}lrrrrrrrr@{}}\toprule
& 3 & 5 & 5 & 5 & 6 & 6 & 8 & 10\\
\midrule
Standard & $-1.5$ & $-0.5$ & $-0.5$ & $-0.5$ & 0 & 0 & 1.0 & 2.0\\
MinMax & 0 & 0.2857 & 0.2857 & 0.2857 & 0.4286 & 0.4286 & 0.7143 & 1.0\\
Robust & $-1.6667$ & $-0.3333$ & $-0.3333$ & $-0.3333$ & 0.3333 & 0.3333 & 1.6667 & 3.0\\
MaxAbs & 0.3 & 0.5 & 0.5 & 0.5 & 0.6 & 0.6 & 0.8 & 1.0\\
\bottomrule\end{tabular}\end{center}

(c) Standard: mean 0, sd 1, \emph{unbounded}. MinMax: exactly $[0,1]$.
Robust: median 0, IQR 1, unbounded. MaxAbs: $[-1,1]$, and zeros stay zero.
\end{answerbox}

\A{U4.63} \textbf{Verified.}
\begin{answerbox}
(a) With 400 appended ($n = 9$): mean $6.0 \to \mathbf{49.7778}$;
population sd $2.0 \to \mathbf{123.8366}$; median $5.5 \to \mathbf{6.0}$;
IQR $1.5 \to \mathbf{3.0}$.

(b) Span occupied by the eight good values:
\begin{center}\begin{tabular}{@{}lrrr@{}}\toprule
scaler & before & after & compression\\
\midrule
Standard & 3.5000 & 0.0565 & $\mathbf{61.9\times}$\\
MinMax & 1.0000 & 0.0176 & $\mathbf{56.7\times}$\\
Robust & 4.6667 & 2.3333 & $\mathbf{2.0\times}$\\
\bottomrule\end{tabular}\end{center}

(c) \textbf{RobustScaler} survives. Mean and sd have 0\% breakdown --- one
point moves both without limit --- and min/max are the extremes themselves, so
one bad value \emph{defines} MinMax's range. Median and IQR are order
statistics: the typo moved the median by 0.5 and the IQR by 1.5, so the good
values keep most of their spread.

(d) Its output is \emph{guaranteed} to lie in $[0,1]$, so it passes every sanity
check you might run on it. Nothing in the output announces that eight of nine
rows are now crammed into a window of width 0.018 --- the numbers look
perfectly scaled, and the information has already been destroyed.
\end{answerbox}

\A{U4.64} \textbf{Verified.}
\begin{answerbox}
(a) $d(A,B)^2 = 25^2 + 5000^2 = 625 + 25{,}000{,}000 = \mathbf{25{,}000{,}625}$
--- income supplies \textbf{99.998\%}.
$d(A,C)^2 = 1^2 + 400{,}000^2 = 1 + 1.6\times10^{11} =
\mathbf{160{,}000{,}000{,}001}$ --- income supplies \textbf{100.000\%}.

(b) Column means $(38.6667,\ 635{,}000)$, population sds
$(11.5566,\ 187{,}394.4)$. Standardised:
$A = (-0.7499,\ -0.7204)$, $B = (1.4133,\ -0.6937)$, $C = (-0.6634,\ 1.4141)$.
$d(A,B)^2 = 4.6797 + 0.0007 = \mathbf{4.6804}$;
$d(A,C)^2 = 0.0075 + 4.5562 = \mathbf{4.5637}$.

(c) Raw: $A$'s nearest neighbour is $\mathbf{B}$ ($2.5\times10^7$ against
$1.6\times10^{11}$). Standardised: $\mathbf{C}$ ($4.5637 < 4.6804$).
\textbf{The neighbour changed.}

(d) Choosing not to scale is choosing to weight every feature by its own
standard deviation. Here that means declaring one rupee of income as important
as one year of age --- a claim nobody would defend if stated aloud, but which is
made silently by every unscaled distance computation.
\end{answerbox}

\A{U4.65} \textbf{Verified.}
\begin{answerbox}
(a) $\bar x = 1$; $\bar y = 115/3 = \mathbf{38.3333}$.
$S_{xy} = (-1)(30-38.3333) + 0 + (1)(15-38.3333) = 8.3333 - 23.3333 =
\mathbf{-15}$. $S_{xx} = 1 + 0 + 1 = \mathbf{2}$.

(b) Slope $= -15/2 = \mathbf{-7.5}$;
intercept $= 38.3333 + 7.5 = \mathbf{45.8333}$.

(c) Predictions: blue $\mathbf{45.8333}$, green $\mathbf{38.3333}$,
red $\mathbf{30.8333}$ --- against true means 30, 70, 15.
$\text{SSE} = 15.8333^2 + 31.6667^2 + 15.8333^2 = \mathbf{1504.1667}$;
$\text{SST} = 8.3333^2 + 31.6667^2 + 23.3333^2 = \mathbf{1616.6667}$;
$R^2 = 1 - 1504.1667/1616.6667 = \mathbf{0.0696}$.

(d) One-hot achieves $R^2 = \mathbf{1.0}$ (three dummies can reproduce three
group means exactly). The ordinal code asserted that green sits exactly halfway
between blue and red on the target scale, and that the step blue$\to$green
equals green$\to$red. Both are false --- green is the \emph{highest} group.
A linear model can only add its inputs up, so it believes the assertion and
fits a straight line through means that are not collinear, recovering 7\% of
the variance instead of all of it. A decision tree on the same ordinal column
would score 1.0, because it can split twice and never uses the spacing.
\end{answerbox}

\A{U4.66} \textbf{Verified.} (a) \textbf{5} columns; \textbf{4} with
\texttt{drop='first'}. (b) Intercept plus 5 dummies is a $n\times6$ matrix of
rank \textbf{5} --- the dummies sum to the all-ones column in every row, so one
column is a linear combination of the others. (c) It breaks coefficients,
standard errors, $t$-statistics and $p$-values --- they become non-unique. It
does \textbf{not} break predictions or $R^2$: two entirely different coefficient
vectors gave a maximum prediction difference of 0.00 and identical
$R^2 = 0.986231$. (d) Drop when anyone will \emph{interpret} the coefficients,
or when the model is unregularised OLS. Do not drop for regularised models
(the penalty is no longer symmetric across levels) or for trees (which do not
have an intercept to collide with).

\A{U4.67} \textbf{Verified.} $\binom{d+2}{2}-1$ and $\binom{d+3}{3}-1$:
\begin{center}\begin{tabular}{@{}rrr@{}}\toprule
$d$ & degree 2 & degree 3\\
\midrule
2 & \textbf{5} & 9\\ 10 & \textbf{65} & 285\\ 30 & \textbf{495} & 5455\\
100 & \textbf{5150} & 176850\\
\bottomrule\end{tabular}\end{center}
The count grows like $d^k$. At $d = 30$, degree 2 already produces 495 columns
--- and on the diabetes data, expanding 10 raw features to 65 degree-2 features
made ridge \emph{worse}, $R^2$ falling from 0.4822 to 0.4179. Generating all
interactions is not feature engineering; it is buying variance in bulk and
hoping regularisation pays for it. Engineer a feature because you have a reason.

\A{U4.68} \textbf{Verified.}
\begin{answerbox}
(a) Mean $= (3,3)$. Centred:
$(-2,-2),\ (-1,0),\ (0,-1),\ (1,2),\ (2,1)$.

(b) $\sum x_c^2 = 4+1+0+1+4 = 10$; $\sum y_c^2 = 4+0+1+4+1 = 10$;
$\sum x_cy_c = 4+0+0+2+2 = 8$. With divisor $n-1 = 4$:
$S = \begin{pmatrix} 2.5 & 2.0\\ 2.0 & 2.5\end{pmatrix}$.

(c) The diagonals are equal, so the eigenvectors are
$\frac{1}{\sqrt2}(1,1)$ and $\frac{1}{\sqrt2}(1,-1)$ and the eigenvalues are
$a \pm b = 2.5 \pm 2.0$: $\lambda_1 = \mathbf{4.5}$,
$\lambda_2 = \mathbf{0.5}$. Trace check: $2.5 + 2.5 = 5.0 = 4.5 + 0.5$.
\checkmark

(d) Variance explained: $4.5/5 = \mathbf{0.90}$ and $0.5/5 = \mathbf{0.10}$.

(e) $u_1 = \frac{1}{\sqrt2}(1,1)$, so a score is $(x_c + y_c)/\sqrt2$:
$\mathbf{-2.8284},\ \mathbf{-0.7071},\ \mathbf{-0.7071},\ \mathbf{2.1213},
\ \mathbf{2.1213}$. Their sample variance is
$(8 + 0.5 + 0.5 + 4.5 + 4.5)/4 = 18/4 = \mathbf{4.5} = \lambda_1$. \checkmark

(f) Reconstruction $= \text{score}\cdot u_1 + \text{mean}$:
$(1,1),\ (2.5,2.5),\ (2.5,2.5),\ (4.5,4.5),\ (4.5,4.5)$.
Note $(2,3)$ and $(3,2)$ both collapse onto $(2.5,2.5)$ --- one component
cannot tell them apart.
Total squared error $= 0 + 0.5 + 0.5 + 0.5 + 0.5 = \mathbf{2.0}$, and
$(n-1)\lambda_2 = 4 \times 0.5 = \mathbf{2.0}$. \checkmark
\end{answerbox}

\A{U4.69} \textbf{Verified.} Total $= 10.0$.
\begin{answerbox}
(a) Ratios: $\mathbf{0.48,\ 0.24,\ 0.13,\ 0.08,\ 0.04,\ 0.03}$.

(b) Cumulative: $\mathbf{0.48,\ 0.72,\ 0.85,\ 0.93,\ 0.97,\ 1.00}$.

(c) 80\% needs $k = \mathbf{3}$; 90\% needs $k = \mathbf{4}$; 95\% needs
$k = \mathbf{5}$.

(d) Relative reconstruction error $= 1 - (\text{variance kept})$, so
$\mathbf{0.15},\ \mathbf{0.07},\ \mathbf{0.03}$. No computation is needed
because variance kept and relative error removed \emph{add to exactly one}:
total squared reconstruction error is $(n-1)\sum_{j>k}\lambda_j$ and the total
variance is $(n-1)\sum_j \lambda_j$, so their ratio is
$1 - \sum_{j\le k}\lambda_j/\sum_j\lambda_j$.
\end{answerbox}

\A{U4.70} \textbf{Verified.} (a) All \textbf{three} positives are in the test
set and \textbf{none} in training: the model cannot learn the positive class at
all, and the test set is 75\% positive against a 25\% population. (b)
$P(0) = \binom94/\binom{12}4 = 126/495 = \mathbf{0.2545}$ --- better than one
chance in four. (c) Exactly one positive and three negatives in the test set,
\emph{every} time: $3 \times \frac13 = 1$ and $9 \times \frac13 = 3$.
(d) Recall is $0/0$; scikit-learn returns \textbf{0.0} and emits a warning,
which is easy to miss --- so a fold with no positives silently drags the mean
recall down while accuracy reads 1.0000.

\A{U4.71} \textbf{Verified.}
\begin{answerbox}
(a) $n = 10{,}000$.
accuracy $= 9910/10000 = \mathbf{0.9910}$;
precision $= 10/70 = \mathbf{0.1429}$;
recall $= 10/40 = \mathbf{0.2500}$;
$F_1 = 20/110 = \mathbf{0.1818}$;
balanced $= \tfrac12(9900/9960 + 10/40) = \tfrac12(0.9940 + 0.2500) =
\mathbf{0.6220}$.

(b) Always-negative: $9960/10000 = \mathbf{0.9960}$.

(c) The \emph{always-negative} model scores higher --- 0.9960 against 0.9910.
A model that has learnt something real is beaten on accuracy by a model that
has learnt nothing, because 99.6\% of the answer is ``no''. Accuracy on an
imbalanced problem measures the base rate, not the model.

(d) \textbf{Recall} at a stated threshold, alongside \textbf{PR--AUC} (chance
baseline 0.004 here, not 0.5). Justification: a false negative is a patient
with the disease sent home; a false positive is a patient given a further test.
The costs differ by orders of magnitude, so the headline metric must be one
that penalises missed positives specifically --- and recall of 0.25 says three
in four cases are being missed, which accuracy of 0.991 completely conceals.
\end{answerbox}

\A{U4.72} \textbf{Verified.} $n = 5000$; row-fits $= n(k-1)$.
\begin{answerbox}
\begin{center}\begin{tabular}{@{}lrrrr@{}}\toprule
scheme & fits & rows/fit & row-fits & time\\
\midrule
2-fold & 2 & 2500 & 5{,}000 & 0.5\,s\\
5-fold & 5 & 4000 & 20{,}000 & \textbf{2.0\,s}\\
10-fold & 10 & 4500 & 45{,}000 & 4.5\,s\\
LOOCV & 5000 & 4999 & 24{,}995{,}000 & \textbf{2499.5\,s $\approx$ 42 min}\\
\bottomrule\end{tabular}\end{center}
At $0.4$\,s per 4000-row fit the unit cost is $10^{-4}$\,s per row-fit.

\emph{Choice:} 5-fold. LOOCV costs 1250$\times$ more for an estimate that is
not better --- measured on breast\_cancer, 5-fold and LOOCV both returned
0.9789. Use $k = 5$ or $10$; do not use LOOCV because it sounds thorough.
\end{answerbox}

\A{U4.73} \textbf{Verified.} (a) Validation fold: $5000/5 = \mathbf{1000}$
rows with $\mathbf{80}$ positives. Training portion: $\mathbf{4000}$ rows with
$\mathbf{320}$ positives. (b) Oversampling the 4000 training rows
(3680 majority, 320 minority) to balance gives $\mathbf{7360}$ rows,
$[3680, 3680]$; $3360$ of the $3680$ minority rows are duplicates, i.e.\
$\mathbf{91.3\%}$. (c) Undersampling gives $\mathbf{640}$ rows,
$[320, 320]$, discarding $3360/3680 = \mathbf{91.3\%}$ of the majority.
(d) The validation fold is \textbf{left alone} in both cases --- it keeps its
natural 80/1000 balance, because that is the distribution the model will meet
in deployment and the only one against which a metric means anything.

\A{U4.74} \textbf{Verified.} (a) $B - A = (4, 8)$, so
$A + u(B-A)$ gives $u=0 \to \mathbf{(2,3)}$; $0.25 \to \mathbf{(3,5)}$;
$0.5 \to \mathbf{(4,7)}$; $0.75 \to \mathbf{(5,9)}$; $1 \to \mathbf{(6,11)}$.
(b) The straight \textbf{line segment} joining $A$ and $B$. (c) SMOTE asserts
that the straight line between two minority points is itself full of minority
points. (d) Any categorical or one-hot feature --- a synthetic ``0.4 of the way
from red to blue'' is not a category. Also a multi-modal minority class, where
the segment between two clusters passes through majority territory.

\A{U4.75} \textbf{Verified.} (a)
$\text{SE} = \sqrt{0.91 \times 0.09/44} = \sqrt{0.001861} = \mathbf{0.0431}$.
(b) $1/44 = \mathbf{0.0227}$ --- one row is worth 2.3 accuracy points, so the
score can only move in steps larger than the difference between many real
models. (c) SE $\propto 1/\sqrt{n}$, so halving it needs \textbf{four times}
the test set: $\mathbf{176}$ rows, giving SE $= 0.0216$. (d) With SE $\approx
0.043$, the difference between two models must exceed roughly $0.12$ before a
single split can distinguish them --- and measured over 500 seeds on wine, the
same pair of models produced verdicts ranging from $+0.1556$ to $-0.1333$. You
would be measuring the split, not the models.

\A{U4.76} (a) \textbf{Resampling before the split} --- a leak.
(b) Random oversampling duplicates minority rows; the split then scatters those
duplicates across both sides, so a row the model memorised in training
reappears in the test set. The model is being asked to recognise rows it has
already seen, and it does so perfectly. The measured overstatement was
$F_1\ +0.7638$ for random oversampling and $+0.7809$ for SMOTE --- larger than
any other leak in the course. (c) Close to \textbf{100\%}: when the minority is
oversampled by a factor of roughly 100, almost every minority row in the data
is a copy of some original, so almost every minority test row has its twin in
training. The measured figure was 990 of 990, i.e.\ 100.0\%. (d) Move the
resampling inside the fold --- in practice, use
\texttt{imblearn.pipeline.Pipeline}, because scikit-learn's own
\texttt{Pipeline} cannot hold a step that changes the number of rows.
\end{enumerate}

\subsection*{§C --- Ten-mark questions}

\noindent\small These are the long questions. Mark for coverage, structure and
at least one worked number. \textbf{An answer with no numbers anywhere caps at
6 of 10.}\normalsize

\begin{enumerate}[leftmargin=1.9cm]
\A{U4.77} See the worked mock Section C for
the full four-part breakdown --- this is the same question. Key figures: 1000
validation rows / 80 positives; 4000 training rows / 320 positives; everything
refitted per fold; resampling last, inside the fold; PR--AUC or recall, never
accuracy at a 92\% base rate.

\A{U4.78} Coverage: $(1-p)^d$ with the 0.2146 figure; mean/median/mode and
their situations; the $\sqrt{m/n}$ derivation; MCAR/MAR/MNAR with examples;
MICE recovers 88\% of MAR bias and nothing of MNAR (the sign flip
$+5.603 \to -0.413$ is the memorable case); the indicator and its leakage
condition; and the recommendation --- worry about whether you are \emph{deleting
rows}, not about which imputer, since the imputers differed by 0.003 and
deletion by 5.1 at 30\% missing.

\A{U4.79} Coverage: the three rules and their breakdown points 0\%/25\%/50\%;
masking and swamping; a worked demonstration (U4.56 numbers); the
$(n-1)/\sqrt n$ ceiling \emph{with} the proof; skew is not contamination, so
transform first; the four disposal options. Conclusion required: detection is
statistics, disposal is a judgement you must defend in writing.

\A{U4.80} Coverage: the three mechanisms (distance, gradient/conditioning,
regularisation penalties); the four scalers with formulae and ranges; a worked
neighbour flip (U4.64); which models change and which provably do not, with the
$x \mapsto (x-a)/b$ split argument; the one-typo compression table; and the
training-split rule.

\A{U4.81} Coverage: the maximum-variance question; the four steps; $\lambda_j$
as variance and the trace identity; a full by-hand example (U4.68) including
$(n-1)\lambda_2$; choosing $k$ and why there is no universal 95\% rule (5 of 30
on breast\_cancer, 40 of 64 on digits); standardise first, with the
\texttt{worst area} failure; PCA is not selection (26 of 30 loadings above
0.10); and PCA is unsupervised (0.9867 against 0.5217).

\A{U4.82} Coverage: a single split has SE $\approx 1/\sqrt{n_{\text{te}}}$ and
was measured over 500 seeds at 0.7556 to 1.0000 on the same code; the $k$-fold
recipe and the $n(k-1)$ cost model; $k$ trading bias (training-set size) against
time, with the learning curve flattening; LOOCV's $\sqrt{p(1-p)}$ artefact; the
three failure modes and their splitters; and nested CV, with optimism growing
from 0.0000 at one grid point to $+0.0783$ at sixty.

\A{U4.83} Coverage: the always-negative confusion matrix; accuracy, precision,
recall, $F_1$, balanced accuracy, ROC--AUC and PR--AUC with their chance
baselines; the four remedies; SMOTE's arithmetic and its assumption; the
measured result that resampling raises recall from 0.04 to about 0.78 while
precision falls from 0.13 to 0.05; threshold moving as the cheaper alternative;
and the inside-the-fold rule with the $+0.78$ $F_1$ overstatement.

\A{U4.84} Expected ranking, largest error first: \textbf{resampling before the
split} ($F_1\ +0.78$); \textbf{grouped rows split randomly} ($+0.35$);
\textbf{feature selection on all the data} ($+0.31$); \textbf{time-ordered rows
split randomly} ($R^2\ +3.81$, a ranking inversion); \textbf{row-dropping using
the model then splitting} ($+0.17$); \textbf{scaler on all the data}
($+0.0034$ falling to $+0.0002$ as $n$ grows); \textbf{imputer on all the data}
($+0.0003$, not distinguishable from zero). The correct construction in each
case: put the step inside a \texttt{Pipeline} and let the cross-validator refit
it per fold --- except row-dropping, which cannot live in a
\texttt{Pipeline} at all, and that is exactly why it is dangerous.

\A{U4.85} Coverage: filter/wrapper/embedded with two methods each and the
cost--risk of each; a filter measures the dependence it was built to measure
($F$ vs MI, with the $|b|$ feature ranked 1st and 8th); RFE costs one fit per
elimination; L1 selects and L2 shrinks; the Jaccard overlaps as low as 0.20
between selectors; and the honest finding that on breast\_cancer \emph{not one}
selector beat using all thirty features.

\A{U4.86} Coverage: what feature engineering is; combine/decompose/bin/
transform; the four-plots-of-land example where $\text{corr}(\text{price},
\text{width}) = -0.0948$ but the product column takes $R^2$ from 0.6558 to
1.0000; binning helps a linear model $+0.0914$ and harms a tree $-0.0534$, so a
representation is good only relative to a model; and the combinatorial
blow-up as the reason automatic generation is not a substitute for a reason.
\end{enumerate}

\subsection*{§D --- Statements}

\begin{enumerate}[leftmargin=1.9cm]
\A{U4.87} \textbf{This is the definition, not an answer.}
If the question asked what the rule \emph{does}, it is true but empty; if the
question asked anything about its behaviour, it does not address the question. The substantive
points are the 0\% breakdown point, masking, and the $(n-1)/\sqrt n$ ceiling.
\A{U4.88} \textbf{False.} PCA keeps \emph{linear combinations} of all features,
not a subset. 26 of 30 loadings on PC1 exceeded 0.10. You still need every
original measurement to compute the components.
\A{U4.89} \textbf{False.} Provably no change: a split $x_j \le t$ becomes
$(t-a)/b$ under a positive affine rescaling and sends the same rows left.
Measured: identical structure, identical predictions, accuracy 0.902098 both
ways.
\A{U4.90} \textbf{False.} It preserves the mean and destroys the variance:
$\sigma$ falls by $\sqrt{1-p}$, so every standard error is too small and every
correlation is attenuated (0.5865 $\to$ 0.4536 at 40\% missing). It manufactures
confidence you do not have.
\A{U4.91} \textbf{False.} Measured: 257 levels in 600 rows gave train $R^2$
0.9356 and test 0.6984, against 0.8081 without the id column --- the encoding
made the model \emph{worse}. 82 levels appeared in training and never in test.
\A{U4.92} \textbf{False as stated} --- and this one matters, because the
overclaim is common. Fitting the scaler on everything is genuinely wrong and
should be fixed, but the measured effect was $+0.0034$ at 20 training rows and
$+0.0002$ at 427. Report what you measured, not what you expected. The large
leaks come from steps that see $y$.
\A{U4.93} \textbf{False.} It costs $n$ fits (98$\times$ 5-fold for the same
0.9789), and its across-fold sd of 0.1437 is $\sqrt{p(1-p)}$ --- an artefact of
scoring one row at a time, carrying no information. Use $k = 5$ or 10.
\A{U4.94} \textbf{False.} SMOTE interpolates between existing minority points:
it \emph{asserts} that the segment between two minority examples is full of
minority examples. On categorical features or a multi-modal minority class that
assertion is simply wrong.
\A{U4.95} \textbf{False.} An always-negative model scores 0.99 on the same
problem and has learnt nothing. Measured, a real logistic regression scored
0.9867 --- \emph{worse} than the dummy's 0.9900 --- while its PR--AUC of 0.1075
against a chance baseline of 0.0100 showed it had learnt a great deal.
\A{U4.96} \textbf{False.} It is a coin toss with your result on it. Over 500
seeds on wine the same model scored between 0.7556 and 1.0000 --- a range of
0.2444.
\end{enumerate}

\vspace{0.4em}\hrule

% =====================================================================
\section*{Unit V (part) --- Regression \quad (L8, L9, L10)}

\begin{warnbox}[Scope, and one thing to watch]
Covered: L8 least squares, L9 maximum likelihood / $R^2$ / adequacy /
over-fitting, L10 logistic regression. \textbf{Not covered:} multiple linear
regression (L11) and regularization (L12).

\emph{Note:} this material sits in L8 and L9, and
L10 to \textbf{T5}. A mid-semester paper spanning L8--L10 therefore cuts across
that mapping. Nothing here is wrong, but if a student asks why the bank groups
L10 with L8--L9 when T5 groups it with L11--L12, the answer is that the tests
follow the teaching order and the mid-semester paper follows the calendar.
\end{warnbox}

\subsection*{§A --- Two-mark questions}

\begin{enumerate}[leftmargin=1.9cm]
\A{U5.1} $y_i = \beta_0 + \beta_1 x_i + \varepsilon_i$. $y$ response,
$x$ predictor, $\beta_0$ intercept, $\beta_1$ slope, $\varepsilon_i$ error.
\textbf{Greek} letters are the unknown \emph{population} parameters and the
unobservable errors; \textbf{Roman} ($b_0$, $b_1$, $e_i$) are the
\emph{estimates} and the observed residuals. It matters because writing
``$\varepsilon_i = 0.6$'' claims to have observed something unobservable.

\A{U5.2} Linear \textbf{in the parameters}, not in $x$. (a) linear;
(b) \textbf{not} --- $\beta_0\beta_1$ is a product of parameters; (c) linear;
(d) not.

\A{U5.3} $\text{SSE}(b_0,b_1) = \sum_i (y_i - b_0 - b_1x_i)^2$ --- the sum of
squared \textbf{vertical} deviations of the observed $y$ from the fitted line.
\emph{Insist on ``squared vertical''; ``the distance from the line'' is the
perpendicular distance and is a different method.}

\A{U5.4} $n b_0 + b_1\sum x_i = \sum y_i$ and
$b_0\sum x_i + b_1\sum x_i^2 = \sum x_iy_i$.

\A{U5.5} $b_1 = S_{xy}/S_{xx}$, $b_0 = \bar y - b_1\bar x$, where
$S_{xx} = \sum(x_i-\bar x)^2$ and $S_{xy} = \sum(x_i-\bar x)(y_i-\bar y)$.

\A{U5.6} Any three: the derivative of a square is linear, so the first-order
conditions form a \emph{linear} system with a closed-form solution;
differentiable everywhere; the estimator is unique whenever $S_{xx} > 0$;
it is the maximum likelihood estimator under Gaussian noise; and Gauss--Markov
minimum variance among linear unbiased estimators. Given up: robustness ---
squaring hands a single bad row disproportionate influence.

\A{U5.7} Because the model treats $x$ as fixed and known and $y$ as the thing
with error; the residual is an error in $y$ alone. Minimising perpendicular
distance gives \textbf{total least squares} (orthogonal / major-axis
regression) --- which is the first principal component of the data, linking
back to L6.

\A{U5.8} $H = 2\begin{pmatrix} n & \sum x_i\\ \sum x_i & \sum x_i^2
\end{pmatrix}$, $\det H = 4nS_{xx}$. Positive definite iff $n > 0$ and
$S_{xx} > 0$ --- that is, iff the $x_i$ are \textbf{not all equal}. Then SSE is
strictly convex and the stationary point is the unique global minimum.

\A{U5.9} Any dataset with all $x$ identical, e.g. $x = (3,3,3,3,3)$. Then
$S_{xx} = 0$, $\det(X^\top X) = 0$, and $b_1 = S_{xy}/S_{xx}$ is undefined ---
infinitely many lines fit equally well.

\A{U5.10} $X^\top X\, b = X^\top y$, so $b = (X^\top X)^{-1}X^\top y$, provided
$X$ has \textbf{full column rank}.
\emph{Worth adding: in practice you never compute the inverse at all --- \texttt{lstsq} via SVD, never \texttt{np.linalg.inv}.}

\A{U5.11} $\sum e_i = 0$; the line passes through $(\bar x, \bar y)$;
$\sum x_i e_i = 0$; and $\text{SST} = \text{SSR} + \text{SSE}$. The first
requires an intercept --- it \emph{is} the first normal equation.

\A{U5.12} The first normal equation is
$\partial\text{SSE}/\partial b_0 = -2\sum(y_i - b_0 - b_1x_i) = -2\sum e_i = 0$,
hence $\sum e_i = 0$ directly.

\A{U5.13} From $\sum e_i = 0$: $\sum(y_i - b_0 - b_1x_i) = 0$, so
$n\bar y = nb_0 + b_1 n\bar x$, i.e. $\bar y = b_0 + b_1\bar x$ --- which says
the point $(\bar x, \bar y)$ satisfies the fitted line.

\A{U5.14} $\sum(y_i-\bar y)^2 = \sum(\hat y_i-\bar y)^2 + \sum e_i^2$, i.e.
SST (total) $=$ SSR (explained by the regression) $+$ SSE (residual). The
cross term $2\sum(\hat y_i - \bar y)e_i$ vanishes by the two normal equations.

\A{U5.15} $\operatorname{Var}(b_1) = \sigma^2/S_{xx}$. Therefore
\textbf{spread the $x$ values out} --- put observations at the extremes of the
usable range. The slope is estimated more precisely the larger $S_{xx}$ is, and
this costs nothing extra per observation.

\A{U5.16} $h_i = 1/n + (x_i-\bar x)^2/S_{xx}$; the $h_i$ sum to $p$ (2 here).
Flag $h_i > 2p/n$. Leverage is about \emph{where the point sits in $x$}, not
about whether its $y$ is unusual.

\A{U5.17} Closed form: $O(nd^2 + d^3)$ once, needs $X^\top X$ or all of $X$ in
memory, no tuning, exact answer, fails when $d$ is large or $X^\top X$ is
singular. Gradient descent: $O(nd)$ per step over many steps, one mini-batch in
memory, requires a step size / schedule / batch size / stopping rule, gets as
close as you can afford, almost never fails outright. Closed form extends only
to squared loss on a linear model; GD extends to anything differentiable.

\A{U5.18} $y_i = \beta_0 + \beta_1x_i + \varepsilon_i$ with
$\varepsilon_i \sim \mathcal N(0,\sigma^2)$ \textbf{independently and
identically}, and $x_i$ fixed and known. Equivalently
$y_i \sim \mathcal N(\beta_0 + \beta_1x_i,\ \sigma^2)$.

\A{U5.19} $L = (2\pi\sigma^2)^{-n/2}\exp\!\left(-\frac{1}{2\sigma^2}\sum e_i^2
\right)$;
$\ell = -\frac n2\log(2\pi) - \frac n2\log\sigma^2 - \frac{1}{2\sigma^2}S$,
with $S = \sum(y_i - \beta_0 - \beta_1x_i)^2$.

\A{U5.20} \emph{The sentence to be able to write from memory:}
``\textbf{Least squares is the maximum likelihood estimator of the linear model
under the assumption of independent Gaussian errors with constant
variance.}''

\A{U5.21} Because $\ell = A - cS$ where $A$ does not involve $\beta_0,\beta_1$
and $c = 1/(2\sigma^2) > 0$. Maximising $A - cS$ over $(\beta_0,\beta_1)$ is
minimising $S$, for \emph{every} positive $\sigma^2$ --- the value of $\sigma^2$
scales the objective but cannot move its argmin.

\A{U5.22} $\hat\sigma^2_{\text{ML}} = \text{SSE}/n$; unbiased
$s^2 = \text{SSE}/(n-2)$. \textbf{Reason:} the residuals satisfy two linear
constraints ($\sum e_i = 0$ and $\sum x_ie_i = 0$), so only $n-2$ of them are
free; equivalently $E[\text{SSE}] = (n-2)\sigma^2$.

\A{U5.23} $\ell = -n\log(2b) - \frac1b\sum|e_i|$, so maximising it minimises
$\sum|e_i|$ --- \textbf{least absolute deviations}. Change the noise model and
you change the estimator; the recipe, not the squares, is the general fact.

\A{U5.24} $R^2 = 1 - \text{SS}_{\text{res}}/\text{SS}_{\text{tot}} =
\text{SS}_{\text{reg}}/\text{SS}_{\text{tot}}$. On training data with an
intercept it lies in $[0,1]$, because the fit can always fall back on the
constant $\bar y$. On held-out data it has \textbf{no lower bound}: the model
was not allowed to see $\bar y_{\text{test}}$, so it can do arbitrarily worse
than predicting it.

\A{U5.25} Because the larger model \emph{contains} the smaller one --- setting
the new coefficient to zero reproduces the old fit exactly, so the minimised
SS$_{\text{res}}$ cannot rise.

\A{U5.26} $R^2_{\text{adj}} = 1 - (1-R^2)\frac{n-1}{n-p-1}$. It \emph{is} a
penalty for parameters that makes nested models roughly comparable. It is
\textbf{not} a hypothesis test, not a substitute for held-out validation, and
it can still be fooled --- on 30 rows with 27 noise columns it rose again at
$p = 21$ while the cross-validated score collapsed.

\A{U5.27} Two of: it is a training-set number and says nothing about
generalisation; it cannot distinguish a good fit from a wrong model, a
contaminated dataset or a fit resting on one point (Anscombe); it never falls
when predictors are added, so it rewards complexity; and it has no units and no
comparison across different $y$.

\A{U5.28} Linearity --- fails, and the \emph{coefficients} are meaningless.
Constant variance --- fails, and standard errors, $t$, $F$, $p$ and all
intervals are invalid (the coefficients stay unbiased). Normality --- fails,
and $p$-values and intervals are unreliable in small samples. Independence ---
fails, and standard errors are wrong, usually understated.

\A{U5.29} Fitting the noise rather than the signal: training error falls while
error on new data rises. It depends on both because the \emph{same} degree-20
polynomial gave a test RMSE of 53{,}293 at $n = 15$ and 0.2545 at $n = 1000$.
Capacity is only over-capacity relative to the amount of data.

\A{U5.30} (i) The fitted line is unbounded, so it produces ``probabilities''
outside $[0,1]$ --- a predicted $-16.7\%$. (ii) A single correct but distant
point drags the decision boundary, because squared error keeps charging for
predictions that are ``too right''.

\A{U5.31} Odds $= p/(1-p) \in (0,\infty)$; log-odds
$= \ln\!\big(p/(1-p)\big) \in (-\infty,\infty)$; $p \in (0,1)$. Only the
\textbf{log-odds} has the range of a linear expression, which is exactly why
the model is built there.

\A{U5.32} $\ln\!\frac{p(x)}{1-p(x)} = w^\top x + b = z$. Inverting:
$\frac{p}{1-p} = e^z \Rightarrow p = e^z - pe^z \Rightarrow p(1+e^z) = e^z
\Rightarrow p = \frac{e^z}{1+e^z} = \frac{1}{1+e^{-z}} = \sigma(z)$.

\A{U5.33} $\sigma(0) = 1/2$ (so the boundary is the hyperplane $z = 0$);
$\sigma(-z) = 1-\sigma(z)$; $\sigma'(z) = \sigma(z)(1-\sigma(z))$, maximal at
$z = 0$ with value $0.25$.

\A{U5.34} $\ell = \sum_i[y_i\ln p_i + (1-y_i)\ln(1-p_i)]$. Then $-\ell/n$ is
exactly the binary cross-entropy (\texttt{log\_loss}): maximising the
likelihood and minimising cross-entropy are the same operation.

\A{U5.35} $\nabla_w\ell = X^\top(y - p)$, equivalently
$\nabla_w L = X^\top(p - y)$ for the negative log-likelihood. In words: how
strongly does each feature line up with the errors I am currently making ---
\emph{prediction minus truth}, projected onto the features.

\A{U5.36} The score equations are $X^\top(\sigma(Xw) - y) = 0$ --- the unknown
sits inside a transcendental function, like $xe^x = 1$, and cannot be isolated.
The normal equations $X^\top Xb = X^\top y$ are \emph{linear} in $b$ and so
solve in closed form.

\A{U5.37} $H = -X^\top S X$ with $S = \operatorname{diag}(p_i(1-p_i))$. For
$v \ne 0$, $v^\top H v = -\sum_i p_i(1-p_i)(x_i^\top v)^2 < 0$ when $X$ has
full column rank, so $\ell$ is strictly concave. That buys \textbf{uniqueness}
of the maximum and freedom from local optima --- it does \textbf{not} buy
\emph{existence}: under perfect separation the supremum is approached only as
$\|w\| \to \infty$.

\A{U5.38} $e^{\beta_j}$ is the \textbf{multiplicative change in the odds} for a
one-unit increase in $x_j$, holding everything else fixed --- and it is the
same factor whatever the starting odds. It is \emph{not} a change in the
probability.

\A{U5.39} $\left.\frac{\mathrm dp}{\mathrm dx_j}\right|_{p=0.5} =
\beta_j\sigma'(0) = \beta_j/4$. Since $\sigma' \le 0.25$ everywhere, $|\beta_j|/4$
is an \textbf{upper bound} on the change in probability per unit of $x_j$,
attained only at $p = 0.5$.

\A{U5.40} The model outputs a probability; turning it into a decision requires
knowing the cost of each kind of error, which is not in the data. The
threshold should be chosen by whoever owns those costs --- the clinician, the
bank, the business --- on a \emph{validation} set. Choosing it on the test set
is Lecture-7 leakage in a new hat.

\A{U5.41} $P(y=k\mid x) = e^{z_k}/\sum_{m=1}^{K}e^{z_m}$ with
$z_k = w_k^\top x + b_k$. For $K = 2$, fix $z_0 = 0$ (legitimate, since adding
a constant to every $z$ cancels): $P(y=1) = e^{z_1}/(1 + e^{z_1}) =
\sigma(z_1)$.

\A{U5.42} A hyperplane separates the classes perfectly. The likelihood
increases without bound as $\|w\| \to \infty$ --- probabilities are pushed to
0 and 1 --- so the MLE \textbf{does not exist} (the supremum is not attained).
Fix: an L2 penalty, $\min L(w) + \frac{1}{2C}\|w\|^2$, which is
scikit-learn's default and the subject of L12.

\A{U5.43} (i) It is non-convex in $w$ when composed with the sigmoid, so the
optimisation loses its single-optimum guarantee. (ii) The gradient carries a
$\sigma'$ factor which saturates to zero for badly wrong confident
predictions --- the model stops learning exactly where it is most wrong.
Cross-entropy's $(y-p)$ gradient has no such factor.
\end{enumerate}

\subsection*{§B --- Five-mark questions}

\begin{enumerate}[leftmargin=1.9cm]

\A{U5.44} \textbf{Verified.}
\begin{answerbox}
(a) $\bar x = 2$, $\bar y = 4$.
\begin{center}\begin{tabular}{@{}rrrrr@{}}\toprule
$x$ & $y$ & $x-\bar x$ & $y-\bar y$ & $(x-\bar x)^2$ \quad $(x-\bar x)(y-\bar y)$\\
\midrule
0 & 1 & $-2$ & $-3$ & 4 \quad\quad\quad 6\\
1 & 2 & $-1$ & $-2$ & 1 \quad\quad\quad 2\\
2 & 4 & 0 & 0 & 0 \quad\quad\quad 0\\
3 & 5 & 1 & 1 & 1 \quad\quad\quad 1\\
4 & 8 & 2 & 4 & 4 \quad\quad\quad 8\\
\midrule
& & & & $S_{xx} = \mathbf{10}$ \quad $S_{xy} = \mathbf{17}$\\
\bottomrule\end{tabular}\end{center}
$S_{yy} = 9+4+0+1+16 = \mathbf{30}$.

(b) $b_1 = 17/10 = \mathbf{1.7}$; $b_0 = 4 - 1.7(2) = \mathbf{0.6}$;
$\hat y = 0.6 + 1.7x$.

(c) Fitted $\mathbf{0.6,\ 2.3,\ 4.0,\ 5.7,\ 7.4}$;
residuals $\mathbf{0.4,\ -0.3,\ 0.0,\ -0.7,\ 0.6}$.
$\sum e_i = \mathbf{0}$; $\sum x_ie_i = 0(0.4)+1(-0.3)+2(0)+3(-0.7)+4(0.6)
= -0.3-2.1+2.4 = \mathbf{0}$.

(d) $\text{SSE} = 0.16+0.09+0+0.49+0.36 = \mathbf{1.10}$;
$\text{SSR} = \mathbf{28.90}$; $\text{SST} = \mathbf{30.00}$;
$28.90 + 1.10 = 30.00$ \checkmark.
Checks: $\text{SSR} = b_1^2S_{xx} = 2.89(10) = 28.90$ \checkmark\ and
$= S_{xy}^2/S_{xx} = 289/10 = 28.90$ \checkmark.

(e) $\hat y(6) = 0.6 + 1.7(6) = \mathbf{10.8}$.
\end{answerbox}

\A{U5.45} \textbf{Verified.}
\begin{answerbox}
(a) $X = \begin{pmatrix}1&0\\1&1\\1&2\\1&3\\1&4\end{pmatrix}$,
$X^\top X = \begin{pmatrix}5 & 10\\ 10 & 30\end{pmatrix}$,
$X^\top y = \begin{pmatrix}20\\ 57\end{pmatrix}$.

(b) $\det = 5(30) - 10^2 = \mathbf{50}$, and $nS_{xx} = 5(10) = \mathbf{50}$
\checkmark.

(c) $(X^\top X)^{-1} = \frac{1}{50}\begin{pmatrix}30 & -10\\ -10 & 5
\end{pmatrix} = \begin{pmatrix}0.6 & -0.2\\ -0.2 & 0.1\end{pmatrix}$;
$b = \frac{1}{50}\begin{pmatrix}30(20)-10(57)\\ -10(20)+5(57)\end{pmatrix}
= \frac{1}{50}\begin{pmatrix}30\\ 85\end{pmatrix} = \mathbf{(0.6,\ 1.7)}$
--- identical to U5.44.

(d) Fails iff $X$ is rank deficient, i.e. $S_{xx} = 0$, i.e. every $x_i$ equal:
e.g. $x = (3,3,3,3,3)$, giving $\det(X^\top X) = 0$.
\end{answerbox}

\A{U5.46} The six-part derivation. Mark it in six parts, 0.75--1 each:
\begin{answerbox}
(1) Object: $\text{SSE}(b_0,b_1) = \sum(y_i - b_0 - b_1x_i)^2$.
(2) $\partial/\partial b_0 = -2\sum(y_i-b_0-b_1x_i)$;
$\partial/\partial b_1 = -2\sum x_i(y_i-b_0-b_1x_i)$.
(3) Setting both to zero: $nb_0 + b_1\sum x_i = \sum y_i$ and
$b_0\sum x_i + b_1\sum x_i^2 = \sum x_iy_i$.
(4) First equation $\Rightarrow b_0 = \bar y - b_1\bar x$; substituting into
the second gives $b_1(\sum x_i^2 - n\bar x^2) = \sum x_iy_i - n\bar x\bar y$.
(5) Recognise $\sum x_i^2 - n\bar x^2 = S_{xx}$ and
$\sum x_iy_i - n\bar x\bar y = S_{xy}$, so $b_1 = S_{xy}/S_{xx}$.
(6) Second-order: $H = 2\begin{pmatrix}n & \sum x_i\\ \sum x_i & \sum x_i^2
\end{pmatrix}$, $\det H = 4nS_{xx} > 0$ and $2n > 0$, so $H$ is positive
definite, SSE is strictly convex, and the stationary point is the unique
global minimum.
\end{answerbox}
\emph{The expansion identities in (5) have to be shown or quoted, not assumed --- however neat the work, an answer that stops after (3) has not finished the question.}

\A{U5.47} \textbf{Verified.}
\begin{answerbox}
\begin{center}\begin{tabular}{@{}llrrr@{}}\toprule
& line & SSE & $\sum e_i$ & $\sum x_ie_i$\\
\midrule
A & $0.5 + 1.7x$ & 1.15 & $+0.5$ & $+1.0$\\
\textbf{B} & $\mathbf{0.6 + 1.7x}$ & \textbf{1.10} & \textbf{0.0} & \textbf{0.0}\\
C & $0.6 + 1.8x$ & 1.40 & $-1.0$ & $-3.0$\\
D & $1.0 + 1.5x$ & 1.50 & \textbf{0.0} & $+2.0$\\
\bottomrule\end{tabular}\end{center}
\textbf{B} is the least-squares line: it is the only candidate satisfying
\emph{both} normal equations.

\textbf{Why one is not enough:} candidate \textbf{D} has $\sum e_i = 0$ exactly
--- it passes the first test cleanly --- yet $\sum x_ie_i = +2.0$ and its SSE
is 36\% worse than B's. $\sum e_i = 0$ says only that the line passes through
$(\bar x,\bar y)$; infinitely many lines do. The second condition is what pins
down the slope.
\end{answerbox}

\A{U5.48} \textbf{Verified.}
\begin{answerbox}
(a) $g(b_1) = \sum(y_i - b_1x_i)^2$; $g'(b_1) = -2\sum x_i(y_i-b_1x_i) = 0
\Rightarrow b_1\sum x_i^2 = \sum x_iy_i \Rightarrow b_1 = \sum x_iy_i/\sum x_i^2$.
$g''(b_1) = 2\sum x_i^2 > 0$, so it is a minimum.

(b) $b_1 = 57/30 = \mathbf{1.9}$.

(c) Residuals $\mathbf{1.0,\ 0.1,\ 0.2,\ -0.7,\ 0.4}$;
$\sum e_i = \mathbf{+1.0}$, \emph{not} zero.

(d) There is no $b_0$, so there is no $\partial/\partial b_0$ equation --- and
$\sum e_i = 0$ \emph{is} that equation. Only $\sum x_ie_i = 0$ survives (check:
$0(1.0)+1(0.1)+2(0.2)+3(-0.7)+4(0.4) = 0$ \checkmark).

(e) SSE rises from $1.10$ to $\mathbf{1.70}$. The comparison is fair in the
sense that both minimise the same criterion over their own model class, and the
no-intercept class is a strict subset --- so its SSE can never be smaller. It
would be unfair to read the gap as evidence about the data rather than about
the constraint.
\end{answerbox}

\A{U5.49} \textbf{Verified.}
\begin{answerbox}
(a) $h_i = 0.2 + (x_i-2)^2/10$:
$\mathbf{0.6,\ 0.3,\ 0.2,\ 0.3,\ 0.6}$; sum $= \mathbf{2.0} = p$ \checkmark.

(b) $b_1 = \sum(x_i-\bar x)(y_i-\bar y)/S_{xx}$. Adding $\delta$ to $y_i$ adds
$\delta$ to $y_i$ and $\delta/n$ to $\bar y$; the $\bar y$ shift contributes
$-\frac{\delta}{n}\sum(x_j-\bar x) = 0$, leaving
$\Delta b_1 = \delta(x_i-\bar x)/S_{xx}$.

(c) $\delta = 3$, $\Delta b_1 = 3(x_i-2)/10$:
\begin{center}\begin{tabular}{@{}rrr@{}}\toprule
$x_i$ & $\Delta b_1$ & new $b_1$\\
\midrule
0 & $-0.6$ & 1.1\\ 1 & $-0.3$ & 1.4\\ 2 & \textbf{0.0} & \textbf{1.7}\\
3 & $+0.3$ & 2.0\\ 4 & $+0.6$ & 2.3\\
\bottomrule\end{tabular}\end{center}

(d) The point at $x = 2 = \bar x$. Its deviation $(x_i - \bar x)$ is zero, so
it carries no weight in $S_{xy}$ at all --- you may move its $y$ anywhere and
the slope will not shift by a hair. (The intercept does move.) This is the
cleanest demonstration that influence is about position in $x$, not about
having an unusual $y$.
\end{answerbox}

\A{U5.50} \textbf{Verified.} $\sigma^2 = 4$.
\begin{answerbox}
\begin{center}\begin{tabular}{@{}lrrr@{}}\toprule
design & $S_{xx}$ & $\operatorname{Var}(b_1)$ & $\operatorname{sd}(b_1)$\\
\midrule
three at $x=0$, two at $x=8$ & \textbf{76.8} & \textbf{0.0521} & \textbf{0.2282}\\
equally spaced $0..8$, $n=5$ & 40.0 & 0.1000 & 0.3162\\
clustered on $[3,5]$, $n=5$ & 2.5 & 1.6000 & 1.2649\\
\bottomrule\end{tabular}\end{center}
Ranking: split-ends best, equally spaced second, clustered worst by a factor of
$1.2649/0.2282 = 5.5$ in standard deviation --- for the same five
measurements.

\textbf{What the best design cannot do:} with observations at only two $x$
values it cannot detect \emph{curvature}. Two points determine a line, and any
non-linearity between them is invisible. Precision on the slope has been bought
by giving up the ability to check whether a slope was the right model at all.
\end{answerbox}

\A{U5.51} Mark in five parts:
\begin{answerbox}
(1) Model: $y_i \sim \mathcal N(\beta_0+\beta_1x_i,\ \sigma^2)$ independently.
(2) One row: $p(y_i) = (2\pi\sigma^2)^{-1/2}\exp(-e_i^2/(2\sigma^2))$.
(3) Independence $\Rightarrow$ product:
$L = (2\pi\sigma^2)^{-n/2}\exp\!\big(-\frac{1}{2\sigma^2}\sum e_i^2\big)$.
(4) Logs: $\ell = -\frac n2\log(2\pi) - \frac n2\log\sigma^2 -
\frac{1}{2\sigma^2}S$.
(5) Only the last term contains $\beta_0,\beta_1$; it enters with a negative
sign and a positive coefficient $1/(2\sigma^2)$. Writing $\ell = A - cS$ with
$c > 0$: maximising $\ell$ is minimising $S$, \emph{for every} $\sigma^2 > 0$.
\end{answerbox}

\A{U5.52} \textbf{Verified.} SSE $= 1.10$, $n = 5$.
\begin{answerbox}
(a) $\hat\sigma^2_{\text{ML}} = 1.10/5 = \mathbf{0.22}$.
(b) $s^2 = 1.10/3 = \mathbf{0.3667}$; $s = \sqrt{0.3667} = \mathbf{0.6055}$.
(c) Bias factor $(n-2)/n = 3/5 = \mathbf{0.6}$; the ML estimate understates
$\sigma^2$ by $2/n = \mathbf{40\%}$ in expectation.
(d) The residuals are not $n$ free numbers: they satisfy $\sum e_i = 0$ and
$\sum x_ie_i = 0$, two linear constraints imposed by fitting two parameters.
Only $n-2 = 3$ of the five residuals can be chosen freely, so
$E[\text{SSE}] = (n-2)\sigma^2$ and dividing by $n$ systematically
under-counts.
\end{answerbox}

\A{U5.53} \textbf{Verified.}
\begin{answerbox}
\begin{center}\begin{tabular}{@{}rrrrrr@{}}\toprule
$n$ & SSE & $\hat\sigma^2_{\text{ML}}$ & $s^2$ & $(n-2)/n$ & understatement\\
\midrule
8 & 20 & 2.5000 & 3.3333 & 0.7500 & 25.00\%\\
6 & 12 & 2.0000 & 3.0000 & 0.6667 & 33.33\%\\
12 & 45 & 3.7500 & 4.5000 & 0.8333 & 16.67\%\\
\bottomrule\end{tabular}\end{center}
Understatement is $2/n$, independent of SSE. Below 1\% requires
$2/n < 0.01$, i.e. $n > 200$, so $\mathbf{n \ge 201}$.
\end{answerbox}

\A{U5.54} \textbf{Verified.}
\begin{answerbox}
(a) SS$_{\text{tot}} = \mathbf{30.00}$, SS$_{\text{res}} = \mathbf{1.10}$,
SS$_{\text{reg}} = \mathbf{28.90}$.
(b) $R^2 = 1 - 1.10/30.00 = \mathbf{0.9633}$.
(c) $\text{SS}_{\text{reg}}/\text{SS}_{\text{tot}} = 28.90/30 = 0.9633$
\checkmark; $b_1^2S_{xx}/S_{yy} = 2.89(10)/30 = 0.9633$ \checkmark.
(d) $r = S_{xy}/\sqrt{S_{xx}S_{yy}} = 17/\sqrt{300} = \mathbf{0.9815}$;
$r^2 = 0.9633 = R^2$ \checkmark.
(e) It does not tell you whether a \emph{line} was the right model, whether one
point is carrying the fit, or whether the variance is constant --- Anscombe's
four datasets all have $R^2 \approx 0.666$ and four different diagnoses. Look
at the \textbf{residual plot}.
\end{answerbox}

\A{U5.55} \textbf{Verified.} $n = 25$, $R^2 = 0.55$.
\begin{answerbox}
$R^2_{\text{adj}} = 1 - 0.45\cdot\frac{24}{24-p}$:
\begin{center}\begin{tabular}{@{}rr@{}}\toprule
$p$ & $R^2_{\text{adj}}$\\
\midrule
1 & $\mathbf{0.5304}$\\ 5 & 0.4316\\ 10 & 0.2286\\ 15 & $-0.2000$\\
\bottomrule\end{tabular}\end{center}
Negative when $0.45\cdot\frac{24}{24-p} > 1$, i.e. $24 - p < 10.8$, i.e.
$p > 13.2$: the first integer is $\mathbf{p = 14}$, giving $-0.0800$
($p = 13$ still gives $+0.0182$).

\textbf{Interpretation:} once the penalty for parameters exceeds all the
variance the model explains, adjusted $R^2$ goes below the value a
constant-only model would score. It is saying that on a per-parameter basis
this model has bought nothing --- with 14 predictors on 25 rows you are fitting
noise, and the un-adjusted $R^2$ of 0.55 is an illusion of the parameter count.
\end{answerbox}

\A{U5.56} \textbf{Verified.}
\begin{answerbox}
(a) $\bar y_{\text{test}} = 4$; SS$_{\text{tot}} = 4+0+4 = \mathbf{8}$;
SS$_{\text{res}} = 25+0+25 = \mathbf{50}$;
$R^2 = 1 - 50/8 = \mathbf{-5.25}$.
(b) Constant 4: SS$_{\text{res}} = 4+0+4 = 8$, so $R^2 = 1 - 8/8 = \mathbf{0}$.
(c) \textbf{None} --- held-out $R^2$ is unbounded below.
(d) $R^2 \in [0,1]$ on training data holds because the fitted model
\emph{could} have chosen the constant $\bar y$ and did at least as well; the
constant is inside the model's reach. On held-out data $\bar y_{\text{test}}$
is a quantity the model never saw, so there is no guarantee it beats it. A
negative held-out $R^2$ means precisely: you would have done better predicting
the mean of the test set for every row.
\end{answerbox}

\A{U5.57} \textbf{Verified.} $x = (1,2,3,4,5)$, $y = (2,3,6,5,9)$.
\begin{answerbox}
$\bar x = 3$, $\bar y = 5$; $S_{xx} = \mathbf{10}$, $S_{xy} = \mathbf{16}$,
$S_{yy} = \mathbf{30}$.
Direct slope $= 16/10 = \mathbf{1.6}$; inverse slope
$= S_{yy}/S_{xy} = 30/16 = \mathbf{1.875}$.
Ratio $= 1.6/1.875 = \mathbf{0.8533}$; and
$r^2 = S_{xy}^2/(S_{xx}S_{yy}) = 256/300 = \mathbf{0.8533}$ \checkmark.
($r = 0.9238$.)

\textbf{Which to report:} the \textbf{direct} line, because the question is to
predict $y$ from $x$ and the direct fit is the one that minimises error in $y$.
The two coincide only when $r^2 = 1$.

\textbf{Inverting the direct fit is wrong:} $1/1.6 = 0.625$, but the correct
slope for predicting $x$ from $y$ is $S_{xy}/S_{yy} = 16/30 = 0.533$. The gap
is a factor of $1/r^2$. Regression is not symmetric --- there is a different
line for each direction, and choosing which variable carries the error is part
of stating the problem.
\end{answerbox}

\A{U5.58} (a) Faults: reporting a \emph{training} score as evidence of quality;
285 parameters against 331 rows, so the model can very nearly interpolate;
no held-out or cross-validated number quoted alongside; and treating a high
$R^2$ as sufficient with no residual diagnostic. (b) A degree-3 expansion of 10
features has $\binom{13}{3} - 1 = 285$ terms against 331 rows --- roughly one
parameter per row. (c) Its squared error is about \textbf{23 times} that of
predicting a single constant ($\text{SS}_{\text{res}}/\text{SS}_{\text{tot}}
= 1-(-22.04) = 23.04$); it has learned the training noise and that noise does not recur.
(d) The cross-validated score, the held-out score, the parameter count against
$n$, and the trivial baseline for comparison ($-0.0001$ for the training mean)
--- with the model rejected in favour of the plain linear fit.

\A{U5.59} \textbf{Verified.}
\begin{answerbox}
(a) $\bar x = 3.5$, $\bar y = 0.5$; $S_{xx} = \mathbf{17.5}$,
$S_{xy} = \mathbf{4.5}$; slope $= 4.5/17.5 = \mathbf{0.2571}$;
intercept $= 0.5 - 0.2571(3.5) = \mathbf{-0.4}$.

(b) Predictions $\mathbf{-0.1429},\ 0.1143,\ 0.3714,\ 0.6286,\ 0.8857,\
\mathbf{1.1429}$ --- \textbf{two} lie outside $[0,1]$, one below and one
above. Neither is a probability.

(c) $x = (0.5 - b_0)/b_1 = 0.9/0.2571 = \mathbf{3.5}$ mm.

(d) Adding $(20,1)$: $S_{xx} = \mathbf{250.857}$, $S_{xy} = \mathbf{11.571}$,
slope $= \mathbf{0.0461}$, intercept $= \mathbf{0.3013}$, boundary
$= \mathbf{4.3086}$ mm --- it moved $\mathbf{+0.8086}$ mm. The point at
$x = 4$, a genuine malignant case, now predicts $0.4857 < 0.5$ and is
\textbf{misclassified}.

(e) The added point is already on the correct side and far from the boundary,
so it should be irrelevant --- but the old line predicts $-0.4 + 0.2571(20) =
4.74$ there, giving a residual of $(4.74-1)^2 = 14.0$, larger than the total
squared error of all six original points. Squared error keeps charging for
predictions that are \emph{too right}, so the fit rotates to reduce that one
residual and sacrifices a correct classification to do it.
\end{answerbox}

\A{U5.60} \textbf{Verified.}
\begin{answerbox}
\begin{center}\begin{tabular}{@{}rrr@{}}\toprule
$p$ & odds & log-odds\\
\midrule
0.05 & 0.0526 & $\mathbf{-2.9444}$\\
0.20 & 0.2500 & $-1.3863$\\
0.50 & 1.0000 & \textbf{0.0000}\\
0.80 & 4.0000 & $+1.3863$\\
0.95 & 19.0000 & $\mathbf{+2.9444}$\\
\bottomrule\end{tabular}\end{center}
\textbf{Antisymmetry:} $\operatorname{logit}(1-p) = -\operatorname{logit}(p)$
--- the column is symmetric about zero, and $p = 0.5$ sits exactly at the
origin.

\textbf{Implication:} relabelling the classes ($y \mapsto 1-y$) flips the sign
of every coefficient and of the intercept, and leaves the fitted
probabilities, the accuracy and the AUC completely unchanged. Which class you
call ``positive'' is a labelling convention, not a modelling decision.
\end{answerbox}
\emph{2.5 table, 1 antisymmetry, 1.5 implication.}

\A{U5.61} Mark in five parts:
\begin{answerbox}
(1) $P(y_i\mid x_i) = p_i^{y_i}(1-p_i)^{1-y_i}$ with $p_i = \sigma(w^\top x_i)$
--- the exponents are switches, not arithmetic.
(2) $L = \prod_i p_i^{y_i}(1-p_i)^{1-y_i}$.
(3) $\ell = \sum_i[y_i\ln p_i + (1-y_i)\ln(1-p_i)]$.
(4) Three chain-rule factors:
$\frac{\partial\ell}{\partial p_i} = \frac{y_i}{p_i} -
\frac{1-y_i}{1-p_i}$; $\frac{\partial p_i}{\partial z_i} = p_i(1-p_i)$;
$\frac{\partial z_i}{\partial w_j} = x_{ij}$.
(5) \textbf{The cancellation:}
\[\Big[\tfrac{y_i}{p_i} - \tfrac{1-y_i}{1-p_i}\Big]p_i(1-p_i)
= y_i(1-p_i) - (1-y_i)p_i = y_i - y_ip_i - p_i + y_ip_i = y_i - p_i.\]
Hence $\partial\ell/\partial w_j = \sum_i(y_i - p_i)x_{ij}$, i.e.
$\nabla_w\ell = X^\top(y-p)$.
\end{answerbox}

\A{U5.62} \textbf{Verified.}
\begin{answerbox}
(a) $z = Xw = (0,0,0,0)$, so $p = (0.5,0.5,0.5,0.5)$.
$\ell = 4\ln 0.5 = \mathbf{-2.772589}$; per-row NLL
$= 2.772589/4 = \mathbf{0.693147} = \ln 2$ --- the no-information baseline.

(b) $p - y = (\mathbf{+0.5},\ +0.5,\ -0.5,\ -0.5)$.
Intercept component: $0.5+0.5-0.5-0.5 = \mathbf{0}$.
Feature component: $(-1)(0.5) + 0(0.5) + 1(-0.5) + 2(-0.5) =
-0.5 - 0.5 - 1.0 = \mathbf{-2.0}$.
So $X^\top(p-y) = \mathbf{(0,\ -2.0)}$.

(c) $w \leftarrow (0,0) - 0.4(0,-2.0) = \mathbf{(0,\ 0.8)}$.

(d) $z = (-0.8,\ 0,\ 0.8,\ 1.6)$;
$p = (\mathbf{0.310026},\ 0.5,\ \mathbf{0.689974},\ \mathbf{0.832018})$;
$\ell = \mathbf{-1.619249}$, an improvement of $\mathbf{+1.153339}$.

(e) The intercept gradient is $\sum_i(p_i - y_i)$. At $w = 0$ every
$p_i = 0.5$, and the labels are \emph{balanced} --- two zeros and two ones ---
so the positive and negative discrepancies cancel exactly. The first step
therefore adjusts only the slope. With three ones and one zero it would not
have been zero.
\end{answerbox}

\A{U5.63} \textbf{Verified.} $\beta = \ln 3$.
\begin{answerbox}
(a) Odds ratio $= e^{\ln 3} = \mathbf{3}$ --- the odds triple.

(b)--(c)
\begin{center}\begin{tabular}{@{}rrrrr@{}}\toprule
start $p$ & odds & $\times 3$ & new $p$ & $\Delta p$\\
\midrule
0.1 & 0.1111 & 0.3333 & 0.2500 & $+0.1500$\\
0.2 & 0.2500 & 0.7500 & 0.4286 & $+0.2286$\\
\textbf{0.5} & 1.0000 & 3.0000 & 0.7500 & $\mathbf{+0.2500}$\\
0.8 & 4.0000 & 12.0000 & 0.9231 & $+0.1231$\\
\bottomrule\end{tabular}\end{center}
$\Delta p$ is largest starting from $p = 0.5$ and shrinks towards both ends ---
the same multiplicative change in the odds buys very different changes in
probability.

(d) $\beta/4 = 1.0986/4 = \mathbf{0.2747}$, and the largest observed
$\Delta p$ is $0.2500 < 0.2747$ \checkmark. The bound is not attained because
it is the \emph{instantaneous} slope at $p = 0.5$, while a one-unit step is a
finite move.

(e) The sentence is wrong twice. $e^\beta = 3$ is a change in the
\textbf{odds}, not the probability; and it \textbf{multiplies} rather than
adds, so ``by 200\%'' misdescribes even the odds. Correct: ``a one-unit
increase multiplies the odds by 3; the effect on the probability depends on
where you start, and is at most $0.2747$.''
\end{answerbox}

\A{U5.64} \textbf{Verified.}
\begin{answerbox}
\begin{center}\begin{tabular}{@{}rrrr@{}}\toprule
$z$ & $\sigma(z)$ & $1-\sigma(z)$ & $\sigma(1-\sigma)$\\
\midrule
$-3.0$ & 0.0474 & 0.9526 & 0.0452\\
$-1.5$ & 0.1824 & 0.8176 & 0.1491\\
$0.0$ & \textbf{0.5000} & 0.5000 & \textbf{0.2500}\\
$1.5$ & 0.8176 & 0.1824 & 0.1491\\
$3.0$ & 0.9526 & 0.0474 & 0.0452\\
\bottomrule\end{tabular}\end{center}
$\sigma(-3) = 0.0474 = 1 - 0.9526 = 1 - \sigma(3)$ \checkmark, and likewise at
$\pm1.5$.
The derivative is largest at $z = 0$, where it equals $\mathbf{0.25}$.

\textbf{Link to divide-by-four:} $\mathrm dp/\mathrm dx_j = \beta_j\sigma'(z)$,
and $\sigma' \le 0.25$ everywhere. So $|\beta_j|/4$ is the maximum possible
rate of change of the probability, attained only at the decision boundary
$p = 0.5$. The rule is nothing more than this table's middle row.
\end{answerbox}
\emph{2.5 table, 1 antisymmetry check, 1.5 the link.}

\A{U5.65} \textbf{Verified.}
\begin{answerbox}
(a)--(b)
\begin{center}\begin{tabular}{@{}rrrrrr@{}}\toprule
thr & accuracy & precision & recall & $F_1$ & $10\text{FN}+\text{FP}$\\
\midrule
0.10 & 0.8600 & 0.6610 & \textbf{0.9750} & 0.7879 & \textbf{30}\\
0.25 & 0.9000 & 0.7551 & 0.9250 & 0.8315 & 42\\
0.50 & \textbf{0.9200} & 0.8684 & 0.8250 & \textbf{0.8462} & 75\\
0.75 & 0.9133 & \textbf{0.9655} & 0.7000 & 0.8116 & 121\\
\bottomrule\end{tabular}\end{center}

(c) Accuracy chooses \textbf{0.50}; cost chooses \textbf{0.10}. Deploy
\textbf{0.10}. It has the \emph{worst} accuracy of the four and the lowest
cost by a factor of 2.5, because it misses one cancer instead of seven. The
$F_1$ also disagrees with the cost, since $F_1$ weights precision and recall
equally and the problem does not.

(d) ROC--AUC is computed over \emph{all} thresholds --- it measures only how
well the scores rank positives above negatives. Moving the threshold selects a
point on the curve; it does not change the curve. The same model with a
different threshold has the same AUC by construction.
\end{answerbox}

\A{U5.66} (a) \textbf{Perfect (or quasi-complete) separation}. (b) Once a
hyperplane classifies every training row correctly, scaling $w$ up by any
factor pushes each $p_i$ closer to its label, so $\ell$ increases
monotonically and its supremum (zero) is approached only as
$\|w\| \to \infty$. The supremum is never attained, so there is no maximising
$w$ --- the MLE does not exist. (c) Accuracy is a function of the \emph{sign}
of $w^\top x$, which is unchanged by scaling; it reads 1.0000 at
$\|w\| = 2.9$ and at $\|w\| = 51.8$ alike. (d) Add an L2 penalty:
$\min L(w) + \frac{1}{2C}\|w\|^2$, which makes the objective coercive and the
optimum finite. Cost: the estimates are shrunk and biased, and $C$ becomes a
hyperparameter needing cross-validation. Covered in \textbf{L12}.

\A{U5.67} Side by side:
$X^\top Xb = X^\top y$ (least squares) against
$X^\top(\sigma(Xw) - y) = 0$ (logistic). The first is \emph{linear} in $b$ ---
a $p\times p$ system with a closed-form solution. In the second the unknown
sits inside $\sigma$, a transcendental function; the equation is of the same
character as $xe^x = 1$, which has no solution in elementary functions.
Methods used instead: \textbf{gradient ascent / descent} (first order, cheap
per step, many steps) and \textbf{Newton's method / IRLS} (second order, uses
$H = -X^\top SX$, quadratic convergence --- six steps to machine precision on
a small problem). scikit-learn's default is L-BFGS, a quasi-Newton method.
\end{enumerate}

\subsection*{§C --- Ten-mark questions}

\begin{enumerate}[leftmargin=1.9cm]
\A{U5.68} Coverage: model and assumptions; squared \emph{vertical} deviations
with the differentiability/linear-system argument; the six-part derivation
including the Hessian; the matrix form and the rank condition; the four
properties; $\operatorname{Var}(b_1) = \sigma^2/S_{xx}$ with the design
implication; and the limits --- leverage, Anscombe, no causation. Worked fit:
U5.44's numbers.

\A{U5.69} Coverage: the probability model; $L$ and $\ell$; the
$\ell = A - cS$ argument for why $\sigma^2$ is irrelevant to the argmin;
$\hat\sigma^2_{\text{ML}} = \text{SSE}/n$ with the second-derivative check;
the bias $(n-2)/n$ with the two-constraints reason; $s^2 = \text{SSE}/(n-2)$;
and the Laplace substitution giving LAD. The recipe sentence --- state the
noise model, write the likelihood, take logs, maximise --- should appear.

\A{U5.70} Coverage: SST $=$ SSR $+$ SSE with the vanishing cross term; both
definitions of $R^2$ and the chain
$R^2 = \text{SS}_{\text{reg}}/\text{SS}_{\text{tot}} = b_1^2S_{xx}/S_{yy} =
r^2$; monotonicity via the set-the-coefficient-to-zero argument; adjusted
$R^2$ and its failure on 30 rows with 27 noise columns; negative held-out
$R^2$; Anscombe's four datasets with one $R^2$ and four diagnoses. The
sentence to land: \textbf{never report an $R^2$ without a residual plot beside
it}.

\A{U5.71} Coverage: definition; why training error cannot detect it (it is
monotone in capacity); the degree curve with the turn at degree 3 and the
$548\times$ blow-up at degree 20; the $n$ dependence (53{,}293 at $n=15$ against
0.2545 at $n=1000$); CV as the detector, with training picking degree 15 and
5-fold, LOOCV and the test set all picking 3; why one validation split is not
enough (ten splits chose degrees $[3,3,3,3,4,5,3,3,3,3]$ with the degree-3
estimate ranging 0.2243--0.4415); and the rule that everything learned is
refitted inside every fold.

\A{U5.72} Coverage: the straight-line failure with numbers (U5.59); odds,
log-odds, the four-line inversion to $\sigma$, and $\sigma' = \sigma(1-\sigma)$;
the Bernoulli likelihood and log-likelihood; \textbf{the gradient derivation
with the cancellation written out}; the transcendental score equations against
the linear normal equations; $H = -X^\top SX$ negative definite giving
uniqueness but not existence; and one hand-worked gradient step (U5.62).
\emph{The cancellation and the uniqueness-not-existence distinction are the two
places to be strict.}

\A{U5.73} Coverage: the one-line odds-ratio derivation
$e^{z+\beta_j}/e^z = e^{\beta_j}$; the five-starting-probabilities table
showing $\Delta p$ varying by a factor of several while the odds ratio is
constant; the divide-by-four rule as an upper bound; why raw coefficients are
incomparable across features with different units (and the cm-vs-mm example,
$\beta = 1.034$ per mm becoming 10.34 per cm with $e^\beta \approx 3\times10^4$);
the threshold as a decision made outside the model and on validation data; and
a worked cost-weighted threshold choice (U5.65).

\A{U5.74} A comparison table plus numbers from both. Expect: modelled quantity
--- $E[y]$ against $\ln\frac{p}{1-p}$; loss --- squared error against
cross-entropy; closed form --- yes against no; objective shape --- convex
quadratic against strictly concave $\ell$; solution --- solved against
optimised; coefficient --- additive change in $y$ against multiplicative change
in the odds; pathology --- leverage and outlier sensitivity against perfect
separation. The sentence worth crediting: \emph{least squares is solved,
logistic regression is optimised --- which is why Unit III came before
Unit V}.

\A{U5.75} Three derivations from one recipe.
\begin{answerbox}
\textbf{Gaussian:} $\ell = -\frac n2\log(2\pi\sigma^2) -
\frac{1}{2\sigma^2}\sum e_i^2$ $\Rightarrow$ maximise by minimising
$\sum e_i^2$ --- least squares.
\textbf{Laplace:} density $\frac{1}{2b}e^{-|e|/b}$ gives
$\ell = -n\log(2b) - \frac1b\sum|e_i|$ $\Rightarrow$ minimise $\sum|e_i|$ ---
least absolute deviations.
\textbf{Bernoulli:} $P = p^y(1-p)^{1-y}$ gives
$\ell = \sum[y\ln p + (1-y)\ln(1-p)]$ $\Rightarrow$ maximise, i.e.\ minimise
$-\ell/n$ --- binary cross-entropy.
\end{answerbox}
\end{enumerate}

\subsection*{§D --- Statements}

\begin{enumerate}[leftmargin=1.9cm]
\A{U5.76} \textbf{Incomplete.} Differentiability is
part of it, but the substantive reason is that the derivative of a square is
\emph{linear}, so the first-order conditions form a linear system with a
closed-form solution --- and squares are what Gaussian maximum likelihood
implies. Absolute values are also (almost everywhere) differentiable.
\A{U5.77} \textbf{False as stated.} Only when the model has an intercept ---
$\sum e_i = 0$ \emph{is} the first normal equation. The no-intercept fit of
U5.48 has $\sum e_i = +1.0$.
\A{U5.78} \textbf{Restates the formula; not an answer.} The reason is that the
residuals obey two linear constraints, so only $n-2$ are free and
$E[\text{SSE}] = (n-2)\sigma^2$.
\A{U5.79} \textbf{Must be qualified.} $R^2$ measures the proportion of variance
in $y$ explained by the model \emph{on the data it was computed from}. It says
nothing about generalisation, model correctness, or whether one point is
carrying the fit.
\A{U5.80} \textbf{False.} Training $R^2$ can never decrease when a predictor is
added, so it cannot arbitrate. Measured: 27 pure-noise columns took $R^2$ from
0.3738 to 0.9084 on 30 rows.
\A{U5.81} \textbf{False} --- and it is the error people most often make. $e^{\beta_j}$ is
a \textbf{multiplicative} change in the \textbf{odds}. The change in
probability depends on the starting probability and is bounded by
$|\beta_j|/4$.
\A{U5.82} \textbf{False.} It is linear in the \emph{log-odds}:
$\ln\frac{p}{1-p} = w^\top x + b$. The sigmoid is the inverse link that maps
that linear quantity back to a probability; the decision boundary
$w^\top x + b = 0$ is a hyperplane.
\A{U5.83} \textbf{False, and it is the classic separation tell.} Accuracy
depends only on the sign of $w^\top x$ and is invariant to scaling $w$. The
iris example scored 1.0000 with a coefficient of 51.83 and with 2.90 --- one
of which is a diverging MLE and one of which is a sensible fit.
\A{U5.84} \textbf{False.} \texttt{LogisticRegression()} applies L2
regularisation with $C = 1.0$ \emph{by default}. Pass \texttt{penalty=None} for
the actual MLE --- which is also why most users never encounter the separation
bug.
\end{enumerate}

\vspace{0.5em}\hrule

% =====================================================================
\section*{Mock paper B --- answers}

\subsection*{Section A (5 $\times$ 2 = 10)}
\begin{enumerate}[label=\textbf{\arabic*.}, leftmargin=*]
\item As U2.5. The compromise property: quadratic near zero (smooth, like MSE)
  and linear in the tail (bounded influence, like MAE), joined so that value and
  slope match.
\item As U3.4: $0 < \eta < 2/L$. At exactly $\eta = 2/L$ the multiplier is
  $|1-2| = 1$: the iterate neither converges nor diverges --- it bounces between
  two points forever. (Measured on $(w-3)^2$ at $\eta = 1$: $0 \to 6 \to 0$.)
\item $(1-p)^d$; $0.95^{20} = \mathbf{0.3585}$.
\item As U4.19 and U4.20: \texttt{MinMaxScaler}, because it is defined by
  $x_{\min}$ and $x_{\max}$, which one bad value sets.
\item As U5.11: $\sum e_i = 0$; the line passes through $(\bar x,\bar y)$;
  $\sum x_ie_i = 0$; SST $=$ SSR $+$ SSE. The \textbf{first} requires an
  intercept --- it is the first normal equation.
\end{enumerate}

\subsection*{Section B (any 3 of 4, 5 each)}
\begin{enumerate}[label=\textbf{\arabic*.}, leftmargin=*, start=6]
\item Identical to \textbf{U3.32}. Table, $\eta/\sqrt{1-\gamma} = 1.5811$,
  $s_t$ falls once $g_t^2 < s_{t-1}$; Adagrad's cannot because it is an
  undecayed sum of non-negative terms.
\item Identical to \textbf{U4.56}(a)--(c) plus the naming of \textbf{masking}.
\item Identical to \textbf{U5.44}. $S_{xx} = 10$, $S_{xy} = 17$, $b_1 = 1.7$,
  $b_0 = 0.6$; fitted $0.6, 2.3, 4.0, 5.7, 7.4$; residuals
  $0.4, -0.3, 0, -0.7, 0.6$; $\sum e_i = \sum x_ie_i = 0$;
  SSE $1.10$, SSR $28.90$, SST $30.00$, $R^2 = 0.9633$;
  SSR checked as $b_1^2S_{xx} = 28.90$.
\item Identical to \textbf{U5.62}. $p = (0.5,0.5,0.5,0.5)$,
  $\ell = -2.772589$, $X^\top(p-y) = (0, -2.0)$, $w = (0, 0.8)$,
  new $\ell = -1.619249$. The intercept gradient is zero because the labels
  are balanced and every $p_i = 0.5$, so the discrepancies cancel.
\end{enumerate}

\subsection*{Section C (10)}
\begin{enumerate}[label=\textbf{\arabic*.}, leftmargin=*, start=10]
\item (a) The chain of U3.46, with all seven update rules written correctly ---
  \textbf{4 marks}, roughly 0.6 per correct rule with its problem.
  (b) Curvature and frequency; momentum addresses curvature only ---
  \textbf{2 marks}, 1 each, and the second is the one students miss.
  (c) The derivation of U3.19 and U3.20 --- \textbf{2 marks}, 1 for
  $m_t = (1-\beta_1^t)g$, 1 for the $\pm\eta$ first step.
  (d) For: robustness to $\eta$ (2.24 decades against 1.31), works out of the
  box, handles sparsity. Against: measured \emph{last} of six on diabetes at 30
  epochs (0.4929 against momentum's 0.4846), three extra hyperparameters, and
  the loss got worse without decay. What it buys: a wider window of working
  learning rates, not a better optimum --- \textbf{2 marks}, and the last
  sentence is required for the second of them.
\end{enumerate}

\vspace{0.4em}\hrule

\section*{Mock paper C --- answers}

\subsection*{Section A (5 $\times$ 2 = 10)}
\begin{enumerate}[label=\textbf{\arabic*.}, leftmargin=*]
\item As U1.1. Credit scoring: $T$ = classify an applicant good/bad;
  $E$ = past applicants with repayment outcomes; $P$ = ROC--AUC, or expected
  cost under the bank's error costs.
\item As U2.8. At $p = 0.5$ the loss is $-\log 0.5 = \log 2 = \mathbf{0.6931}$.
\item As U3.8 and U3.9: $\eta_{\text{eff}} \approx \eta/(1-\beta)$.
\item $\sqrt{1-p} = \sqrt{0.75} = \mathbf{0.8660}$ --- the sd falls to 86.6\%
  of its true value.
\item As U5.31: odds $= p/(1-p) \in (0,\infty)$; log-odds
  $\in (-\infty,\infty)$; only the \textbf{log-odds} matches the range of
  $w^\top x + b$, which is why the model is built there.
\end{enumerate}

\subsection*{Section B (any 3 of 4, 5 each)}
\begin{enumerate}[label=\textbf{\arabic*.}, leftmargin=*, start=6]
\item \textbf{U2.18} and \textbf{U2.19} combined: MSE 1.7, MAE 1.0,
  RMSE 1.3038, Huber at $\delta=1.5$ is 0.75; shares 0.7353 and 0.5000.
\item Identical to \textbf{U4.65}. $S_{xy} = -15$, $S_{xx} = 2$, slope $-7.5$,
  intercept 45.8333, $R^2 = 0.0696$ against one-hot's 1.0.
\item \textbf{U5.54}(a)--(b) and \textbf{U5.52} combined. $R^2 = 1 - 1.10/30 = 0.9633$;
  $\hat\sigma^2_{\text{ML}} = 0.22$; $s^2 = 0.3667$; $s = 0.6055$;
  bias factor $0.6$, understatement $40\%$. The reason is the two residual
  constraints, so only $n-2$ are free --- \emph{not} ``because we divide by
  $n$''.
\item Identical to \textbf{U5.63}. Odds ratio 3; new probabilities
  $0.25, 0.4286, 0.75, 0.9231$; $\Delta p$ largest at $p = 0.5$ ($+0.25$);
  divide-by-four bound $0.2747$; the sentence must be corrected to say the
  \textbf{odds} are \textbf{multiplied} by 3.
\end{enumerate}

\subsection*{Section C (10)}
\begin{enumerate}[label=\textbf{\arabic*.}, leftmargin=*, start=10]
\item (a) The rule of U4.37, and five of: a mean or median, a standard
  deviation, a category vocabulary, a feature subset, a hyperparameter, a
  resampled row, a PCA basis --- \textbf{2 marks}.
  (b) The four leaks with their measured sizes, ranked: resampling $+0.78$,
  selection $+0.31$, scaling $+0.0034$, imputation $+0.0003$. The ranking
  principle: a step that sees $y$ leaks far more than one that does not, and a
  step that changes the \emph{rows} leaks most of all --- \textbf{4 marks},
  1 per leak, with the ranking principle required for the fourth.
  (c) U1.18: selection on all the data planted the correlation it then
  ``found''; honest answer 0.50 --- \textbf{2 marks}.
  (d) \texttt{Pipeline} (or \texttt{ColumnTransformer} inside one) passed to
  \texttt{cross\_val\_score}/\texttt{GridSearchCV}, which refits the whole chain
  per fold; \texttt{imblearn.pipeline.Pipeline} where resampling is involved.
  The operation it cannot protect you from is \textbf{row-dropping} --- a
  transformer may change columns but not rows, so outlier removal and
  deduplication live outside the pipeline and must be handled by hand
  --- \textbf{2 marks}, and the last point is the discriminating one.
\end{enumerate}

\vspace{0.6em}\hrule\vspace{0.4em}
\noindent\small\textbf{End of key.} 314 questions, 63 of them numerical. Regenerate every numeric

\end{document}
